 %%%%%%%%%%%%%%%%%%%%%%% file template.tex %%%%%%%%%%%%%%%%%%%%%%%%%
%
% This is a general template file for the LaTeX package SVJour3
% for Springer journals.          Springer Heidelberg 2010/09/16
%
% Copy it to a new file with a new name and use it as the basis
% for your article. Delete % signs as needed.
%
% This template includes a few options for different layouts and
% content for various journals. Please consult a previous issue of
% your journal as needed.
%
%%%%%%%%%%%%%%%%%%%%%%%%%%%%%%%%%%%%%%%%%%%%%%%%%%%%%%%%%%%%%%%%%%%
%
% First comes an example EPS file -- just ignore it and
% proceed on the \documentclass line
% your LaTeX will extract the file if required
\begin{filecontents*}{example.eps}
%!PS-Adobe-3.0 EPSF-3.0
%%BoundingBox: 19 19 221 221
%%CreationDate: Mon Sep 29 1997
%%Creator: programmed by hand (JK)
%%EndComments
gsave
newpath
  20 20 moveto
  20 220 lineto
  220 220 lineto
  220 20 lineto
closepath
2 setlinewidth
gsave
  .4 setgray fill
grestore
stroke
grestore
\end{filecontents*}

\RequirePackage{fix-cm}
%\documentclass{svjour3}                     % onecolumn (standard format)
%\documentclass[smallcondensed]{svjour3}     % onecolumn (ditto)
\documentclass[smallextended]{svjour3}       % onecolumn (second format)
%\documentclass[twocolumn]{svjour3}          % twocolumn
%
\smartqed  % flush right qed marks, e.g. at end of proof

\usepackage{graphicx}
%
% \usepackage{mathptmx}      % use Times fonts if available on your TeX system
%
% insert here the call for the packages your document requires
%\usepackage{latexsym}
% etc.
%
% please place your own definitions here and don't use \def but
% \newcommand{}{}
%
% Insert the name of "your journal" with
% \journalname{myjournal}
\usepackage[utf8]{inputenc}
\usepackage[usenames,dvipsnames]{xcolor}

\usepackage{amsmath}
\usepackage{amsfonts,mathabx}
\usepackage[all]{xy}
\usepackage[english]{babel}
\usepackage{graphicx}
\usepackage{enumerate}
\usepackage{wrapfig}
\usepackage[colorlinks]{hyperref}
\newtheorem{thm}{Theorem}

\newtheorem{lem}{Lemma}[section]      % Lemma environment with numbering 
\renewcommand{\thelem}{\arabic{section}.\arabic{lem}}
\newtheorem{prop}[lem]{Proposition}
\newtheorem{cor}[lem]{Corollary}

\newtheorem{defi}[lem]{Definition}    % Definition environment with 
%                                    % numbering consecutive to theorems
\newtheorem{rem}[lem]{Remark}             % Unnumbered environment for remarks.
\newtheorem{ex}[lem]{Example}     
\newtheorem{cex}[lem]{Counter-Example}  
%

  \def\ZZ{\mathbb{Z}}
  \def\r{\mathfrak{r}}
  \def\D{\mathcal{D}}
  \def\Z{\mathbb{Z}}
  \def\H{\mathbb{H}}
  \def\M{\mathcal{M}}
  \def\s{\mathfrak S}
  \def\v{\mathfrak v}


  \def\stab{\mathrm{Stab}}
\def\Kbar{\overline{K}}
\def\au{\overrightarrow{u}}
  \def\U{\mathcal U}
  \def\ab{ \mathrm{ab}}
  \def\Hom{\mathrm{Hom}}
  \def\isom{Isom}
  \def\lsem{[\![}
  \def\rsem{]\!]}
\def\V{\mathcal V}
\def\F{\mathcal F}
\def\C{\mathcal C}
\def\A{\mathcal A}
\def\TT{\mathcal T}
\def\R{\mathbb{R}}
\def\E{\mathbf E}
\def\G{\mathbf G}
\def\S{\mathcal S}
\def\N{\mathbb N}
\def\T{\mathcal{T}}
\def\rot{\mathbf{rot}}
\def\div{\mathbf{div}}
\def\grad{\overrightarrow{\mathbf{grad}}}
\def\d{\mathrm{d}}
\def\S{\mathcal{S}}
\def\K{\mathbf K}
\def\D{\mathcal{D}}
\def\SO{\mathrm{SO}}
\def\E{\mathbb{E}}
\def\BTZ{{\E_{0}^{1,2}}}
\def\sing{\mathrm{Sing}}
\def\RR{\mathbb{R}}
\newcommand{\mass}[1]{{\mathbb E}^{1,2}_{#1}}
\def\isom{\mathrm{Isom}}
\def\L{\mathrm{L}}
\def\graph{\mathrm{Graph}}
\def\Teich{\mathrm{Teich}}
\def\Hol{\mathrm{Hol}}
\def\Reg{\mathrm{Reg}}
\def\fix{\mathrm{Fix}}
\def\teich{\mathrm{Teich}}
\def\susp{\mathrm{susp}}
\def\pgcd{M_1\wedge M_2}
\def\ppcm{M_1\vee M_2}
\def\reg{\Reg}

\def\DD{\mathbb{D}}
\def\isomH{\isom^+(\H)}
\def\V{\mathcal V}
\def\F{\mathcal F}
\def\T{\mathcal{T}}
\def\BTZtilde{{\widetilde \E}^{1,2}_0}
\def\G{\mathbf G}
\def\S{\mathcal S}
\def\W{\mathcal W} 
\def\MM{\overline M}
\def\stab{\mathrm{Stab}}
\def\RR{\mathbb R}
\def\R{\mathcal R}
\def\CC{\mathbb{C}}
\def\J{\mathcal{J}}
\def\SS{\mathbb{S}}

\def\H{\mathbb{H}}
\def\C{\mathcal C}
\def\E{\mathbb{E}}
\def\sing{\mathrm{Sing}}
\def\isom{\mathrm{Isom}}
\def\d{\mathrm{d}}
\def\Z{\mathbb Z}
\def\N{\mathbb{N}}
\def\reg{\mathrm{Reg}}
\def\D{\mathcal{D}}
\def\BTZ{\mathbb{E}^{1,2}_0}
\def\U{\mathcal{U}}
\def\fix{\mathrm{Fix}}
\def\BTZext{\mathrm{Ext}_{\mathrm{BTZ}}}
\def\Reg{\mathrm{Reg}}
\def\mink{\mathbb{E}^{1,2}}
\def\Dbar{\overline{\D}}
\def\O{\mathcal{O}}
\def\adm{\mathcal{P}_\Sigma}
\def\EH{\mathcal H}
\def\grad{\mathrm{grad}}
\def\Hol{\mathrm{Hol}}
\def\Cat{\mathbf{C}}
\def\Top{\mathbf{Top}}
\def\M{\mathcal M}
\def\cotan{\mathrm{ct}}
\def\catGX{\mathbf{Var_{G,X}}}
\newcommand{\catE}[1]{\mathbf{Var_{\mass{#1}}}}
\newcommand{\catX}[1]{\mathbf{Var_{X_{#1}}}}
\def\Im{\mathrm{Im}}
\def\et{~\text{et}~}
\def\DD{\mathbb{D}}
\def\r{\mathfrak{r}}
\def\Tame{\mathrm{Tame}}
\def\WTame{\mathrm{WTame}}
\def\massA{\mass{\geq 0}}
\def\pgcd{M_1 \wedge M_2}
\def\ppcm{M_1 \vee M_2}
\def\teich{\mathrm{Teich}}
\def\fix{\mathrm{Fix}}
\def\dec{\mathrm{dec}}
\def\sing{\mathrm{Sing}}
\def\Hol{\mathrm{Hol}}
\def\so{\mathfrak{so}}
\def\pen{\mathcal{P}}
\def\susp{\mathrm{susp}}
\def\UU{\mathbb{U}}
\def\PP{\mathbb{P}}
\def\P{\mathcal{P}}
\def\Hom{\mathrm{Hom}}
\def\a{\lambda}
\def\b{\mu}
\def\cc{\beta}
\def\dd{\alpha}
\def\A{\mathcal A}
\def\B{\mathcal B}
\def\GG{\mathrm{O}_{0}(1,2 ; \RR)}
\def\Hbar{\overline{\H^2}}
\def\RP{\mathbb{RP}^1}
\def\PSL{\mathrm{PSL}}
\def\HS{\mathbb{HS}}
\def\struc{\mathbf{Struct}}
\def\GL{\mathrm{GL}}
\def\SL{\mathrm{SL}}
\def\G{\mathcal G}
\def\Aut{\mathrm{Aut}}
\def\supp{\mathrm{supp}}
\def\Homeo{\mathrm{Homeo}}
\def\cool{{cool}}
\def\ord{\mathrm{Ord}}

\def\vecM{\overrightarrow{\E}^{1,2}}
\def\X{\overrightarrow X}
\def\Aut{\mathrm{Aut}}
\def\gradH{{\overrightarrow{\grad} \mathcal H_{\bar\kappa}}}
\def\Wild{\mathrm{Wild}}
\newcommand{\fonction}[5]{\displaystyle#1:\begin{array}{l|rcl}
& \displaystyle #2 & \longrightarrow & \displaystyle #3 \\
    & \displaystyle #4 & \longmapsto & \displaystyle #5 \end{array}}
\newcommand{\fonctiontrois}[9]{\displaystyle#1:\begin{array}{l|rcl}
& \displaystyle #2 & \longrightarrow & \displaystyle #3 \\
    & \displaystyle #4 & \longmapsto & \displaystyle #5\\
    & \displaystyle #6 & \longmapsto & \displaystyle #7\\
    & \displaystyle #8 & \longmapsto & \displaystyle #9\end{array}}    
\newcommand{\fonctionn}[4]{\displaystyle\begin{array}{l|rcl}
& \displaystyle #1 & \longrightarrow & \displaystyle #2 \\
    & \displaystyle #3 & \longmapsto & \displaystyle #4 \end{array}}

\newcommand{\fonctiondeux}[7]{\displaystyle#1:\begin{array}{l|rcl}
& \displaystyle #2 & \longrightarrow & \displaystyle #3 \\
	& \displaystyle #4 & \longmapsto & \displaystyle #5 \\ 
	& \displaystyle #6 & \longmapsto & \displaystyle #7\end{array}}	

\newcommand{\cyl}[1]{\CC_\infty^{(#1)}}

\definecolor{MyGreen}{rgb}{0.0,.5,0.0}
\newcommand{\edl}[1]{{\color{MyGreen}{#1}}}
\newcommand{\leo}[1]{{\color{MyGreen} [\textbf{Leo: } #1]}}

\definecolor{MyDarkRed}{rgb}{0.7,0,0}
\newcommand{\edt}[1]{{\color{MyDarkRed}{#1}}}
%\usepackage{float}
%%% Start of metadata

%

\title{On branched coverings of singular (G,X)-manifolds}
\author{L\'eo Brunswic}
\institute{L\'eo Brunswic \at
              Univ Lyon, Ens de Lyon, Univ Lyon1, CNRS, Centre de Recherche Astrophysique de Lyon UMR5574, F--69007, Lyon, France \\
              \email{leo@brunswic.fr}
}
%  
\def\CatSpread{\mathbf{Spr}}
\def\CatSingGX{\mathbf{SingGX}}
\smartqed
%%% End of user-defined macros %%%
% \renewcommand{qedsymbol}{QED}
\begin{document}

\maketitle
\begin{abstract}    % type your abstract below
Branched coverings have a long history from the ramification of Riemann surfaces to the realization of 3-manifolds as coverings branched over knots; from geometrical topology to algebraic geometry. 
The present work investigates branched covering "à la Fox" of $(G,X)$-manifolds and is  three fold. First, we introduce a general class of singular $(G,X)$-manifolds and present elementary theory with examples to convince the reader of its efficiency and generality. Second, we recall and enrich the current state of the art (based upon Montesinos) and develop a Galois theory for such branched coverings together with a description  of the fiber above branching points. 
Finally, find local properties ensuring the existence of developing maps for singular $(G,X)$-manifolds based upong the branched coverings framework; in particular, we localize the existence of developing maps for such singular manifolds.  The present work  is particularly relevant for non-metric singular $(G,X)$-manifolds such as Lorentzian or projective ones considered by Barbot, Bonsante, Suhyoung Choi, Danciger, Seppi, Schlenker or the author among others. Examples are provides along the way but the main application we provide is a uniformization theorem ``à la Mess'' for singular localLy Minkowski manifolds with BTZ-like singularities.
\end{abstract}

%%%%%%%%%%%%%%%%%%%%   Start of main body of article
\tableofcontents

        \section{Introduction}
\subsection{Some motivational examples}
Consider $\H$ the hyperbolic plane and $
\isomH\simeq\GG\simeq \mathrm{PSL}(2;\RR)$\footnote{The index $0$ indicates the connected component of identity} its group of direct isometries . The quotient of $\H$ by a geometrically finite discrete torsion-free subgroup of $\isomH$, say $\Gamma$, gives a locally hyperbolic manifold, ie  a $\H$-manifold, say $\Sigma:=\Gamma\backslash\H$. Such  a $\Sigma$ admits a developing map $\D:\widetilde \Sigma\rightarrow \H$ and an holonomy homomorphism
$\rho:\pi_1(\Sigma)\rightarrow \isomH$ where $\widetilde \Sigma$ is the universal covering of $\Sigma$ and $\pi_1(\Sigma)$ its fundamental group. If furthermore the holonomies of the ends of $\Sigma$ are parabolic then one can associate to each end an ideal point ``at infinity", effectively compactifying $\Sigma$ to some $\overline \Sigma$ (which is the Freudenthal compactification \cite{MR931673} of $ \Sigma$). The geometric model $(\isomH,\H)$ is an instance of a geometrical structure in the sense of Ehresmann \cite{MR794193_loc_hom}, popularised by Thurston \cite{thurston-3mani} and others  \cite{MR957518,mess,MR2170277,barbot_globally_2004,MR2522946,MR1299730,MR755162} ie a couple $(G,X)$ 
where $X$ is a locally path-connected Hausdorff topological space and $G$ a group acting analytically on $X$ by homeomorphisms. Here ``analytically'' means that for all $g,h\in G$, if the restrictions to any non trivial open subset $\U$  of $g$ and $h$ agree then $g=h$. A $(G,X)$-manifold is then a second countable Hausdorff topological space $M$ endowed with an atlas whose change of charts lies in $G$. The existence of a developing map from the universal covering of such a manifold $M$ to the model space $X$ and of an holonomy from its fundamental group to the model group $G$ is a general property as long as the universal covering exists. \footnote{In \cite{MR794193_loc_hom} elements of proof are presented for $G$ a Lie group  acting on a manifold $X$, a careful proof for the group  of conformal transformations or isometries of a Riemannian  manifold see \cite{MR0031310} or \cite{MR755162}. The generality considered here is contained in neither references but the arguments still work \cite{thesis}.
}


The manipulation of such points ``at infinity" (or even irrational conical singularities), though legal and quite simple in this specific context, requires some additional care in the general setting of $(G,X)$-manifolds especially if one considers non-metric analytical structures such as conformal structures, Lorentzian structures or affine structures ie $(\mathrm{Aff}(\RR^n),\RR^n)$-structures. Compactifications and completions are not $(G,X)$-manifolds anymore but rather ''singular" $(G,X)$-manifolds. The notion of singularity is not a unified concept: singularities are described in specific context.
Most ``singularities" arise either via 
\begin{itemize}
 \item completion or compactification procedures as presented above;
 \item quotienting a regular $(G,X)$-manifold by a discrete group acting properly discontinuously by automorphisms;
 \item gluings of simplices of the model space $X$ (when it makes sense);
 \item suspensions such as presented in \cite{MR1668340,Particules_1}.
\end{itemize}


In all instances, the regular locus is open and dense. Additional properties are however required for any weak form of analyticity to hold. Indeed consider $X=\RR^2$ and $G$ either $
\mathrm{Isom}(\mathbb E^2) \simeq\mathrm{O}_2(\RR) \ltimes \RR^2$ or $\mathrm{Aff}(\RR^2)\simeq \mathrm{GL}_2(\RR)\ltimes \RR^2$ acting as usual on  $\RR^2$ ; then consider the manifold $M = \RR^2$ together with the $(G,X)$-structure induced by the atlas composed of a unique chart $\mathrm{Id} : \RR^2\setminus S \rightarrow \RR^2\setminus S$ for some subset $S$.  If $S$ is a singleton, using the metric it is easy to check one can extend uniquely the flat metric structure to the whole $\RR^2$. The same statement is true in the affine case.
However, if $S$ is a  line segment, say  $S=\{0\}\times [-1,1]$, uniqueness of atlas extension can be proved only in the metric case while in the affine case, uniqueness fails. Indeed, say in the affine case, one can add any one  of  the charts
$$\fonctiondeux{\varphi_{\lambda,\mu}}{\RR\times ]-1,1[}{\RR\times ]-1,1[}{(x,y),x\leq 0}{(\lambda x,y)}{(x,y),x\geq 0}{(\mu x,y)},\quad \quad \lambda,\mu>0.$$
The new regular locus is then the complement of two points and its holonomy is trivial if and only if $\lambda=\mu$. In any case, different extension of the initial atlas give rise to inequivalent affine structures.\footnote{If however we only consider extensions of the $(G,X)$-atlas which are defined on the whole $\RR^2$ then in both cases the extension is unique.} 
\subsection{Singular $(G,X)$-manifolds and branched coverings}

We show that an efficient definition for singular $(G,X)$-manifold is that they have a $(G,X)$-atlas defined almost everywhere in the sense that its support is {\it portly}: open, dense and locally connected in the whole manifold \footnote{A subset $Y\subset X$ is locally connected in $X$ if for all open connected $\W\subset X$, the intersection $\W\cap Y$ is connected. }. The complement of a portly subset is called  {\it skeletal}.  In section \ref{sec:GXsing} we write a primer on such singular $(G,X)$-manifolds generalizing elementary results to this context,  showing this definition is satisfactory. 

Beyond elementary analyticity statements we would like a good notion of developing map, holonomy and universal covering for singular $(G,X)$-manifolds.  
A natural first step in this direction is to consider, for some singular $(G,X)$-manifold $M$, the universal covering $\widetilde\reg(M)$ of its regular locus $\reg(M)$, which gives rise to a map $\widetilde\reg(M)\rightarrow M$, then tries to complete $\widetilde\reg(M)$ in a natural way.  

We argue that this natural way has been provided by Fox in the 1950's \cite{MR0123298,MR1152367} with the notions of {\it spread}, complete spread and completion of spread and then extended to include branched covering; we refer the interested reader to Montesinos \cite{MR3060537} for an up to date review. We shall recall fundamental properties of coverings {\it à la} Fox as well as extend it in two directions. First, we define a Galois Theory of branched coverings  with  the important addition of a characterisation of Galoisian universal branched coverings above given locii. In particular, we prove in Proposition \ref{prop:Galoisian_semilocal} that for a Hausdorff connected locally path-connected topological space $X$, the universal covering branched  above a skeletal
locus $S$ is Galoisian if and only if $X\setminus S$ is {\it semi-locally simply connected in $X$} in a sense we will define. 
Second, we give a path space description of the universal covering branched above a given locus akin to that of the universal unbranched covering, see Theorem \ref{theo:ramified_covering_path}. To this end, we introduce a notion of {\it almost trivial loop}. 
Most intermediary results and definitions presented are {\it folkloric}; however, to the author's knowledge, most are not presented in the generality considered in the present work and the notion of ''semi-locally simply connected in`` is new as well as the path description of the maximal branched covering.

Coming back to singular $(G,X)$-manifolds we introduce the classes of {\it tame singular $(G,X)$-manifolds} and {\it virtually tame singular $(G,X)$-manifolds} that admit natural developing maps.  We prove localization Theorems to ensure those global properties can be efficiently derived from local ones. See Theorems \ref{theo:local_tame} and \ref{theo:local_virt_tame}.

\subsection{General references}
We use freely standard notions from general topology (separation and countable basis hypothesese, compact-open topology,...), algebraic topology (first fundamental group, semi-local simple connectedness, Galois correspondance,...), group actions and category theory (categories, commutative diagrams, projective limits,...). We refer the reader uncomfortable with some of these notions to reference texts \cite{MR2048350,bourbaki2013general,steen1978counterexamples,MR1700700,hatcher2005algebraic,MR1712872}. To the author knowledge, there is no reference textbook on the fundamentals of $(G,X)$-manifolds we thus refer the reader unfamiliar with the theory to what seems to be the closest equivalent \cite{MR957518} as well as 
Benzecri thesis \cite{Benzecri1960VaritsLA} and \cite{MR1299730,MR2522946}  which include some metric background in the $(G,X)$-manifold context and some elementary account of conical singularities. Having a particular interest for Lorentzian manifolds, the author kindly asks the reader to pardon his tendency to draw non-trivial examples from those. The reader unfamiliar with the geometry of Lorentzian manifolds may find useful \cite{oneil,barbot2010algebre}.


\begin{acknowledgements}
This work is part of a project that has received funding from the European Research Council (ERC) under the European Union's Horizon 2020 research and innovation programme (grant agreement ERC advanced grant 740021--ARTHUS, PI: Thomas Buchert). 
The author has no relevant financial or non-financial interests to disclose otherwise.
Data sharing not applicable to this article as no datasets were generated or analysed during the current study.The author thanks Thierry Barbot and the reviwer for their careful reading, corrections and sugestions, José María Montesinos Amilibia for his interest and advices,  Thomas Buchert for his continuous support,  Masoud Hasani, Andrea Seppi, Uira Noberto Matos de Almeida,  Ivan Izmestiev and Roman Prosanov for valuable discussions.
\end{acknowledgements}



\section{Singular (G,X)-manifolds (I): bases and examples}\label{sec:GXsing}

We give ourselves some analytical structure $(G,X)$, ie a connected, locally connected, Hausdorff topological space $X$ together with a faithful action of a group $G$ by homeomorphisms such that for all $\phi_1,\phi_2\in G$ and all non trivial open subset $\U\subset X$, if the action of $\phi_1$ and $\phi_1$ agree on $\U$ they agree on $X$.
A $(G,X)$-atlas on a topological space $M$ is a family $(\U_i,\V_i,\varphi_i)_{i\in I}$ where for all  $i\in I$, $\U_i$ is an open subset of $M$, $\V_i$ is an open subset of $X$ and $\varphi_i:\U_i\rightarrow \V_i$ an homeomorphism, with the property that for all $i,j\in I$ such that $\U_i\cap \U_j\neq \emptyset$ and for all connected component $\W$ of $\U_i\cap \U_j$, there exists a $\phi \in G$ such that $\phi\circ \varphi_{i|\W} = \varphi_{j|\W}$. By analyticity of $(G,X)$, such a $\phi$ is unique. The support $\supp(\A)$ of such an atlas $\A$ is defined as the union of the $\U_i$. 
A $(G,X)$-manifold is a second countable Hausdorff topological space $M$ endowed with an $(G,X)$-atlas whose support is $M$.
%The fundamental theorem of such manifolds, coming from the analyticity of the action of $G$, is that the exists a map $\D : \widetilde M \rightarrow X$ which is $\rho$-equivariant for some representation $\rho:\pi_1(M)\rightarrow G$ of the fundamental group of $M$ into $G$. Furthermore, such couples $(\d,\rho)$ only differ by conjugation by some element of  $G$. 

The simplest example of singularities occuring in the literature of $(G,X)$-manifolds are conical singularities in homogeneous Riemannian manifolds. 
\begin{ex} \label{ex:conical}Let $\alpha>0$ and $n\geq 2$, define $\mathbb E_\alpha^n = \RR^{n-2}\times \RR^2$ endowed with cylindrical coordinates $(z,r,\theta) \in \RR^{n-2}\times \RR_+\times \RR/\ZZ$ and the Riemannian metric $\mathbf{g} = \d z^2+\d r^2 + \alpha r^2\d \theta^2$. 
This Riemannian metric is well defined and locally Euclidean on the complement of the domain $\{r=0\}$. 
One can thus endow the domain $\{r>0\}$ with a $(\mathbb E^n,\RR^n \rtimes \SO(n))$-structure. However, this geometrical structure cannot be extended to the whole $\RR^n$ except for $\alpha=2\pi$. An obstruction is given by the holonomy $\rho:\pi_1(\{r>0\})=\ZZ \rightarrow \RR^n\rtimes \SO(n)$, indeed $\rho(1)$ is the rotation of angle $\alpha$ which is non trivial  when $\alpha\notin  2\pi\N^*$.
\end{ex}
In the previous example, manifolds with conical singularities have an underlying metric space structure that gives a natural notion of isomorphism class. In a way, the author feels that such an underlying structure distracts from the fact that the notion of isomorphism between singular $(G,X)$-manifolds does not depend on such metrics and is actually more general. Gluing projective polyhedra leads to ''singular" projective manifolds and it would feel "unatural" to define isomorphism classes for such manifolds based upon some metric choice. As illustration of the richness of non metric singularities and their usefulness we wish to refer for instance to \cite{Particules_1} for a fundamental study of conical 2+1 Lorentzian singularities or \cite{MR3190306,MR3814342} for their role in geometric transitions. 



The present section is devoted to the most basic definitions and properties of the notion of singular $(G,X)$-manifolds we introduce.  It encompass most singularities of the literature and we show some elementary properties ensuring the theory of such manifold works well. 




\subsection{Topological preliminaries}

Recall that for $X$ a topological space, a subset $\U\subset X$ is locally connected in $X$ if for all open connected $\W\subset X$, the intersection $\W\cap \U$ is connected. Note in particular that if $X$ is connected then $\U$ is connected.

\begin{defi}[Portly/skeletal subset]\label{def:portly_skeletal_subset} Let $X$ be a Hausdorff locally connected topological space, a subset $\U$ of $X$ is portly (in $X$) if $\U$ is open dense and locally connected in $X$. A subset $S\subset X$ is skeletal if its complement in $X$ is portly.
\end{defi}

\begin{ex}\label{ex:following_example_skeletal_itemize} The following example of $S$ are skeletal in $X$.
\begin{itemize}
 \item $S$ finite  and $X$ a surface;
 \item $S$ a tame knot  in a $3$-manifold $X$;
 \item $S$ the $k$-skeleton of a pure $n$-dimensional simplicial complex $X$ with $k\leq n-2$;
 \item $S$ a Cantor wildly embedded into a $n$-manifold $X$.
\end{itemize} 
 \end{ex}





Portly subsets have nice properties, the proof of which rely on usual connectedness and/or density arguments and are thus left to the reader,
most can be found in \cite{MR0123298,MR1152367,MR2599123,thesis}. Consider a locally connected Hausdorff topological space $X$.
\begin{itemize}
 \item finite intersections of portly subsets are portly;
 \item if $\U\subset \V$, $\U$ is portly in $X$ and $\V$ is open then $\V$ is portly (hence unions of portly subsets are portly);
 \item if $\U\subset \V$ and $\U$ is portly in $X$ then $\U$ is portly in $\V$;
 \item if $\U$ is portly in $\V$ and $\V$ is portly in $X$ then $\U$ is portly in $X$;

 \item locally portly subsets are portly (i.e. if $(\U_i)_{i\in I}$ is a family of open subsets and for each $i\in I$, $\V_i\subset \U_i$ is a portly subset of $\U_i$, then $\bigcup_{i\in I}\V_i$ is portly in $\bigcup_{i\in I}\U_i$);
 \item if $X\xrightarrow{f} Y$ is a local homeomorphism and $\U\subset Y$ is portly in $Y$ then $f^{-1}(\U)$ is portly in $X$.
 \item if $X\xrightarrow{f} Y$ is a continuous open surjective and $\U\subset X$ is portly in $X$ then $f(\U)$ is portly in $Y
 $.
\end{itemize}




\subsection{Almost everywhere $(G,X)$-atlases}



The objects of this section are a.e $(G,X)$-atlases which are the first step toward $(G,X)$-structures defined almost everywhere. 
\begin{defi}\label{def:almost_everywhere_atlas_topological} An almost everywhere $(G,X)$-atlas on a topological space $M$ is a $(G,X)$-atlas supported by some portly subset of $M$. 
\end{defi}
In order to avoid caveats, we need to build suitable maximality properties for such $(G,X)$-atlases.  We can first extend usual properties to this context.
A  $(G,X)$-atlas is complete if for all $(\U,\V,\phi)\in\A$ and all $\U'\subset \U$ we have $ (\U',\phi(\U'),\phi_{|\U'}) \in \A$. We can complete any $(G,X)$-atlas by adding missing sub-charts, the completed atlas is still a $(G,X)$-atlas.
In our context, a $(G,X)$-atlas $\A = (\U_i,\V_i,\phi_i)_{i\in I}$ is {\it thinner} than a $(G,X)$-atlas $\B = (\U_j',\V_j',\phi_j')_{j\in J}$ if the support of $\A$ is a subset of the support of $\B$ and  for all $p\in \mathrm{supp} \A$ and all  $j\in J$ such that $p\in \U_j'$, there exists $g\in G$ and $i\in I$ such that $p\in \U_i\subset \U_j'$ and $\phi_{i} = g\circ \phi_{j|\U_i}'$. Notice that if such an $\A$ is thinner than such a $\B$ then the completion of $\A$ is still thinner than $\B$.

What we need and now provide is the flexibility to manipulate atlases that are defined up to some "negligible set" akin to the similar notion in measure theory. 



    \begin{lem} \label{lem:GX_atlas_union}
        Let $N$ be a locally connected topological space.
        Let $(\B_k)_{k\in K}$ be a family of a.e.$(G,X)$-atlases on $N$. Assume for  all $(k,k')\in K^2$, there exists an a.e $(G,X)$-atlas $\mathcal A$ such that $\A$ is thinner than both $\B_{k}$ and $\B_{k'}$.


            Then the union $\bigcup_{k\in K} \B_k$ is an a.e.$(G,X)$-atlas.
    \end{lem}
					\begin{proof}
					Let $k,k'\in K$ and write $\B_k=(\U_i,\V_i,\phi_i)_{i\in I(k)}$ and
					$\B_{k'}=(\U_i,\V_i,\phi_i)_{i\in I(k')}$.
					Let $i\in I(k), j\in I(k')$ such that $\U_i\cap \U_j \neq \emptyset$ and let $\widehat\W \subset\U_i\cap\U_j$ be an open connected non empty subset.
					Let $\A$ be an a.e.$(G,X)$-atlas supported by some portly subset $M\subset N$ which is thinner that both $\B_{k}$ and $\B_{k'}$. Without loss of generality, we assume that $\A$ is complete. 
					
					\begin{itemize}
					 \item Consider $p\in \widehat\W\cap M$ and a chart $(\U,\V,\phi)$ of $\A$ around $p$ such that $\U\subset \widehat\W$. The existence of such a chart is guaranteed since $\A$ is  complete and thinner than $\B_k$, $\B_{k'}$.
					Let $\W \subset \U$ be an open connected neighborhood of $p$.
					There exists unique $g,h\in G$ such that
					$$\phi_{|\W} =g\circ \phi_{i|\W} \quad \quad  \phi_{|\W} =h\circ \phi_{j|\W},$$
					hence
					$$\phi_{j|\W}=(h^{-1} g)\circ \phi_{i|\W}.$$
					 Define $g_{p,\U,\W}:=h^{-1}g$ so that $\phi_{j|\W}= g_{p,\U,\W}\circ \phi_{i|\W}$.
				
					\item  We show that $g_{p,\U,\W}$ depends neither on $\U$ nor on $\W$ nor on $p$.
					
					Fixing $p\in \widehat \W\cap M$, consider
					 $(\U,\V,\phi)$ and $(\U',\V',\phi')$ two charts of $\A$ on a neighborhood of $p$ 
					as well as $\W$ and $\W'$ two open connected neighborhoods of  $p$ in $\U$ and $\U'$ respectively.
					Since $$\forall x\in \phi_i(\W\cap \W'),\quad g_{p,\U,\W}(x)= \phi_j(\phi_i^{-1}(x))=g_{p,\U',\W'}(x)$$  we thus have $g_{p,\U,\W}=g_{p,\U',\W'}$ so
					$g_p:=g_{p,\U,\W}$ only depends on $p$. 
					Furthermore, if $q\in \W$, we prove the same way that
					$g_{q,\U,\W}=g_{p,\U,\W}$.
					Finally, the map $\W\cap M\rightarrow G, p\mapsto g_p$ is locally constant.
					Since $M$ is portly and since $\widehat \W$ is connected, the intersection $\widehat \W \cap M$ is connected and $p\mapsto g_p$ is constant on $\widehat \W\cap M$.
					
					\item We proved there exists $g\in G$  such that
					$$\forall x\in M\cap \widehat \W, \quad  \phi_j(x)=g\circ \phi_i(x)$$
					Since $\phi_i,\phi_j$ and $g$ are continuous, $X$ is Hausdorff and $M$ is dense, 
					the intersection $\widehat \W\cap M$ is then dense in 
					$\widehat \W$ and
					$$\phi_{j|\widehat\W}=g\circ \phi_{i|\widehat\W}.$$
					\end{itemize}
					Finally, $\B_k\cup \B_{k'}$ is an a.e. $(G,X)$-atlas, moreover  $k,k'\in K$ are arbitrary thus $\bigcup_{k\in K} \B_k$ is an a.e. $(G,X)$-atlas.
					
					\qed\end{proof}
					
					
					\begin{cor} \label{cor:maximal_atlas}
					Let $M$ be a locally connected Hausdorff topological space. 
					Consider the set $\mathcal E$ of  a.e $(G,X)$-atlas $\mathcal A$ on $M$ ordered by $\A\leq \B$ if the completion of $\A$ is thinner than $\B$.
					
					Then, for every $\mathcal A \in \mathcal E$ the maximum $$ \max\,(\B ~:~ \A\leq \B \in \mathcal E\})$$
					is well defined.
					\end{cor}
                    \begin{proof}
                      By Lemma \ref{lem:GX_atlas_union}, this maximum is  $\bigcup_{\B\geq \A}\B$.
                    \qed\end{proof}					
                    
                    \subsection{Category of Singular $(G,X)$-manifolds}
                    
                    \begin{defi}\label{def:structure_topological_space_atlas} An a.e. $(G,X)$-structure on a topological space $M$ is an a.e. $(G,X)$-atlas on $M$.                     
                    \end{defi}
                    
					\begin{defi}\label{def:singular_manifold_locally_connected} A singular $(G,X)$-manifold is a locally connected Hausdorff second countable topological space $M$ endowed with an a.e. $(G,X)$-structure maximal in the sense of Corollary \ref{cor:maximal_atlas}.
					
					The support of its a.e. $(G,X)$-atlas is called its regular locus denoted $\reg(M)$ and the complement of $\reg(M)$ is called the singular locus denoted $\sing(M)$.
					\end{defi}
                    
                    
					\begin{defi}[a.e.$(G,X)$-morphism]\label{def:morphism} Let
					$M$ and $N$ be two singular $(G,X)$-manifolds.

					An a.e.morphism $\phi:M\rightarrow N$ is a continuous map such that there exists portly subsets $\U\subset \reg(M)$ and $\V\subset \reg(N)$ such that
					$\phi_{|\U}^{|\V}:\U\rightarrow \V$ is a $(G,X)$-morphism.
					\end{defi}

					\begin{rem}\label{rem:composition_morphisms_morphisms_category} Composition of a.e. $(G,X)$-morphisms is an a.e. $(G,X)$-morphisms. The category $\CatSingGX$ of singular $(G,X)$-manifold which morphisms are a.e $(G,X)$-morphisms is then well defined. 
					Isomorphisms of singular $(G,X)$-manifolds are then maps that are both  homeomorphism and a.e. $(G,X)$-morphism.
					\end{rem}

                    \begin{rem}\label{rem:morphisms_analytical_two_equal} A.e. $(G,X)$-morphisms are analytical: if two are equal on an open set, they are equal everywhere.
                     
                    \end{rem}

					
					
					\begin{rem}\label{prop:GX_sing_morph_local}
					Let $M\rightarrow N$ be a continuous map between singular $(G,X)$-manifolds.
					The following properties are equivalent.
					\begin{enumerate}[(i)]
					 \item $f$ is an a.e. $(G,X)$-morphism;
					 \item $f$ is locally an a.e. $(G,X)$-morphism ;
					 \item $f_{|\U}$ is an a.e. $(G,X)$-morphisms for some portly set $\U\subset M$.
					\end{enumerate}
					\end{rem}

                \begin{lem}\label{lem:singular_manifolds_morphism_set} 
                   Let $M$ and $N$ be singular $(G,X)$-manifolds and let $f:M\rightarrow N$ an be a.e $(G,X)$-morphism. 
                   The set $\mathcal E$ of open subset $\U$ such that $\U\subset \reg(M)$ and $f$ induces a $(G,X)$-morphism $\U\rightarrow \reg(N)$.   
                   
                   Ordered by the inclusion, $\mathcal E$ has a maximum.
                    
                \end{lem}
                \begin{proof}
                  Let $\U_0:=\bigcup_{\U\in \mathcal E}\U$, the map $f_{|\U_0}^{|\reg(N)}$ is locally a $(G,X)$-morphism hence a $(G,X)$-morphism. Finally, $\U_0\in \mathcal E$ and is certainly the maximum.
                \qed\end{proof}

                \begin{defi}\label{def:singular_manifolds_morphism_maximal} Let $M$ and $N$ be singular $(G,X)$-manifolds and let $f:M\rightarrow N$ an a.e $(G,X)$-morphism.
                The maximal open subset of $M$ on which $f$ induces a $(G,X)$-morphism is the regular locus of $f$ which we denote by $\reg(f)$.
                \end{defi}

                The following properties give characterisations of $\reg(f)$ in common situations, limiting what can go wrong.
                    \begin{lem}\label{lem:regf1}
                     Let $M$ and $N$ be  singular $(G,X)$-manifolds and let $f:M\rightarrow N$ an a.e $(G,X)$-morphism. Then, 
                     $$\reg(f)=\reg(M)\cap f^{-1}(\reg(N)).$$
                    \end{lem}
                    \begin{proof}
                     The left hand side is clearly included in the right hand side.                   
                     Via restrictions to charts, it suffices to prove the special case of $N=X$ hence $M=f^{-1}(\reg(N))$.
                    
                     Let $x\in \reg(M)$ and let $(\U,\V,\varphi)$ be a connected chart around $x$. The map $f\circ\varphi^{-1}$ induces a $(G,X)$-morphisms from $\varphi(\U\cap\reg(f))$ to $X$. Since $\reg(f)$ is locally connected in $M$ and $\U$ is connected, $\varphi(\U\cap\reg(f))$ is also connected and there exists a unique $g\in G$ such that $$\forall y\in \varphi(\U\cap\reg(f)),\quad  f\circ \varphi^{-1}(y)=g(y).$$
                     Since $\varphi:\U\rightarrow \V$ is a homeomorphism and $\U\cap\reg(f)$ is dense in $\U$, then  $\varphi(\Reg(f)\cap \U)$ is dense in $\V$. Since $f\circ \varphi^{-1}$ is continuous, we have
                     $$\forall y\in \U,\quad  f(y)=g\circ \varphi(y).$$
                     Finally, $f$ is a regular on $\U$. The result follows.
                    \qed\end{proof}
       

                \begin{rem}\label{rem:though_may_tempted_include}
                 Though we may be tempted to include Riemann surfaces in the a.e. $(G,X)$-manifold framework, the previous Lemma shows there is no simple definition of Riemann surfaces as $(G,X)$-manifolds that allows branched map to be a.e. $(G,X)$-morphisms. Indeed, those maps are locally modeled on $z\mapsto z^p$ and $0\notin \reg(z\mapsto z^p)$ except if $p=1$. The above Lemma would thus force $p=1$ hence does not allow ramification.  
                \end{rem}
                \begin{lem}\label{lem:aepullback}
                    Let $M\xrightarrow{f} N$ be a continuous map with $M$ locally connected Hausdorff topological space and $N$ a singular $(G,X)$-manifold. 
                    
                    If $f$ is a.e. a local homeomorphism there exists a unique a.e. $(G,X)$-structure on $M$ such that $f$ is an a.e. $(G,X)$-morphism.
                \end{lem}
                \begin{proof} 
                 Let $\mathcal O\subset M$ and $\mathcal P\subset \reg(N)$  both portly such that $f_{|\mathcal O}^{|\P}$ is a local homeomorphism.
                 
                 By usual results, there exists a unique a.e. $(G,X)$-structure $\A$ supported on $\mathcal O$ such that $f_{|\mathcal O}^{|\P}$ is a $(G,X)$-morphism. 
                 Since $\mathcal O, \P$ are portly, we thus defined an a.e $(G,X)$-structure on $M$ such that $f$ in an a.e. $(G,X)$-morphism. 
                 
                 Assume that $M$ is endowed with an a.e $(G,X)$-structure $\B$ supported on some portly subset $\mathcal O'\subset M$ such that $f_{|\mathcal O'}^{|\P}$ is a $(G,X)$-morphism.  In particular $f_{|\mathcal O'\cap \mathcal O}^{|f(\mathcal O'\cap \mathcal O)}$ is a $(G,X)$-morphism, therefore the $\A_{|\mathcal O'\cap \mathcal O}=\B_{|\mathcal O'\cap \mathcal O}$. By Corollary \ref{cor:maximal_atlas} and maximality $\A=\B$ and $\O=\O'$.
                \qed\end{proof}
                \begin{defi}\label{def:notations_lemma_lem_aepullback}
                 With the notations of Lemma \ref{lem:aepullback}, the unique  a.e. $(G,X)$-structure $f^*\A$ on $M$ is called the pullback by $f$ of the a.e. $(G,X)$-structure $\A$ of $N$.
                 
                \end{defi}
                \begin{cor}\label{cor:locally_connected_hausdorff_topological}
                 Let $L\xrightarrow{f}M\xrightarrow{g} N$ be locally connected Hausdorff topological space. And let $\A$ be an a.e $(G,X)$-structure on $N$. 
                 Then $f^*g^*\A = (g\circ f)^*\A$.
                \end{cor}


% 
%                 \begin{lem}\label{lem:aepushforward} Let $M$ be a singular $(G,X)$-manifold, let $N$ be a locally connected Hausdorff topological space and let $f: M \rightarrow N$ be a continuous map. 
% 
%                 If $f$ is surjective, open and almost everywhere a local homeomorphism, then there exists at most one a.e. $(G,X)$-structure on $N$ such that $f$ is an a.e. $(G,X)$-morphism.
%                 Furthermore, one must have
%                 $$\reg(f)=\reg(M)\cap \{x\in M  ~|~ f \text{ is a local homeomorphism around }x\}.$$
%                 \end{lem}
%                 \begin{proof}
%                 Let $$\U = \reg(M)\cap \{x\in M  ~|~ f \text{ is a local homeomorphism around }x\}$$
%                 by assumption $\U$ is an intersection of portly subsets thus portly.
%                 
%                 
%                 There exists a unique $(G,X)$-structure on $f(\U)$ such that $f_{|\U}^{|f(\U)}$ is a $(G,X)$-morphism. Since $f$ is open and surjective, $f(\U)$ is portly; the pushforward by $f$ of the $(G,X)$-structure on $\U$ defines an a.e $(G,X)$-atlas $\A$ inducing an a.e. $(G,X)$-structure  on $N$. We endow $N$ with this a.e. $(G,X)$-structure so that $f$ is an a.e. $(G,X)$-morphism.
%                 Furthermore, since $(G,X)$-morphisms are local homeomorphism we have $\reg(f)\subset \U$; by construction, $\reg(f) \supset \U$ so we have $\reg(f)=\U$. 
%                 
%                 If $\B$ is the maximal a.e. $(G,X)$-atlas of an a.e. $(G,X)$-structure on $N$ such that $f$ is an a.e. $(G,X)$-morphism, then there exists some portly subsets $\V \subset \reg(M)$ and $\W\subset \reg(N)$ such that $f_{|\V}^{|\W}$ is a $(G,X)$-morphism. In particular, $f_{|\V}^{|\W}$ is a local homeomorphism so $\V\subset \U$ and we may choose $\W\subset f(\U)$. 
%                 By uniqueness, $\B_{|\W} = \A_{|\W}$.
%                 In other word, $\A_{|\W}$ is thinner than both $\A$ and $\B$; therefore, $\B=\A$ by maximality of $\A$ and $\B$. 
%                 
%                 
%                 \qed\end{proof}
%                 \begin{defi}
%                  With the notations of Lemma \ref{lem:aepushforward}, the unique a.e. $(G,X)$-structure on $N$ is called the pushforward by $f$ of the a.e. $(G,X)$-structure of $N$.
%                 \end{defi}



					
                \begin{lem}{}\label{lem:GX_sing_homeo_reg}
						 Let $M$ and $N$ be singular $(G,X)$-manifolds and let $f:M\rightarrow N$ an a.e $(G,X)$-morphism. Then,
					$$\reg(f)=\reg(M)\cap \{x\in M  ~|~ f \text{ is a local homeomorphism around }x\}$$
					\end{lem}
					\begin{proof}
					By definition, $\reg(f)\subset \reg(M)$, furthermore a $(G,X)$-morphism is always a local homeomorphism; therefore the left hand side is included in the right hand side. Furthermore, let $\A,\B$ be the a.e. $(G,X)$-atlases of $M,N$ respectively. Since $f$ is a.e. a $(G,X)$-morphism, then it is a.e. a local homemorphism; we may thus apply Lemma \ref{lem:aepullback} to show that $f^*\B = \A$. In particular, if $f_{|\U}^{|\V}$ is an homemorphism for some $\U,\V$ then 
					$\reg(\U) = f^{-1}(\reg(\V))\cap \U$ and $f:\reg(\U)\rightarrow \reg(\V)$ is a $(G,X)$-morphism. 
					Therefore, if $f$ is a local homemorphism around $x$ and in $\reg(M)$ then $x\in \reg(\U)$ for an $\U$ such as above thus $f$ is $(G,X)$-morphism around $x$ and $x\in \reg(f)$.
					\qed\end{proof}


					\begin{cor}\label{cor:morphisms_also_local_homemorphism} If an a.e. $(G,X)$-morphisms  is also a local homemorphism then it preserves both regular and singular locii.
					\end{cor}
					
					\begin{rem}\label{rem:consequence_isomorphism_singular_manifolds} As a consequence, 
					 an isomorphism $M\xrightarrow{f} N$ of singular $(G,X)$-manifolds is an isomorphism of almost everywhere $(G,X)$-atlas in the sense that given the maximal almost everywhere $(G,X)$-atlases $\mathcal A$ and $\mathcal B$ of $M$ and $N$ respectively, the pullback $f^*\mathcal B$ and $\mathcal A$ are equal. It is then an isomorphism in every natural ways relative to singular $(G,X)$-manifolds.
					 
					 
                    \end{rem}

\subsection{Common construction of  singular $(G,X)$-manifolds}
\label{sec:construction_example}

From the literature we can list common construction leading to "informal" singular $(G,X)$-manifolds:
\begin{itemize}

 \item branched coverings;
 \item gluings of projective polyhedron; \cite{MR2127379,MR2170277,Particules_2,brunswic2021cauchy,spinexp}
  \item a quotient of some $(G,X)$-manifold $M$ by some discrete group $\Gamma$ acting via $(G,X)$-morphisms \cite{MR1177168};
 \item suspensions \cite{MR1668340,Particules_1};
 \item completions \cite{brunswic2021cauchy} 
\end{itemize}

Branched covering will be studied more thoroughly in the following sections. When $X\subset \mathbb RP^n$ and $G\subset \mathrm{PGL}(n+1; \RR)$, gluings of projective polyhedra via elements of $G$ have a natural $(G,X)$-structure on the complement of the $(n-2)$-skeleton. The latter being skeletal, such gluings are singular $(G,X)$-manifolds with our definitions. This includes in particular most usual Lie group actions such as semi-Riemannian isometry groups as well as conformal and simplectic actions. We will not cover completions.

\subsubsection{Quotients}
The first example that comes to mind is the following.
\begin{ex}\label{ex:consider_identifying_point_wise_stabilizer} 
  Consider $X=M = \mathbb E^n$ and $G=\isom^+(\mathbb E^n)$ identifying $\mathbb E^n$ with $\RR^n$ and $n\in \N_{\geq 2}$. The point-wise stabilizer of $S = \RR^{n-2}\times \{(0,0)\}$ is isomorphic to $\mathrm {SO}(2;\RR)$,  we choose $\Gamma \subset \stab(S)$ a finite group of rotations fr order $p$.
 The group $\Gamma$ acts freely on $M\setminus S$ which is portly, therefore the quotient $\Gamma \backslash M$ has a natural singular $(G,X)$-structure. 
 The metric on $\reg(\Gamma\backslash M)$ extends continuously to $\Gamma\backslash M$. As a \textbf{metric space} $\Gamma \backslash M$ is isomorphic to $\mathbb E^n_ {2\pi/p}$ described in Example \ref{ex:conical}.
 
  Variations on this example can be done either by replacing $X$ by another Riemannian spaceform $\mathbb H^n$ or $\mathbb S^n$. More generaly one can take submodel of the projective model space :
 $X \subset \mathbb K P^n$ and $G \subset \mathrm{PGL}(n;\mathbb K)$ for $\mathbb K \in \{\mathbb R,\mathbb C\}$.
\end{ex}

We can relate this example to a fokloric result for regular $(G,X)$-manifolds.
\begin{prop}\label{prop:GXquotient} Let $M$ be a $(G,X)$-manifold and let $\Gamma$ be a discrete group acting freely properly discontinuously on $M$. Then there exists a unique $(G,X)$-structure on $\Gamma\backslash M$  such that the natural projection $M\rightarrow \Gamma\backslash M$ is a $(G,X)$-morphism.
\end{prop}
In the example above, the action is not free and the singular locus of the quotient corresponds to the "non-free locus" the projection of the set of points of non trivial stabilizer. Nonentheless, we may extend the last Proposition to a.e. $(G,X)$-manifolds.
\begin{prop}\label{prop:aeGXquotient}
 Let $M$ be a singular $(G,X)$-manifold, let $\Gamma$ be a discrete group acting properly discontinuously on $M$ by a.e. $(G,X)$-morphisms. 
 
 If there exists a portly $\Gamma$-invariant subset $\U$ on which $\Gamma$ acts freely then $\Gamma\backslash M$ has a unique a.e. $(G,X)$-structure such that the natural projection $M \rightarrow \Gamma\backslash M$ is an a.e. $(G,X)$-morphism.
\end{prop}
\begin{proof}
 By assumption, the quotient $\Gamma \backslash M$ is Hausdorff. Futhermore,
 the projection $M \xrightarrow{\pi} \Gamma \backslash M$  is open, hence as an open quotient of a locally connected and second countable space, $\Gamma\backslash M$ is locally connected and second countable. 
 Let $\U \subset M$ be a portly $\Gamma$-invariant subset on which $\Gamma$ acts freely, since $\Gamma$ acts by a.e. $(G,X)$-morphisms, in particular $\Gamma$ preserves $\reg(\U)$. By Proposition \ref{prop:GXquotient}, the quotient $\Gamma \backslash \reg(\U)$ admits a unique $(G,X)$-structure such that  
 $\pi_{|\reg(\U)}^{|\Gamma \backslash \reg(\U)}$ is a $(G,X)$-morphism.
 Since $\pi$ is open and surjective, the image of a portly subset of $M$ is a portly subset of $\Gamma\backslash M$; in particular, $\pi(\reg(\U))$ is portly. Therefore, the $(G,X)$-structure  on $\Gamma\backslash \reg(\U)$ induces an a.e. $(G,X)$-structure $\A$ on $\Gamma\backslash M$ for which $\pi$ is an a.e. $(G,X)$-morphism.
 
 If $\Gamma\backslash M$ is endowed with some other a.e. $(G,X)$-structure $\B$ for which $\pi$ is an a.e. $(G,X)$-morphism then $\pi_{|\V}^{|\reg(\B)}$ is a $(G,X)$-morphism for some portly subset $\V\subset \reg(M)$ and thus $\pi_{|\U\cap\V}^{|\reg(\B)\cap \reg(\A)}$ is a $(G,X)$-morphism both for $\A$ and $\B$. In particular, on the portly subset $\pi(\U\cap \V)\subset \supp(\A)\cap \supp(\B)$, the two $(G,X)$-structures agree. By Lemma \ref{lem:GX_atlas_union} and maximality of $\A$ and $\B$, we conclude that $\A=\B$.
 
\qed\end{proof}



\begin{ex}\label{ex:quotient_base}
 Consider $\mathbb D_{\infty}$ the topological space $(\RR_{+}^*\times \RR \cup \{(0,0)\},\T)$ where $\T$ is the topology generated by the one of $\RR_+^*\times \RR$  and open subsets $\{(r,\theta) \in \mathbb D_{\infty} ~|~ r<r_0\}$ for $r_0>0$. We may refer to the space $\DD_\infty$ as  the "infinite angle disc".
 
 The infinite angle cylinder eg  the space $\cyl{n} :=\DD_\infty \times \RR^n$  can be endowed with an a.e. affine structure for some choice $G\subset \mathrm{Aff}(\RR^n)$ and $X\subset \RR^n$   via pullback by a local homeomorphism $\D : \{r>0\}\rightarrow X$. We call this singular $(G,X)$-manifold $\widetilde M$ and  naturally, $\sing(\widetilde M)\subset \{r=0\}$.
 
 Take $\Gamma_\alpha =\Z$ acting on $\widetilde M$ via $k\cdot (r,\theta,z)= (r,\theta+k\alpha,z)$ for some $\alpha\in \RR_+^*$, assume that $\Gamma_\alpha$ acts via a.e. $(G,X)$-morphisms, then $M:=\Gamma_\alpha\backslash \widetilde M$ is naturally a singular $(G,X)$-manifold homeomorphic to $\RR^n$ endowed with cylindrical coordinates.  Its singular locus is $\sing(M) = \{r=0\}$.
 
 Consider the special case $X=\mass{}$ eg $\RR^3$ endowed  with the bilinear form $Q(t,x,y)=x^2+y^2-t^2$ of signature $(2,1)$ ; with $G = \isom_0(X)$ the identity component of the affine isometry group.  $\mass{}$ is the 3-dimensional Minkowski space and $G$ the group of isometries preserving both orientation and time orientation (see Sections \ref{sec:Minkowski} and \ref{sec:mink_caus} for more details). We can recover several instances of Lorentzian singularities \cite{Particules_1,barbot2012particles} for different choices of $\D$. 
 \begin{itemize}
  \item   $\D(r,\theta,z) = (z+\sigma\theta/\alpha, r\cos(\theta/\alpha),r\sin(\theta/\alpha))$ then $M$ is the so-called massive particle of mass $m=1-\alpha/2\pi$ and spin $\sigma$.
  \item  Consider $$\D\left(r,\theta,t\right)= 
P_\theta\cdot (x_0+rv + t u)
$$
with $x_0\in X$ arbitrary,  $(P_\theta)_{\theta\in \RR}$ the 1-parameter family of parabolic isometries fixing point wise $x_0+\RR u$ some future lightlike vector and $\langle v|u\rangle<0$.  
Then $M$ is a so-called extreme BTZ whitehole that we shall focus on in Section \ref{sec:BTZ_uniform}. Choosing $\D'(r,\theta,t) = \D(r,\theta,-t)$ gives an extreme BTZ blackhole instead.
\item Variations allow to construct so-called Misner particles, BTZ white/blackholes. It suffices to choose  $u$ spacelike instead and $v$ spacelike or future/past timelike  respectively.

 \end{itemize}

So-called tachyons (resp. photons) can be obtained as variations of the above example with a hyperbolic (resp. parabolic)  action of $\Gamma=\ZZ$ on the coordinate $\theta$ with the identication $\DD_\infty = \RR_+^*\times \widetilde {\RR\mathbb P^1} \cup \{(0,0)\}$.

All these examples can also be obtained via gluings of polyedra.
\end{ex}

\begin{rem}\label{rem:previous_examples_instances_developing}
 The previous examples are instances of developing maps defined on the universal covering of some singular space branched over its singular locus. This is the archetype we will describe with more precisions in section~\ref{sec:GXsing2}.
\end{rem}


\subsubsection{Suspensions} \label{sec:susp}

One can construct conical singularities (without spin) iteratively following a method of Thurston \cite{MR1668340}. 
\begin{enumerate}
 \item Let say that any compact $(\SS^1,\SO(2))$-manifold is regular so there is no singular model to construct. 
 \item Assume you have constructed local models of singularities in $(\SS^n,\SO(n+1))$-manifolds. A $(\SS^n,\SO(n+1))$-manifolds with conical singularities is then a metric space locally isometric to either of local models.
 \item You can then construct local model of singularity for $(\SS^{n+1},\SO(n+2))$.
 \begin{enumerate}
 \item Take a $(\SS^n,\SO(n+1))$-manifold $M$ with conical singularities and with $M$ homeomorphic to $\SS^n$,  denote its Riemannian metric by $\mathbf{g}$ which is defined on the complement of the singularities;  
 \item Define the {\it suspension} of $M$ by 
 $\widehat M = (M\times \RR_+)/\sim$ the set of $(\theta,r)$ where $(\theta_1,0)\sim (\theta_2,0)$ for all $\theta_1,\theta_2\in M$; 
 \item Endow $\widehat M$ with a Riemannian metric $$\widehat{\mathbf{g}}(\theta,r) = \mathbf{d r}\otimes \mathbf{d r}+  \sin(r)^2 \mathbf{g}(\theta)$$ which is "regular", eg locally isometric to $\SS^{n+1}$, on rays through "regular" points of $M$.  
 \end{enumerate}
\end{enumerate}
This way, we obtain usual conical singularities in dimension 2, then so-called "collisions" of conical lines in dimension 3, etc.

From $(\SS^{n-1},\SO(n))$-manifolds with conical singularities one can construct conical local models for $(\mathbb H^n,SO(n,1))$ and $(\mathbb E^n, \RR^n\times \SO(n))$ replacing $\sin(r)^2$ by $\sinh^2(r)$ and $r^2$ respectively in the 
formula for Riemannian metric of the suspension above. 
One can show that gluings of finitely many Spherical (or Euclidean or hyperbolic) polyhedra give rise to such singularities on the codimension $k\geq 2$ skeletons of the underlying simplicial complex. 

We are led to the following definition.
\begin{defi}[Topological suspension]\label{def:topological_suspension}
 Let $\Sigma$ be a  topological space, the topological suspension of $\Sigma$ is
 $$ \susp(\Sigma):= (\Sigma \times \RR_+)/{\sim}$$
where $(\sigma,r) \sim (\sigma',r')$ if $r=r'=0$; endowed with the topology generated by the  quotient topology on $\{r>0\}$ and the image of $\Sigma\times [0,\varepsilon[$ for $\varepsilon>0$.
We also define $O:= \{r=0\}$.
 \end{defi}
 \begin{rem}\label{rem:one_may_check_suspension} One may check the suspension preserves Hausdorff as well as first and second countable properties. 
 \end{rem}
 \begin{rem}\label{rem:infinite_angle_disc_suspension} The infinite angle disc $\DD_\infty$ is the suspension of $\RR$.
  
 \end{rem}

\begin{cex}\label{cex:topology_chose_same_quotient} The topology we chose is the same as the quotient topology if $\Sigma$ is compact, however the suspension
 $\susp(\RR)$ endowed with the quotient topology  is not first countable. Indeed, a basis of open neighborhoods of $O$ is given by $\U_f := \{(x,r)~|~ r<f(r)\}$ for $f:\RR\rightarrow \RR_+^*$ continuous. 
 For any sequence $(f_n)_{n\in \N}$ we may construct $g:\RR\rightarrow \RR$ continuous and such that
 $\forall n\in \N, g(n):= \frac{1}{2}\min(f_k : k\leq n )$. 
 This way,  $\forall n\in \N, \U_{f_n} \nsubset \U_g$ which proves $(\U_{f_n})_{n\in \N}$ is not a basis of neighborhoods of $O$.
\end{cex}

\begin{rem}\label{rem:topological_suspension_fact_functor}
 The topological suspension is in fact a functor from the category of topological space to itself as continuous maps can be extended radially in a natural way. Suspension preserves injectivity, surjectivity and embeddings.
\end{rem}

 
Geometrical suspension as above can be described abstractly as a lift of the topological suspension functor $\susp$
 $$\xymatrix{ \mathbf{Top} \ar[rr]^\susp&& \mathbf{Top} \\
 \mathbf{SingGX} \ar[u]\ar@{-->}[rr]&& \mathbf{SingHY}\ar[u] 
 }$$ 
 where $(G,X)$ and $(H,Y)$ are analytical structures and the vertical arrows are the forgetful functor. Thurston method provides a functors $\mathbf{Sing}\boldsymbol{\mathbb S}^n \rightarrow \mathbf{Sing}\boldsymbol{\mathbb S}^{n+1}$ and Bonsante, Barbot, Schlenker \cite{Particules_1} generalized to a functor
  $\mathbf{Sing}\boldsymbol{\mathbb {HS}}^2 \rightarrow \mathbf{Sing}\boldsymbol{\E}^{1,2}$ where $\mathbb{HS}^2:=(\mass{}\setminus\{O\})/\RR_+^*$ is the space of rays from the origin in $\mass{}$ with automorphism group $G:=\SO_0(1,2)$.
For these suspensions, the image of regular $(G,X)$-manifolds are regular except possibly at $O$ and are regular at $O$ if and only if $\Sigma$ is homeomorphic to a sphere.  

To ensure that iterated suspensions described above are of this type we need to prove the following.
 
 
\begin{prop} \label{prop:susp_portly} Let $\Sigma$ be a Hausdorff, connected, locally connected topological space and let $\U\subset \Sigma$ be a portly subset. 
Then $\susp(\U)\setminus \{O\}$ is a portly subset of $\susp(\Sigma)$.
\end{prop}
 Allow us to recall a classical connectivity Lemma the proof of which is omitted.
\begin{lem}\label{lem:connected_union}
 Let $M$ be a topological space and let $(\U_i)_{i\in I}$ be a non-empty family of non-empty connected open subset. Define on $I$ the equivalence relation $\sim$ generated by $i\sim j$ if $\U_i\cap \U_j \neq \emptyset$.
 
 The $\bigcup_{i\in I} \U_i$ is connected if and only if $I/\sim$ is a singleton.
 
\end{lem}

\begin{proof}[of Proposition \ref{prop:susp_portly}]
 Let $\pi : \Sigma \times \RR_+ \rightarrow \susp(\Sigma)$ be the natural projection and define $\V := \susp(\U)\setminus\{O\}$.  We also introduce $\psi : \susp(\Sigma)\setminus \{O\} \rightarrow \Sigma$ induced by the projection $\Sigma\times \RR_+^* \rightarrow \Sigma$.
 Notice that $\pi_{|\Sigma\times \RR_+^*}$ is an embedding, therefore $\V = \pi(\U\times \RR_+^*)$ is open.  Since $\U$ is dense in $\Sigma$ and since $\psi$ is open, then $\V = \psi^{-1}(\U)$ is dense in $\susp(\Sigma)\setminus \{O\}$. For any neighborhood $\W$ of $\{O\}$ and for all $\sigma \in \Sigma$, for $r$ small enough $\pi(\sigma,r)\in \W$; then $O$ is also in the closure of $\V$ and $\V$ is dense in $\susp(\Sigma)$. 
 
 Let $\W$ be a connected open subset of $\susp(\Sigma)$.
 \begin{itemize}
  \item Assume that $O\notin \W$. We may write $\W = \bigcup_{i\in I} \W_i$ with $\W_i = \pi(\widehat \W_i\times ]r_i,r_i'[)$, $0<r_i<r_i'$ and $\widehat \W_i$ open subset of $\Sigma$. Since $\Sigma$ is locally connected, we assume without loss of generality that each $\widehat \W_i$ is connected. Denote by $\sim_1$ the  equivalence relation defined in Lemma \ref{lem:connected_union} for the family $(\W_i)_{i\in I}$ so that $I/\sim_1$ is singleton. 
  
  Then consider the family $(\W_i\cap \V)_{i\in I}$ and define $\sim_2$ the associated equivalence relation on $I$. 
 On the one hand, if $\W_i \cap \W_j\neq \emptyset$ then by density of $\V$ 
 we have $(\W_i\cap \V) \cap (\W_j\cap\V) = (\W_i\cap \W_j)\cap \V \neq \emptyset$. Therefore, $\sim_1$ and $\sim_2$ are equal, so $I/\sim_2=I/
 \sim_1$ is a singleton.
 
 Furthermore, $\W_i\cap \V = \pi((\widehat \W_i \cap \U) \times ]r_i,r_i'[)$; since $\W_i$ is connected and $\U$ portly, $\widehat \W_i\cap \U$ is connected and so is $\W_i\cap \V$. Again by Lemma \ref{lem:connected_union}, $\W\cap \V = \bigcup_{i\in I} (\W_i\cap \V)$ is connected.
 
 \item Assume that $O\in \W$, to reduce to the previous case, it suffices to prove that $\W\setminus \{O\}$ is connected too.
 Consider $\mathcal O = \pi(\{(\sigma,r) ~:~ \sigma\in \Sigma,~ r\in [0,f(\sigma)[\})$ with $f(\sigma):=\sup \{r^*>0 ~|~ \forall r<r^*,~\pi(\sigma,r)\in \W\}$. Note that, since $\W$ is open, $\forall \sigma \in \Sigma,~ (\sigma,f(\sigma))\notin \W$. 
 
 Assume that $f$ is not lower semi-continuous, then (since $\Sigma$ is first countable) there exists some $\sigma\in \Sigma$ and some sequence  $\sigma_n\rightarrow  \sigma$ such that $\ell:=\lim f(\sigma_n)<f(\sigma)$. The sequence $\pi(\sigma_n,f(\sigma_n))$ converges toward $(\sigma,\ell) \in \W$ but $\forall n\in \N, \pi(\sigma_n,f(\sigma_n)) \notin \W$. Contradiction, therefore $f$ is lower semi-continuous.

 We may thus construct some neighborhood $\W_O$ of $O$ in $\W$ such that $\W_O = \bigcup_{i\in J} \pi(\widehat \W_j \times [0,r_j[)$ with $\widehat \W_j$ connected and open. We notice that $\bigcup_{k\in J} \widehat \W_j = \Sigma$ is connected, that the relation $\sim$ on $j$ in Lemma \ref{lem:connected_union} 
 for the families $(\widehat \W_j)_{j\in J}$, $(\W_j)_{j\in J}$ and $(\W_j\setminus \{O\})_{j\in J}$ are equal. Since every $\widehat \W_j, j\in J$ is connected, so are the $\W_j\setminus \{O\}$;    we deduce that $\W_O\setminus \{O\}$ is connected.  
 
 Write $\W\setminus \{O\} = \bigcup_{k\in K\cup \{O\}} (\W_k\setminus \{O\})$ 
 with $(\W_k)_{k\in K}$ the family of connected component of $\W\setminus \{O\}$, in particular $\forall k\in K, \W_k\setminus \{O\}=\W_k$. Since $\W$ is connected, for all $k\in K$, we have $(\W_k\setminus \{O\}) \cap (\W_O \setminus \{O\})=\W_k\cap \W_O \neq \emptyset$. We may apply Lemma \ref{lem:connected_union} again to conclude that $\W\setminus \{O\}$ is connected.
 
 \end{itemize}
\qed\end{proof}
\begin{rem}\label{rem:proof_above_works_quotient} The proof above works for the quotient topology.
 
\end{rem}

 



 
% This suspension method is relatively staightforward in the Riemannian setting since the space of geodesic rays outgoing from given point is naturally identified with $(\SO(n;\RR),\mathbb S^{n-1})$ for a $n$-dimensional Riemannian spaceform. Studying singularities obtained by suspensions in Riemannian spaceform reduces to the study of the hierarchy of spherical suspensions. 
% The situation for more general affine structure is more sophisticated since in general the space of outgoing rays is a projective space. \footnote{The author is not certain of the right definition of projective suspensions.}
% Even in the Riemannian setting, singularities "at infinity", such as cusps, can be obtained via variations of the construction above. \footnote{Such variations are probably avatar of a more general definition of projective suspension.}

% 
%   \subsubsection{Completions}
%  The simplest completions are purely topological such as the Alexandrov or the  Freudenthal completion. 
%  \begin{prop}
%   Let $X$ be a Hausdorff, locally connected topological space and let $\overline X$ be its Feudenthal (or Alexandrov) completion. Then $X$ portly in $\overline X$.
%   \end{prop}
%   \begin{proof}
%    Clearly $X$ is open and dense. 
%   \end{proof}







\section{Branched Covering à la  Fox : spreads, paths and Galois}

Although some of the results of this section may be proven with more generality, all topological spaces are Hausdorff, locally path-connected and first countable, unless explicitely stated otherwise.

In this section we recall fundamentals on spreads and branched coverings with some enrichments. Considering those will be needed later on, we prefer to lay out solid foundation and provide most of the proofs. Some results presented in section \ref{sec:spread_base}, and \ref{sec:branched_base}, though not found in the literature by the author, are probably ``folkloric''. Our presentation is mostly redundant with \cite{MR3060537,MR3309706,MR2599123,MR2198590} but our different purposes leads to divergences. For instance, to our knowledge, Lemma \ref{lem:subcovering_general} and Proposition \ref{prop:branched_covering_composition} are new.


To the author knowledge, the path representation of maximal branched coverings\footnote{The path representation is not used in the second part  of the work on singular $(G,X)$-manifolds but fits nicely both in the present study of branched coverings and the subsequent study of developing maps of singular $(G,X)$-manifolds. It's the author belief that this will prove useful for future studies.} presented in section \ref{sec:maximal_branched_path} as well as the topological caracterisation of Galoisian maximal branched coverings presented in section \ref{sec:galoisian_branched} are both new. The former is a natural extension of the usual description of the universal covering as a space of paths to branched coverings. The key point being the introduction of the right homotopy equivalence relation. An extension of the classical semi-local simple connectedness notion based on the same homotopy equivalence is then introduced and shown to caracterise branching locus giving rise to Galoisian coverings is described  in section \ref{sec:galoisian_branched}. Example and counter-examples are given throughout the section.

\subsection{Preliminaries on spreads}\label{sec:spread_base}


\begin{defi}[Category of spreads]\label{def:category_spreads} A spread $X\xrightarrow{p} Y$ is a continuous map such that the topology of $X$ is generated by the connected components of the preimages of open subsets of $Y$.

We define the category of spreads $\CatSpread$ whose objects are spreads and morphisms in $\mathrm{Hom}(X\xrightarrow{p} Y, X'\xrightarrow{p'} Y')$ are couples of continuous functions $(f,g)$  such that the following diagram commutes

$$\xymatrix{
X\ar[d]^p \ar[r]^g & X' \ar[d]^{p'} \\ Y \ar[r]^{f} & Y'
}.$$

\end{defi}

\begin{rem}\label{rem:spread_open_respectively_co} Let $X\xrightarrow{p} Y$ be a spread, 
with $\U,\V$ open of $X$  and $Y$ respectively such that $p(\U)\subset \V$, the (co)restriction $p_{|\U}^{|\V}$ is a spread.
\end{rem}
\begin{defi}\label{def:continuous_map_ordinary_locus}
  Let $X\xrightarrow{p} Y$ be a continuous map, the  ordinary locus $\ord_p(Y)$ of $Y$  is the set of point $y\in Y$  for which there exists an open neighborhood $\U$ evenly covered by $p$ ie such that $p$ maps each connected component of $p^{-1}(\U)$ homeomorphically onto $\U$. The ordinary locus $\ord_p(X)$ of $X$ is the preimage of the ordinary locus of $Y$.
  
  We will write $\ord$ instead of $\ord_p$ when there is no ambiguity.
\end{defi}
\begin{rem}\label{rem:ordinary_locii_open_but}
 The ordinary locii are open but may not be connected. The only obstruction for the restriction $\ord(X)\rightarrow \ord(Y)$ to be a covering is the connectedness of $\ord(X)$.
\end{rem}
\begin{rem}\label{rem:singular_on_singular}
 Let $X\xrightarrow{p}Y$ be a surjective continuous map. Then $$p(X\setminus \ord(X))  = Y\setminus \ord(Y).$$
\end{rem}


\begin{lem}\label{lem:spread_partial_covering} Let $X\xrightarrow{p} Y$ be a continuous map inducing a covering $X_1\rightarrow Y_1$ for some $X_1\subset X$ and $Y_1\subset Y$ with $X_1$ portly and $Y_1$ open. 
 
 Then $Y_1\subset \ord(Y)$, $X_1 \subset \ord(X)$ and $p^{-1}(Y_1)=X_1$. 
\end{lem}



\begin{proof}
 Let $y\in Y_1$ and let $\U\subset Y_1$ be an open neighborhood 
 of $y$ evenly covered by $p_{|X_1}^{|Y_1}$ ie
 $$p^{-1}(\U)\cap X_1 = (p_{|X_1}^{|Y_1})^{-1}(\U)=\bigsqcup_{i\in I} \V_i$$
 with $\V_i$ open connected and $p_{|\V_i}^{|\U_i}$ homeomorphisms for all $i\in I$. 
 
 Assume by contradiction there exists some $x\in p^{-1}(\U)\setminus X_1$, consider some connected open subset  $\V$ such that $x\in \V\subset p^{-1}(\U)$. Since $X_1$ is portly, $\V\cap X_1$ is connected and dense in $\V$, thus  $\V\cap X_1 \subset \V_i$ for some $i\in I$ and $x\in \overline{\V_i}$. Let $x':= (p_{|\V_i}^{|\U})^{-1}\circ p(x)$ and let $\W,\W'$ disjoint open neighborhoods of $x$ and $x'$ respectively with $\W'\subset \V_i$. Since $p(x) \in p(\W')$ which is open, there exists some open neighborhood $\W''\subset \W$  of $x$ such that $p(\W'')\subset p(\W')$. Therefore, $\emptyset \neq p(\W''\cap \V_i) \subset p(\W')=p(\W'\cap \V_i)$. This contradicts injectivity on $p_{|\V_i}$, hence $p^{-1}(\U)\subset X_1$ and $\U$ is evenly covered by $p$. In particular, $p^{-1}(Y_1) = X_1$.
 
 Therefore, $Y_1\subset \ord(Y)$ and $\ord(X) = p^{-1}(\ord(Y)) \supset p^{-1}(Y_1) = X_1$
\qed\end{proof}

Though not the core of the present work, for pedagogical purposes we provide counter-examples to the previous Lemma we feel help understand some caveats when dealing with spreads.
\begin{cex}[Openess of $Y_1$]\label{cex:openess}

 Consider $X = \RR\times \{0\} $ and $Y = \RR^2$ with $p$ the inclusion map. $p$ is a continuous map (it is even a spread) and induces a covering $X\rightarrow X \subset Y$ but $\ord(Y)=\emptyset$.
\end{cex}
\begin{cex}[Portlyness of $X_1$]\label{cex:portlyness} \ 
\begin{itemize}
 \item If $X_1$ is not open, then take $x \in X_1\setminus Int(X_1)$, for neighborhood $\V$ of $x$ small enough, the set $\V\cap X_1$ is mapped homeomorphically to some open neighborhood $\U$ of $p(x)$. For $\V$ small enough, $p(\V)\subset \U$ thus $p(\V\setminus X_1) \cap p(\V\cap X_1)\neq \emptyset$. In otherword, $p$ is not injective in any neighborhood of $x$. 
 In particular, $p(x)\notin \ord(Y)$ and  $x\notin \ord(X)$.
 \item  Take $X\subset Y=\RR^2$ with the Euclidean norm,
 $X=\{x~|~\|x\|\geq 1\} \cup \{\lambda x ~|~ \lambda \in ]0,1], \|x\|=1, x\in \mathbb Q^2\}$.  Endow $Y$ with its usual topology but $X$ with the coarsest topology making the natural injection $X\rightarrow Y$ a spread : the subset $\{\lambda x ~|~ \lambda \in ]0,1[, \|x\|=1, x\in \mathbb Q^2\}$ is homeomorphic to $\coprod_{i\in \N} ]0,1[$. 
  Define piecewise $p(x) = \frac{x}{\|x\|^2}$ if $\|x\|\leq 1$ and $p(x) = x$ if $\|x\|>1$. Take $X_1=\{x~|~\|x\|>1\}=Y_1$. Notice that $X_1\xrightarrow{p} Y_1$ is an homeomorphism hence a covering, $X$ is Hausdorff, second countable and locally path-connected, $X\xrightarrow{p}Y$ is a spread, $Y_1$ is open and $X_1$ is open locally connected in $X$.  However, $\ord(Y)=\emptyset$ since for every open neighborhood $\U$ of some $y\in p(X)$, the preimage $p^{-1}(\U)$ contains points $x$ with $\|x\|<1$ but $p$ is not a local homeomorphism around such a point. 
\end{itemize}

\end{cex}





\begin{lem}\label{lem:ord_iso} Let $X_i\xrightarrow{p_i} Y_i$ be spreads for $i\in \{1,2\}$ and let $(f,g)$ be a spread isomorphism from $X_1\rightarrow Y_1$ to $X_2\rightarrow Y_2$. Then $g(\ord(X_1)) = \ord(X_2)$ and $f(\ord(Y_1)) = \ord(Y_2)$.
 
\end{lem}
\begin{proof}
 Let $\U$ be an open connected neighborhood of $p_1(x)$ evenly covered by $p_1$, let $(\hat\U_i)_{i\in I}$ the connected components of $p_1^{-1}(\U)$ and define $\V:=f(\U)$ and $\hat \V_i:= g(\hat \U_i)$ for $i\in I$. 
 Since $f$ and $g$ are bijective and $f\circ p_1=p_2\circ g$, we have $$p_2^{-1}(\V)=g\circ p_1^{-1}\circ f^{-1}(\V) = g\circ p_1^{-1}(\U) = \bigsqcup_{i\in I} \hat \V_i$$   
 and each $\hat \V_i$ is both open and connected.
 
 Furthermore, $p_{2|\hat\V_i}^{|\V}\circ g^{|\hat\V_i}_{|\hat\U_i} = f_{|\U}^{|\V}\circ p_{1|
 \hat\U_i}^{|\U}$ for all $i\in I$ and $g^{|\hat\V_i}_{|\hat\U_i}$ as well as $f_{|\U}^{|\V}\circ p_{1|
 \hat\U_i}^{|\U}$ are homeomosphims; therefore, for all $i\in I$, $p_{2|\hat\V_i}^{|\V}$ is an homeomorphism.
 Finally, $\V$ is evenly covered by $p_2$. The result follows.
 \qed\end{proof}

\begin{defi}[Complete spread, spread completion]\label{def:complete_spread_spread_completion} Let  $X\xrightarrow{p} Y$ be a spread. For $y\in \overline{p(X)}$, consider the set $X_y$ of maps $$\chi : \{\U\subset Y ~ \text{ open and connected neighborhood of }y\} \rightarrow \{\V\subset X ~ \text{open connected}\}$$ increasing for the inclusion and such that  $\chi(\U)$ is a connected component of $p^{-1}(\U)$ for all $\U$.
The spread $X\xrightarrow{p} Y$ is complete if for all $y\in \overline{p(X)}$ and all $\chi \in X_y$, the intersection of $\chi(\U)$ for $\U$ going through all open connected neighborhood of $y$ is non empty.


A completion of a spread $X\xrightarrow{p} Y$ is a spread $X'\xrightarrow{p'} Y$ together with a spread morphism $(Id_Y,\iota)$
$$\xymatrix{
X\ar[d]^p \ar[r]^\iota & X' \ar[d]^{p'} \\ Y \ar@{=}[r] & Y
} $$
such that $\iota$ is an embedding, $X'$ is complete and the image of $\iota$ is portly.
\end{defi}
Recall that a functor $T$ is fully faithful if the induced function $\Hom(X,Y) \rightarrow \Hom(T(X),T(Y))$ is bijective.
\begin{thm}[\cite{MR0123298,MR2198590}] \label{theo:Fox_completion_functor}
 Every spread admits a unique completion up to isomorphism, furthermore the completion of spread is a fully faithful functor  from $\CatSpread$ to itself.
\end{thm}



Let $X\xrightarrow{p} Y$ be a spread and let $b\in Y$. Let $\boldsymbol{\U}$ be a connected neighborhood basis of $b$. For $\U \in \boldsymbol{\U}$, let $X_\U$ be the space of connected components of $p^{-1}(\U)$ endowed with the discrete
 topology. Define bonding maps $X_\U\xrightarrow{\pi_{\U\V}} X_\V$ by setting for $\U\subset \V$ the image $\pi_{\U\V}(\widehat\U)$ of $\widehat \U\in X_\U$ to be the unique $\widehat \V\in X_\V$ such that $\widehat \U \subset \widehat \V$. For $x\in p^{-1}(b)$ and $\U\in \boldsymbol\U$, denote by $x_\U$ the connected component of $x$ in $p^{-1}(\U)$.
 We have a natural continuous map $$\fonctionn{p^{-1}(b)}{\varprojlim_{\U\in \boldsymbol{\U}}X_\U}{x}{(x_\U)_{u\in \boldsymbol \U}}.$$

\begin{prop}[\cite{MR2198590}] \label{prop:spread_fiber}
 Let $X\xrightarrow{p} Y$ be a spread and let $b\in Y$. Let $\boldsymbol{\U}$ be a connected neighborhood basis of $b$. 
 
 Then, the  natural map $p^{-1}(b) \rightarrow \varprojlim_{\U\in \boldsymbol{\U}}X_\U$ 
 is injective open. 
 Furthermore, it is an homeomorphism if and only if the spread $X\xrightarrow{p}Y$ is complete.

\end{prop}
\begin{proof}
 Denote by $\phi$ the natural map
 $p^{-1}(b) \xrightarrow{\phi}\varprojlim X_\U$. For every $x\in p^{-1}(b)$, since $p$ is a spread, a neighborhood basis of $x$ is given by the connected components of $x$ in each $p^{-1}(\U)$ for $\U\in \boldsymbol{\U}$. 
 If two elements $x_1,x_2 \in X$ have the same image by $\phi$ then $x_1,x_2 \in \phi(x_1)_\U$ for all $\U\in \boldsymbol{\U}$. Hence, $x_1,x_2$ have a common neighborhood basis. The space $X$ so $x_1=x_2$; we deduce that $\phi$ is injective.
 

 Notice that for all $\U\in \boldsymbol{\U}$ and all $\V\in X_\U$, $$\phi(\V\cap p^{-1}(b)) = \pi_\U^{-1}(\V)$$
 {with $\pi_\U$ the natural map $\varprojlim_{\W\in \boldsymbol{\U}} X_\W \xrightarrow{~~~} X_\U$.}
 On the one hand, such $\pi_\U^{-1}(\V)$ form a basis of the topology of $\varprojlim X_\U$. On the other hand, since $p$ is a spread, such $\V\cap p^{-1}(\V)$ form a basis for the topology of $p^{-1}(b)$. 
 Therefore, $\phi$ is open.
 
  Let $\hat{\boldsymbol{\U}}=\left(\hat\U\right)_{\U\in \boldsymbol{\U}} \in \varprojlim_{i\in I}X_i$. The intersection $\bigcap_{\U\in \boldsymbol{\U}}\hat\U$ contains at most a point $x \in p^{-1}(b)$  and for such an $x$, we have $\phi(x)=\hat{\boldsymbol{\U}}$.  Reciprocally, if $\hat{\boldsymbol \U} \in \phi(p^{-1}(b))$ then  $\phi^{-1}(\hat{\boldsymbol{\U}}) \in \bigcap_{\hat\U\in \hat{\boldsymbol{\U}}}\hat\U$. 
  
  Therefore, $\phi$ is surjective if and only if $\forall \hat{\boldsymbol \U}\in \varprojlim X_\U, ~ \bigcap_{\hat\U\in \hat{\boldsymbol{\U}}}\hat\U \neq \emptyset$. Which is merely a predantic way to write that the spread $X\xrightarrow{p}Y$ is complete. 
\qed\end{proof}
\begin{cor}[\cite{MR2198590}]\label{cor:totdiscont} Let $X\xrightarrow{p} Y$ be a spread, then the fibers of $p$ are totally discontinuous. 
\end{cor}






\subsection{Primer on Branched covering} \label{sec:branched_base}

The object of this section is to recall fundamentals on branched covering from definition, 
 to path-lifting characterisation, basic properties and composition. We end this section 
 with existence and uniqueness of universal branched covering with constrained 
 branching locus.  This will be the basis for futher analysis.

As before, unless explicitely stated otherwise, all topological spaces are Hausdorff, 
locally path-connected and first countable.
In the present section as well as the following, we will make intensive use of the following Lemmata which, for brievety sake, we will not refer to systematically.


\begin{lem}\label{lem:path_covering} Let $X$ be a connected first countable, locally path-connected, Hausdorff topological space and let $\U\subset X$ be a portly subset. 
For all $a\in \U$ and all $b\in X$, there exists a path $\gamma: [0,1]\rightarrow X$  from $a$ to $b$ such that $\gamma([0,1[)\subset \U$.
\end{lem}
\begin{proof}
 Let $(a,b) \in \U\times X$ and let $(\V_n)_{n\in \N}$ be a decreasing base of open connected neighborhood of $b$. Since $\U$ is portly, in particular $\U$ is dense and there exists a sequence $(b_n)_{n\in \N}$ such that  $\forall  n\in \N, b_n \in \V_n\cap \U$. Since $\U$ is portly and $X$ is connected then $\U$ is connected and we can choose a path $\gamma_0$ in $\U$ from $a$ to $b_0$. Then for $n\in \N$, $b_n$ and $b_{n+1}$ are both in $\U\cap \V_n$ and since $\U$ is portly and $\V_n$ open connected, $\U\cap \V_n$ is open connected (hence path-connected) and we can choose a path $\gamma_{n+1}$ in $\U\cap \V_n$ from $b_n$ to $b_{n+1}$. The concatenation of the sequence of path $(\gamma_n)_{n\in \N}$ is a path $\gamma:[0,1[\rightarrow \U$ such that $\lim_{1^-}\gamma = b$. The result follows.
\qed\end{proof} 

\begin{lem}\label{lem:portly_first_countable} Let $X$ be a connected first countable, locally path-connected Hausdorff topological space and let $X_1$ a subset containing a portly subset. Then, $X_1$ is first countable, locally path connect Hausdorff.
\end{lem}
\begin{proof} To begin with, $X_1$ inherits Hausdorff and first countable properties from $X$ as any subspace. It remains to show local path-connectedness.
Take $\U\subset X_1$ some portly subset of $X$. 
 Let $x\in X_1$ and let $\V_1$ be an open neighborhood of $x$ in $X_1$. Choose some connected open $\V$ of $X$ such that $x\in \V_2:=\V\cap X_1 \subset \V_1$. 
 It suffices to show that $\V_2$ is path-connected, so we choose some $a,b\in \V_2$.
 Since $\U$ is portly in $X$, the intersection $\U\cap \V$ is portly in $\V$. Take any $c\in \U\cap \V$ and apply Lemma \ref{lem:path_covering} to construct paths $\gamma_1$ from $c$ to $a$ and $\gamma_2$ from $c$ to $b$ such that $\gamma_{i}([0,1[)\subset \U\cap \V\subset \V_2$ for $i\in \{1,2\}$. Concatenating the reverse of the former with the latter, we obtain a path from $a$ to $b$ contained in $\V_2$. 
 
\qed\end{proof}





\begin{defi}[Branched covering]\label{def:branched_covering}  A {\it branched covering} is a complete surjective spread $X\rightarrow Y$, with $X$ and $Y$ connected, whose ordinary locii in $X$ and $Y$ are portly subsets of $X$ and $Y$ respectively. We call $Y$ the {\it base space}, $X$ the {\it total space} and the inverse image of a point $a$ of $Y$ is called the {\it fiber} above $a$.
 
 We say that $X\rightarrow Y$ is branched over $Y\setminus \ord(Y)$. The latter is the branching locus of $X\rightarrow Y$ and for $S\subset Y$, we say that $X\rightarrow Y$ is possibly branched over $S$ if $S$ contains its branching locus.

 \end{defi}
 


 
 
 \begin{rem}\label{rem:branching_locii_branched_covering} The branching locii of a branched covering are skeletal thus closed with empty interior.
 \end{rem}
\begin{rem}\label{rem:ordinary_part_branched_covering}
 The ordinary part $\ord(X)\rightarrow \ord(Y)$ of a branched covering $X\rightarrow Y$ is a covering. Indeed, $X$ is open connected and $\ord(X)$ is portly so $\ord(X)=\ord(X)\cap X$ is connected. 
\end{rem}
\begin{lem}\label{lem:portly_preimage_branched} If $X\xrightarrow{p} Y$ is a branched covering, then the preimage of  portly subset of $Y$ is portly in $X$.
\end{lem}
\begin{proof} Let $\U$ be a portly subset of $Y$.
The map $\ord_p(X) \xrightarrow{p} \ord_p(Y)$ is a local homeomorphism  and $\ord_p(Y)$ is portly in $Y$ thus $p^{-1}(Y\cap \ord_p(Y))$ is portly in $\ord_p(X)$. Since $\ord_p(X)$ is portly in $X$, then $p^{-1}(Y\cap \ord_p(Y))$ is portly in $X$. 
The subset $p^{-1}(Y)\subset X$ is open and contains a portly subset of $X$, it is thus also portly.
 
\qed\end{proof}



 We begin by a branched covering version of the lifting property for complete spreads ("Corollary of the extension Theorem" in Fox \cite{MR0123298}).
\begin{lem}\label{lem:complete_path_carac} 
 Let $X\xrightarrow{p} Y$ be a spread with $X$ connected which induces a covering $X_1\rightarrow Y_1$ for some $X_1$ and $Y_1$ portly subsets of $X$ and $Y$ respectively. Then the following are equivalent:
 \begin{enumerate}[(i)]
  \item $X\rightarrow Y$ is a branched covering;
  \item $X\rightarrow Y$ is a complete spread
  \item for every Hausdorff locally connected topological space $Q$ and every homotopy $H:Q\times [0,1]\rightarrow Y$ such that $H(Q\times [0,1[)\subset Y_1$, every partial lift 
  $\hat H : Q \times [0,1[ \rightarrow X_1$ of $H$ extends continuously to some lift $Q \times [0,1] \rightarrow X$ of $H$;
  \item for every path $\gamma:[0,1]\rightarrow Y$ such that $\gamma([0,1[) \subset Y_1$, every partial lift 
  $\hat \gamma : [0,1[ \rightarrow X_1$ extends continuously to some lift $[0,1] \rightarrow X$ of $\gamma$.
  
 \end{enumerate}
\end{lem}
\begin{proof}
To begin with, by Lemma \ref{lem:spread_partial_covering} we have $X_1\subset \ord(X)$ and  $Y_1\subset \ord(Y)$ so both $\ord(X)$ and $\ord(Y)$ are portly.



 One has $(i)\Rightarrow (ii)$ by definition, $(ii)\Rightarrow (iii)$ is a direct consequence of the corollary of Fox mentionned above and $(iii)\Rightarrow (iv)$ is trivial (it suffices to take $Q$ a singleton). 
 Assume $(iv)$ and consider some $y_0 \in Y_1$. Since $Y_1$ is portly in $Y$, for every $y\in Y$ by Lemma \ref{lem:path_covering} there exists a path $\gamma:y_0\rightsquigarrow y$ such that $\gamma([0,1[)\subset Y_1$. Since $X_1\rightarrow Y_1$ is a covering, $\gamma_{|[0,1[}$ lifts to some $\hat\gamma :[0,1[\rightarrow X_1$. By assumption, $\hat\gamma$ admits a limit $x$ at $1$ and by continuity, $p(x)=y$. Therefore, $p$ is surjective.  
 Consider $\overline X \rightarrow Y$ the completion of $X\rightarrow Y$; without loss of generality, we may assume that $X\subset \overline X$ and also denote by $p$ the map $\overline X\rightarrow Y$.   Notice that since $X_1$ is portly in $X$ and $X$ is portly in $\overline X$ then $X_1$ is portly in $\overline X$. Let $x \in \overline X$ and let $\hat y_0$ be some lift of $y_0$ in $X_1$. Let  $\hat \gamma$ be some path from $\hat y_0$ to $x$ in $\overline X$ such that $\hat\gamma([0,1[)\subset X_1$. Since $\hat\gamma_{|[0,1[}$ is a lift in $X_1$ of  the path $p\circ \hat \gamma_{|[0,1[} \subset Y_1$, assuming $(iv)$ we deduce that $\hat\gamma_{|[0,1[}$ has limit at $1$ in $X$. Therefore $x = \hat\gamma(1)= \lim_{1^-} \hat\gamma_{|[0,1[} \in X$. Finally, $\overline X\subset X$ hence $\overline X = X$ and $X\rightarrow Y$ is complete. 
\qed\end{proof}



 


\begin{lem}[\cite{MR2198590}]\label{lem:branched_covering_restriction}
 Let $X\xrightarrow{p} Y$ be a branched covering, let $\U$ be some open connected subset of $Y$ and let $\V$ be a connected component of $p^{-1}(\U)$. Then $\V\xrightarrow{p} \U$ is a branched covering.
\end{lem}
\begin{proof} To begin with, sing $\U,\V$ are open, the restriction $p_{|\U}^{|\V}$ is a spread.
 Since $\ord(X)$ is dense, we can choose some $\hat x\in \ord(X)\cap \V$; define $x:=p(\hat x)$. Since $Y$ is locally pathwise connected and since $\ord(Y)$ is portly, then $\ord(Y)\cap \U$ is path-connected.  We then can and do consider a path $\gamma:[0,1]\rightarrow \ord(Y)\cap\U$ from $x$ to some $y\in \ord(Y)\cap \U$. Consider the unique lift $\hat\gamma$ of $\gamma$ in $\ord(X)$ such that $\hat\gamma(0)=\hat x$. On the one hand, $\hat\gamma([0,1])\subset p^{-1}(\U)$; on the other end, $\hat\gamma([0,1])$ is connected. Therefore, $\hat \gamma$ is actually a path in $\V$ and $\gamma([0,1])\subset p(\V)$. 
 The path $\gamma$ is arbitrary, so $p(\ord(X)\cap \V)=\ord(Y)\cap \U$ and  $$\ord(X)\cap \V \xrightarrow{ ~~p~~} \ord(Y)\cap \U$$ is a covering. 
 
 Since $X\xrightarrow{p} Y$ is a branched covering, it satisfies item $(iv)$ of Lemma \ref{lem:complete_path_carac}. Since $X,Y$ are locally path-connected, this item is clearly satisfied by $\V\xrightarrow{p}\U$. Finally, $\V\rightarrow \U$ is a branched covering 
 
\qed\end{proof}


\begin{cor}\label{cor:spread_connected_image}
   Let $X\xrightarrow{p} Y$ be a branched covering. Then $p$ is open, furthermore for every connected open $\U\subset Y$ and every connected component $\V$ of $p^{-1}(\U)$ 
   the map $\V\xrightarrow{p} \U$ is surjective.
\end{cor}
\begin{proof}
Surjectivity is part of the definition of Branched covering and thus follows from Lemma \ref{lem:branched_covering_restriction}. 
  For all such $\U,\V$ we thus have $p(\V)=\U$. Since $p$ is a spread, such $\V$ is a basis of the topology of $X$ hence $p$ is open.
\qed\end{proof}

\begin{cor} \label{cor:subcovering_portly}
  Let $X\xrightarrow{p} Y$ be a branched covering, let $\U$ be some portly subset of $Y$ then $p^{-1}(\U)\xrightarrow{p}\U$ is a branched covering.
\end{cor}
\begin{proof}  From Lemma. \ref{lem:portly_preimage_branched} the subset $p^{-1}(\U)$ is portly in $X$ and since $X$ is connected so is $p^{-1}(\U)$. The result then follows from Lemma \ref{lem:branched_covering_restriction}.
\qed\end{proof}


\begin{lem}\label{lem:subcovering_general}
 Let $X\xrightarrow{p} Y$ be a branched covering, let $Y_1\subset Y$ containing a portly subset of $Y$ then $p^{-1}(Y_1)\xrightarrow{~p~}Y_1$ is a branched covering.
\end{lem}
\begin{proof} Denote $X_1:=p^{-1}(Y_1)$ and $q:= p_{|X_1}^{|Y_1}$. Take some $\U\subset Y_1$ portly in $Y$.
\begin{itemize}

 \item By Lemma \ref{lem:portly_preimage_branched}, the subset $\widehat \U :=p^{-1}(\U)\subset X_1$ is portly in $X$.  
  \item By Lemma \ref{lem:portly_first_countable},  $X_1,Y_1$ are first countable, Hausdorff, locally path-connected.
 \item Let $\V$ an open subset of $X_1$ and $x\in \V$. Consider $\V_1$ open subset of $X$ such that $\V=\V_1\cap X_1$. Since $p$ is a spread, there exists some open subset $\U_1\subset Y$ and some connected component $\W_1$ of $p^{-1}(\U_1)$ such that $x\in \W_1 \subset \V_1$. Since $\widehat \U\subset X_1$ is portly in $X$, the subset  $\W_1\cap \widehat \U$ is connected and dense in $\W_1\cap X_1$. It follows that $\W_2:=\W_1\cap X_1$ is connected and thus a connected component of $q^{-1}(\U_1\cap Y_1)$.
 Moreover, $\W_2$ is open in $X_1$ and  $x\in \W_2\subset \V$. 
 Therefore, $X_1 \xrightarrow{q} Y_1$ is a spread. 
 \item Define $r:=p_{|\widehat \U}^{|\U}$. By Corollary \ref{cor:subcovering_portly}, $\widehat \U\xrightarrow{r}\U$ is a branched covering, in particular $\ord_{r}(\widehat \U)\xrightarrow{p} \ord_r(\U)$ is a covering. 
 Furthermore $\ord_r(\widehat \U)$ is portly in $\widehat \U$ which is portly in $X_1$, so $\ord_r(\widehat \U)$ is portly in $X_1$. The same way, $\ord_r(\U)$ is portly in $Y_1$.
 \item Let $\gamma : [0,1] \rightarrow Y_1$ such that $\gamma([0,1[) \subset \ord_r(\U)$ and let $\hat \gamma : [0,1[\rightarrow \ord_r(\U)$ be a lift of $\gamma$. Since $X\xrightarrow{p} Y$ is a branched covering, $\hat\gamma$ extends continuously at $1$ and $$p\circ \hat \gamma(1)=\lim_{1^-}p\circ \hat\gamma =\lim_{1^-}\gamma =\gamma(1)\in Y_1.$$
 Consequently, $\hat\gamma(1)\in p^{-1}(Y_1)=X_1$.
\end{itemize}

We checked the assumptions of Lemma \ref{lem:complete_path_carac} as well as point $(iv)$ of the equivalence. We thus conclude that $X_1\xrightarrow{q} Y_1$ is a branched covering.
\qed\end{proof}



\begin{prop} \label{prop:branched_covering_composition}Let $X\xrightarrow{p} Y \xrightarrow{q} Z $ be two branched coverings.   Then the following are equivalent: 
\begin{enumerate}[(i)]
 \item $X\xrightarrow{q\circ p} Z$ is a branched covering,
     \item $q(Y\setminus \ord_p(Y))$  is skeletal in $Z$.
 \end{enumerate}


\end{prop}
\begin{proof}
To begin with, composition of spread is a spread so $q\circ p$ is a spread. 
 
 \begin{itemize}
 \item $\underline{(i)\Rightarrow (ii)}$.  Let $\U\subset \ord_{q\circ p}(Z)$ be a connected open  evenly covered subset and define 
 $$ (q\circ p)^{-1}(\U) = \bigsqcup_{i\in I} \hat \U_i   \quad ; \quad \forall i\in I,~\V_i := p(\hat \U_i) $$
respectively the connected component decomposition of the preimage of $\U$ and their image by $p$. 
For all $i\in I$, on the one hand $\V_i$ is open connected by Corollary 
\ref{cor:spread_connected_image}, on the other hand 
$(q\circ p)_{|\hat \U_i}^{|\U}$ is injective, furthermore $q$ is open and $q(\V_i) = \U$.
Therefore, $q_{|\V_i}^{|\U}$ is bijective open thus an homeomorphism.

Denote by $\hat \V_i$ the connected component of $q^{-1}(\U)$ containing $\V_i$. 
If $\V_i\neq \hat \V_i$ we may find some $j\in I$ such that $\hat \V_i = \hat\V_j$ 
and such that $\V_i\cap \V_j \neq \emptyset$; as well as some 
$y\in (\overline \V_i\setminus \V_i) \cap \V_j$. 
An argumentation alike the one of Lemma \ref{lem:spread_partial_covering}
yields a contradiction showing that in fact $\forall i\in I, \V_i = \hat \V_i$.  
We conclude that $q^{-1}(\U)\xrightarrow{q} \U$ and 
for all $i\in I$, $p^{-1}(\V_i)\xrightarrow{p} \V_i$ are coverings; 
with the preimages being disjoint union of connected components 
among the $\V_i$ and $\hat\U_i$ respectively. 
As a consequence 
$$\ord_{q\circ p}(X)\xrightarrow{p} p(\ord_{q\circ p}(X))=q^{-1}(\ord_{q\circ p}(Z)) \xrightarrow{q} \ord_{q\circ p}(Z)$$
is a composition of coverings.
Therefore: 
$$q(Y\setminus \ord_p(Y)) \subset q(Y\setminus q^{-1}(\ord_{q\circ p}(Z))) 
= Z\setminus \ord_{q\circ p}(Z)$$
The later is skeletal, thus so is the former.




 
  \item $\underline{(ii)\Rightarrow (i)}$. Consider $Y_0 :=\ord_p(Y)\cap \ord_q(Y)$, and consider $Z_1 := Z\setminus q(Y\setminus Y_0)$. We note that $$Z_1= Z\setminus (q(Y\setminus \ord_p(Y)) \cup q(Y\setminus \ord_q(Y)) =  \ord_q(Z) \setminus q(Y\setminus \ord_p(Y))$$ 
 Since $\ord_q(Z)$ is portly and $q(Y\setminus \ord_p(Y))$ is skeletal, $Z_1$ is portly. Since $q$ is a covering, it is a surjective local homeomorphism so $Y_1:=q^{-1}(Z_1)$ is portly in $Y$. Furthermore, $Y_1\subset Y_0$.
 Since $r:=q_{|\ord_q(Y)}^{|\ord_q(Z)}$ is a covering then its (co)restriction is also $q_{|Y_1}^{|Z_1} = r_{|r^{-1}(Z_1)}^{|Z_1}$ is a covering.
 The same way, $p_{|X_1}^{|Y_1}$ is a covering with $X_1:= \left(p^{|\ord_p(Y)}_{|\ord_p(X)}\right)^{-1}(Y_1)$ and $X_1$ is portly. 
 Therefore, $(q\circ p)_{|X_1}^{|Z_1}$ is a a composition of covering hence a covering. 
 
 Applying Lemma \ref{lem:complete_path_carac}, we see that item $(iv)$ is stable by composition of branched covering. Therefore, $q\circ p$ is a branched covering.
 \end{itemize}
\qed\end{proof}





\begin{lem}\label{rem:completion_covering} Let $X\xrightarrow{p} Y$ be a branched covering and $\U\subset Y$ some portly subset. The spread $X\rightarrow Y$ is isomorphic to the completion of $p^{-1}(\U)\rightarrow Y$.
\end{lem}
\begin{proof}
 $X\rightarrow Y$ is a completion of the spread $p^{-1}(\U)\rightarrow Y$ which is unique up to isomorphism by Theorem~\ref{theo:Fox_completion_functor}.
\qed\end{proof}




Recall that a connected, locally path-connected Hausdorff topological space $X$ is {\it semi-locally simply connected} if every $x\in X$ admit a neighborhood $\U$ such that every loop $\gamma$ in $X$ of base point $x$ is trivial in $X$. By standard results \cite{MR1700700}, it is a necessary and sufficient condition for such a topological space $X$ to admit a universal covering.

\begin{prop}\label{prop:universality} Let $Y$ be a connected, locally path-connected, Hausdorff topological space.  Let $S$ be a skeletal subset of $Y$.

If $Y\setminus S$ is semi-locally simply connected then  there exists a covering $\widetilde Y^S$ of $Y$ possibly branched over $S$ which is maximal among such branched covering.
Furthermore, $\widetilde Y^S$
 is unique up to isomorphism and universal in the sense that for any covering $X\xrightarrow{p} Y$ possibly branched over $S$, there exists a covering $\widetilde Y^S\rightarrow X$ possibly branched over $p^{-1}(S)$.
\end{prop}
\begin{rem}\label{rem:one_could_sum_up}
One could sum up the proof below by saying Fox completion induces a category equivalence between the category of covering of $Y\setminus S$  and the category of covering of $Y$ possibly branched over $S$.
\end{rem}
\begin{proof} 


 Assume $Y\setminus S$ is semi-locally simply connected. Then define $\widetilde Y^S$ the completion of spread $\widetilde {Y\setminus S} \rightarrow Y$ with $\widetilde {Y\setminus S}$ the universal covering of $Y\setminus S$. For any branched covering $X\xrightarrow{p}Y$ possibly branched over $S$, since $Y\setminus S\subset \ord(Y)$ is portly, $p^{-1}(Y\setminus S)$ is portly (hence connected) and $p^{-1}(Y\setminus S) \rightarrow Y\setminus S$ is a covering. Hence, there exists a covering morphism $\widetilde {Y\setminus S} \rightarrow p^{-1}(Y\setminus S) $. Since the completions of 
 $p^{-1}(Y\setminus S) \rightarrow Y$ and $\widetilde {Y\setminus S}\rightarrow Y$ are respectively $X$ and $\widetilde Y^S$, by Theorem \ref{theo:Fox_completion_functor} we obtain a morphism of branched covering 
 $\widetilde Y^S \rightarrow X$ which is itself a branched covering and its ordinary locus in $X$ contains $p^{-1}(Y\setminus S)$. Finally, $\widetilde Y^S \rightarrow X$ is possibly branched over $p^{-1}(S)$. 
 
 If $X\rightarrow Y$ is a maximal  covering possibly branched over $S$ then the induced covering $p^{-1}(Y\setminus S)\rightarrow Y\setminus S$  is maximal among covering of $Y\setminus S$ hence its universal covering. Therefore, $X$ is isomorphic to $\widetilde Y^S$.
\qed\end{proof}
\begin{rem}\label{rem:considering_maximal_covering_possibly}
 Considering the maximal covering $\widetilde Y^S$ of some $Y$ possibly branched over some subset $S$, there is an ambiguity: the branching locus of $\widetilde Y^S$ might be smaller than $S$. One can consider the example where $Y=\mathbb S^2$ the 2-dimensionnal sphere and $S$ is a singleton, since the complement of a point is simply connected, the maximal covering of $\mathbb S^2$ possibly branched above $S$ is trivial and the branching locus is empty. 
 \end{rem}


\begin{cex}[to Proposition \ref{prop:branched_covering_composition}]\label{cex:proposition_prop_branched_covering}

Taking notions from example \ref{ex:quotient_base},  Consider $Z=\DD$, $Y=\DD_\infty$ with $Y\xrightarrow{q} Z$ the quotient map which is a branched covering. Consider $Y_0 = \{(r,\theta) ~|~ \theta \in ]0,1[\}$ and choose $(r_n,\theta_n)_{n\in \ZZ}$ a countable dense family in $Y_0$. Consider $S := (r_n,\theta_n+n)_{n\in \ZZ}$. We see that $S$ is skeletal and that $Y\setminus S$ is semi-locally simply connected. We thus consider $X:= \widetilde Y^{S} \xrightarrow{p} Y$ the maximal covering possibly branched over $S$. 

Notice that $q(Y\setminus \ord_p(Y)) = q(S)$ is dense in $Z$ hence not skeletal, therefore the composition $q\circ p$ is not a branched covering. 
\end{cex}


\subsection{Paths and maximal branched covering}\label{sec:maximal_branched_path}
The abstract description of Fox completion of a spread, though very general,
may benefit from a path point of view.
This section is devoted to a path description 
of the maximal covering possibly branched above some locus akin to 
the path description (as a set of homotopy classes of paths)
of the universal covering of a Hausdorff, connected, locally path-connected,
semi-locally simply connected topological space. 

To begin with, we introduce some notations as well as a refinement of the
homotopy equivalence relation that can keep track of branched locus of a branched covering.
Considering some connected Hausdorff, locally path-connected topological space $X$, 
some connected open subset $\U$ and some $x,y\in \overline\U$ we denote by
$\Omega_X(x,y,\U)$ the set of path $\gamma:[0,1]\rightarrow X$ such that
$\gamma(0)=x$,$\gamma(1)=y$ and $\gamma(]0,1[)\subset \U$. The index $X$ is dropped whenever clear from context. 
This set will always be endowed with its compact-open topology and two paths 
$\gamma_1,\gamma_2$ of $\Omega(x,y,\U)$ are homotopic with respect to $\U$ if they are in the same connected component of $\Omega(x,y,\U)$ or, equivalently, if there exists an homotopy $H:[0,1]^2\rightarrow X$ with fixed ends such that $H(0,\cdot)=\gamma_1$, $H(1,\cdot)=\gamma_2$ and such that $H([0,1]\times ]0,1[)\subset \U$. 
Composition of path is denoted $*$ so that $\gamma_1*\gamma_2$ represents any parameterization of the concatenation of $\gamma_1$ then $\gamma_2$. We also use the notation ${\gamma}^{-1}:t\mapsto \gamma(1-t)$.

The homotopic relation is not suitable for our purpose as this section as well as the following will make clear, we thus introduce a weaker equivalence relation.

\begin{defi}[Almost trivial loop]\label{def:almost_trivial_loop} Let $X$ be a Hausdorff, locally path-connected topological space, let $\U\subset X$ be an open subset and let $x\in \overline \U$.

A loop $\gamma\in \Omega(x,x,\U)$ is  almost trivial with respect to $\U$ if its connected component in $\Omega(x,x,\U)$ intersects $\Omega(x,x,\V\cap \U)$ for all neighborhood $\V$ of $x$. 
\end{defi}
\begin{ex}\label{ex:consider_carthesian_coordinates_usual}
 Consider $Y = \RR^3$ with carthesian coordinates $x,y,z$ and its usual Euclidean distance; let  $S_1= \{y=z=0\}$ and $$S_2=\{y=sin(1/x) ~:~ x\in \RR^*\}\cup \{x=z=0, |y|\leq 1 \};$$ denote $O=(0,0,0)$ the origin. 
 We make the following claims.
 \begin{itemize}
  \item  The fundamental group of $Y\setminus S_i$ is isomorphic to $\ZZ$ for both $i=1$ and $i=2$. The isomorphism is given by "the number of turns" around $S_i$. 
  \item Every loop  $\gamma\in \Omega(O,O,Y\setminus S_1)$ is homotopically trivial hence almost trivial.
  \item Every loop  $\gamma\in \Omega(O,O,Y\setminus S_2)$ is almost trivial. 
  \item There exists non homotopically trivial loops $\gamma\in \Omega(O,O,Y\setminus S_2)$.
 \end{itemize}
\end{ex}
\begin{proof}
The first two claims are left to the reader. Let us prove the last two.
 \begin{itemize}
\item 
 Consider any ball $B=B(O,r)$ centered at $O$, we may choose some $t_1,t_2$ such that $\gamma([0,t_1]),\gamma([t_2,1])  \subset B$. By compactness $\gamma([t_1,t_2])$ there exists  $\varepsilon\in [0,r[$ such that $\gamma_{|[t_1,t_2]}$ stays  at distance at least $2\varepsilon$ of $S_2$.  Since $Y\setminus \overline{B}(S_2,\varepsilon)$ is homeomorphic to $Y\setminus S_1$, we see that $\gamma([t_1,t_2])$ is homotopic in $Y\setminus S_2$ with fixed extreminites to a path in $B\setminus B(S_2,\varepsilon)$. Finally, $\gamma$ is homotopic in $\Omega(O,O,Y\setminus S_2)$ to a path  in $\Omega(O,O,B\setminus S_2)$. 

 \item Consider some $\gamma \in \Omega(O,O,Y\setminus S_2)$ such that
 $$\forall t< 1/4,\quad  \gamma(t) = (0,0,t),~ \gamma(1-t) = (0,0,-t) $$
 and assume by contradiction that $\gamma$ is homotopically trivial.   
  Let $H$ be an homotopy from $\gamma$ to $O$ in $\Omega(O,O,Y\setminus S_2)$ and let $c : [0,1]^2 \rightarrow [0,1]^2$ such that for all $s$, $c(s,\cdot)$ is piecewise affine with 
 \begin{eqnarray*}
c(s,0)=  (0,0)  &\quad&
c(s,1/5)=  (0,1/2-s/2)  \\
c(s,2/5)=  (s,1/2-s/2)  &\quad&
c(s,3/5)=  (s,1/2+s/2)  \\
c(s,4/5)=  (0,1/2+s/2 )  &\quad&
c(s,1)=  (0,1).  
\end{eqnarray*}
We see that $H'=H\circ c$ is still an homotopy from $\gamma$ to $O$ in $\Omega(O,O,Y\setminus S_2)$  with the additional property that 
\begin{eqnarray*} &&\forall s \in [0,1],  \exists! \lambda_s>0, \forall t<1/5, \\ 
&&\quad \quad \quad H'(s,t) = (0,\lambda_s t,0)\quad H'(s,1-t) = (0,-\lambda_s t,0).  \end{eqnarray*}
Denoting $H''(s,t) = H'(s,\frac{3t + 1}{5})$ and  $I(X,Y)$ the mod 2 intersection  numbers of $X$ and $Y$  \cite{MR2680546},  this property implies that 
$$ \forall s \in [0,1], \quad \quad
I(H''(s,\cdot),\{z=0\}) = 1.$$
 
Let $\U=B(O,1/2)$ be the ball of center $0$ and radius $1/2$ and let $\eta>0$ such that for $s>1-\eta$, $H''(s,\cdot)\subset \U$. We notice that $\U\cap \{z=0\}\setminus S_2$ is disconnected, we denote by $\mathbf{C}$ the set of its connected components (which are all closed submanifolds of dimension 2 of $\{z=0\}\setminus S_2$) and define 
$$\fonction{f}{\mathbf{C}}{\ZZ/2\ZZ}{C}{I(H''(s,\cdot),C)}$$
 which is well defined, does not depends on $s>1-\eta$ since the intersection number is a homotopic invariant and has finite support by compactness of $H''(s,[0,1])\cap \{z=0\}$. On the one hand, $$\sum_{C\in \mathbf{C}} f(C) = I(H(s,\cdot),\{z=0\})=1.$$
 On the other hand,  
 $\forall C\in \mathbf{C}, O \notin \overline C$ thus for all $C\in {\mathbf{C}}$, there exists $\omega>0$ such that for all $s>1-\omega$, we have $H(s,\cdot)\cap C = \emptyset$. In other word, $\forall C\in \mathbf{C}, f(C)=0$. Contradition. 


 
 
   
 \end{itemize}
\qed\end{proof}

\begin{defi}[Almost homotopic paths]\label{def:almost_homotopic_paths}
 Let $X$ be a Hausdorff, locally path-connected topological space, let $\U\subset X$ be an open subset and let $(x,y)\in \U\times \overline \U$.

Two paths $\gamma_1,\gamma_2 \in \Omega(x,y,\U)$ are almost homotopic with respect to $\U$ if the loop $(\gamma_2)^{-1}*\gamma_1$ is almost trivial with respect to $\U$.
\end{defi}

\begin{rem}\label{rem:implicitely_two_almost_homotopic} Implicitely, two almost homotopic paths $\gamma_1,\gamma_2$ always satisfy $\gamma_1(1)=\gamma_2(1)$ and $\gamma_1(0)=\gamma_2(0)$. 
\end{rem}

\begin{lem} \label{lem:almost_homotopic_carac}
 Let $X$ be a Hausdorff, first countable, connected, locally path-connected topological space, let $\U\subset X$ be an open subset and let $x \in  \overline \U$.
 
 For $\gamma \in \Omega(x,x,\U)$,  the following are equivalent:
 \begin{enumerate}[(i)]
  \item $\gamma$ is almost trivial with respect to $\U$;
  \item for all open neighborhood $\V$ of $x$ and for all $t_0,t_1\in ]0,1[$ such that $t_0<t_1$ and $\gamma([0,t_0]\cup [t_1,1])\subset \V$, there exists a path $\omega : \gamma(t_0)\rightsquigarrow \gamma(t_1)$ in $\V\cap \U$ such that $\gamma$ is homotopic to $\gamma_{|[0,t_0]}*\omega*\gamma_{|[t_1,t_0]}$ in $\Omega(x,x,\U)$.
 \end{enumerate}
\end{lem}
\begin{proof}
  $(ii)$ trivially implies $(i)$, we thus assume that $\gamma$ is almost trivial with respect to $\U$. Let $\V$ be an open neighborhood of $x$ and let $t_0<t_1$ such that $\gamma([0,t_0]\cup[t_1,1])\subset \V$. Let $H:[0,1]\times [0,1]\rightarrow X$ be an homotopy from $\gamma$ to some path $\gamma_0\in \Omega(x,x
  ,\V)$ such that $\forall s\in [0,1], H(s,\cdot)\in \Omega(x,x,\U)$. 
  Consider the following closed connected subset $B$ of the boundary of $[0,1]\times [0,1]$:$$B := \{0\}\times [0,t_0]\ \cup\ [0,1]\times \{0\}\ \cup\ \{1\}\times [0,1]\ \cup\ [0,1]\times\{1\}\ \cup\ \{0\}\times [t_1,1].$$
  Since $H(B)\subset \V$, there exists a connected open neighborhood $\W$ of $B$ in $[0,1]\times [0,1]$ such that $H(\W)\subset \V$. Take some path $\eta:[t_0,t_1]\rightarrow \W$ such that $\eta(t_0)=(0,t_0)$, $\eta(t_1)=(0,t_1)$ and $\eta(]t_0,t_1[)\subset ]0,1[\times ]0,1[$ and define $\omega:=H\circ \eta$. 
  With $\eta(\cdot)=(\eta_1(\cdot),\eta_2(\cdot))$, the homotopy $$\fonctiontrois{J}{[0,1]\times [0,1]}{X}{s\in [0,1],t\in [0,t_0]}{H(0,t)}{s\in [0,1], t\in [t_0,t_1]}{H\left(s\cdot\eta_1(t),(1-s)t+s\cdot\eta_2(t)\right)}{s\in [0,1],t\in [t_1,1]}{H(0,t)}$$
  is such that $J(0,\cdot)=\gamma$ and $J(1,\cdot)=\gamma_{|[0,t_0]}*\omega*\gamma_{|[t_1,t_0]}$, moreover  for all $ s\in[0,1]$, we have $J(s,\cdot)\in \Omega(x,x,\U)$. The result follows.
\qed\end{proof}


\begin{lem}\label{lem:hausdorff_first_countable_connected} 
 Let $X$ be a Hausdorff, first countable, connected, locally path-connected topological space, let $\U\subset X$ be an open subset and let $(x,y)\in \U\times \overline \U$.
 
 Then, the "almost homotopic" relation is an equivalence relation in $\Omega(x,y,\U)$.
 
\end{lem}
\begin{proof}
The relation is clearly reflexive and symmetric, we shall then prove that it is also transitive. Let $\gamma_1, \gamma_2,\gamma_3 \in \Omega(x,y,\U)$ such that $\gamma_1,\gamma_2$ are almost homotopic with respect to $\U$ and such that
$\gamma_2,\gamma_3$ are almost homotopic with respect to $\U$. Let $\V$ be an open neighborhood of $x$ and let $t \in ]0,1[^3$ such that $\forall i\in\{1,2,3\},\ \gamma_i([t_i,1])\subset \V$. By Lemma \ref{lem:almost_homotopic_carac}, there exists a path $\omega_{1}$ (resp. $\omega_2$)  in $\V\cap \U$ from $\gamma_1(t_1)$ to $\gamma_2(t_2)$ (resp. from $\gamma_2(t_2)$ to $\gamma_3(t_3)$) such that $\gamma_1^{-1}*\gamma_2$ is homotopic to $(\gamma_{1|[t_1,1]})^{-1}
*\omega_1*\gamma_{2|[t_2,1]}$ (resp. such that $\gamma_2^{-1}*\gamma_1$ is homotopic to $(\gamma_{2|[t_2,1]})^{-1}*\omega_2*\gamma_{3|[t_3,1]}$ ). Then $(\gamma_{1|[0,t_1]})^{-1}*\gamma_{2|[0,t_2]}*\omega_1^{-1}$ is 
homotopically trivial and the loop
 $\gamma_{1}^{-1}*\gamma_3$ is homotopic to 
 $$(\gamma_{1|[t_1,1]})^{-1}*\left((\gamma_{1|[0,t_1]})^{-1}*\gamma_{2|[0,t_2]}*\omega_1^{-1}\right)*\omega_1*(\gamma_{2|[0,t_2]})^{-1}*\gamma_3$$ which is thus homotopic to $(\gamma_{1|[t_1,1]})^{-1}*\omega_1*(\gamma_{2|[0,t_2[})^{-1}*\gamma_3$. The same way, the loop $(\gamma_{2|[0,t_1]})^{-1}*\gamma_{3|[0,t_3]}*\omega_2^{-1}$ is homotopically trivial and $\gamma_{1}^{-1}*\gamma_3$ is then homotopic to $(\gamma_{1|[t_0,1]})^{-1}*\omega_1*\omega_3*\gamma_{3|[t_3,1]} \in \Omega(x,x,\V\cap \U)$. 
 Finally, $\gamma_{1}^{-1}*\gamma_3$ is almost trivial with respect to $\U$ and $\gamma_1$ is almost homotopic to $\gamma_3$.
\qed\end{proof}

\begin{lem} \label{lem:almost_homotopic_carac2}
 Let $\widetilde Y^S\xrightarrow{p} Y$ be a maximal covering possibly branched above some $S$, let $(x,y)\in (Y\setminus S)\times Y$ and let $\hat x \in p^{-1}(x)$.
Let $\gamma_1, \gamma_2\in \Omega(x,y,Y\setminus S)$ and let $\hat \gamma_1,\hat\gamma_2$ be the respective lifts of $\gamma_1$ and $\gamma_2$ such that $\hat \gamma_1(0)=\hat \gamma_2(0)=\hat x$. 

Then, $\hat\gamma_1(1)=\hat\gamma_2(1)$ if and only if $\gamma_1$ and $\gamma_2$ are almost homotopic.
\end{lem}
\begin{proof}
 Assume $\gamma_1$ and $\gamma_2$ are almost homotopic. For  $\U$ some standard neighborhood of $\hat\gamma_1(1)$. By Corollary \ref{cor:spread_connected_image}, $p(\U)$ is a connected open neighborhood of $y$, and since $\gamma_1$ and $\gamma_2$ are almost homotopic, $\gamma_1$ is homotopic to a path $\gamma_1'$ of the form $\gamma_{2|[0,t_2]}*\omega*\gamma_{1|[t_1,1]}$ with $t_1\in ]0,1[$
 such that $\gamma_i([t_i,1])\subset p(\U)$ for $i\in \{1,2\}$ and $\omega\in  \Omega(\gamma_2(t_2),\gamma_1(t_1),Y\setminus S)$. 
 Since $\gamma_1$ and $\gamma_1'$ are homotopic and $p^{-1}(y)$ is totally disconnected, then $\gamma_1(1)=\gamma_1'(1)$. Since $\omega*\gamma_{1|[t_1,1]}\subset p(\U)$, the lift $\eta$ of $\omega*\gamma_{1|[t_1,1]}$ such that $\eta(0)=\gamma_2(t_2)$ stays in $\overline \U$. In particular $\hat\gamma_1(1)=\eta(1)\in \overline \U$. Therefore, $\hat\gamma_1(1)$ is in the closure of all neighborhood of $\hat\gamma_2(1)$, since $\widetilde Y^S$ is Hausdorff we conclude that $\hat\gamma_1(1)=\hat\gamma_2(1)$.
 
 Assume $\hat\gamma_1(1)=\hat\gamma_2(1)$. Let $\U\subset \widetilde Y^S$ be some standard neighborhood of  $\hat\gamma_1(1)$, let $t_1,t_2\in ]0,1[$ such that $\hat\gamma_i([t_i,1])\subset \U$ for $i\in \{1,2\}$ and let $\omega \in \Omega(\hat\gamma_1(t_1),\hat\gamma_2(t_2),\U\cap p^{-1}(Y\setminus S))$. Since $p^{-1}(Y\setminus S)$ is simply connected, the loop $\hat\gamma_{1|[0,t_1]} *\omega * (\hat\gamma_{2|[0,t_2]})^{-1}$ is homotopically trivial. Therefore, $\hat\gamma_2$ is homotopic to $\hat\gamma_{1|[0,t_1]}*\omega*\hat\gamma_{2|[t_2,1]}$, hence $\gamma_2$ is homotopic to $\gamma_{1|[0,t_1]}*(p\circ\omega)*\gamma_{2|[t_2,1]}$. Finally, the loop $(\gamma_1)^{-1}*\gamma_2$ is homotopic to  $(\gamma_1)^{-1}*\gamma_{1|[0,t_1]}*(p\circ\omega)*\gamma_{2|[t_2,1]}$ which is homotopic to $(\gamma_{1|[t_1,1]})^{-1}*(p\circ \omega)*\gamma_{2|[t_2,1]} \subset \U$. 
The loop $(\gamma_1)^{-1}*\gamma_2$ is thus almost trivial and $\gamma_1$ is almost homotopic to $\gamma_2$. 
\qed\end{proof}


\begin{thm} \label{theo:ramified_covering_path} Let $Y$ be a first countable connected locally path-connected Hausdorff topological space and let $S\subset Y$ be a skeletal subset. 
Let $y_0\in Y\setminus S$ and let $\Omega:=\bigcup_{x\in Y}\Omega(y_0,x,Y\setminus S)$. Define $X$ as the set of almost homotopy classes of $\Omega$ with respect to $Y\setminus S$ and define the map $\varphi:\Omega\rightarrow Y, \gamma \mapsto \gamma(1)$.


Then, $\varphi$ induces a branched covering $X \xrightarrow{~~\overline\varphi~~ } Y$  isomorphic to $\widetilde Y^S \xrightarrow{~~p~~} Y$.
 
\end{thm}
\begin{proof}

Let $\hat y_0$ be some lift of $y_0$ in $\widetilde Y^S$.
Consider the function $\psi : \Omega\rightarrow \widetilde Y^S$ which associate to $\gamma\in \Omega$ the point $\hat\gamma(1)$ where $\hat\gamma$ is the unique lift of $\gamma$ in $\widetilde Y^S$ such that $\hat\gamma(0)=\hat y_0$.
By Lemma \ref{lem:almost_homotopic_carac2}, $\psi$ induces an injective function $\overline \psi:X\rightarrow \widetilde Y^S$,
the following diagram commutes
$$\xymatrix{X\ar[d]^{\overline \varphi}\ar[r]^{\overline\psi}& \widetilde Y^S \ar[d]^p \\ Y\ar@{=}[r]&Y}$$
and it suffices to prove that $\overline \psi$ is an homeomorphism.

\begin{itemize}
 \item Since $\widetilde Y^S$ is first countable and $p^{-1}(Y\setminus S)$ is portly, for all $x\in \widetilde Y^S$ there exists a path $\hat\gamma\in \Omega(\hat y_0,x,p^{-1}(Y\setminus S))$ and, for such a path,  $\psi(p\circ \hat\gamma)=\hat\gamma(1)=x$. The map $\psi$ is thus surjective, hence bijective.
 
 \item Let $\U$ be a neighborhood of some $x\in \widetilde Y^S$, since $p$ is a spread and by Corollary \ref{cor:spread_connected_image} without loss of generality, we may assume $\U$ is a connected component of $p^{-1}(p(\U))$
 and consider some $\gamma\in \Omega$ such that $\psi(\gamma)=x$.  
 Recall that $p(\U)$ is also open.
 
 Take some $t_0\in ]0,1[$ such that $\gamma([t_0,1])\subset p(\U)$.
 Since $Y\setminus S$ is semi-locally simply connected, for all $t\in [0,t_0]$ there exists some path-connected open neighborhood $\V\subset Y\setminus S$ of $\gamma(t)$ such that all loop in $\V$ are homotopically trivial in $Y\setminus S$. By compactness, we can choose finitely many such neighborhoods $(\V_i)_{i\in\lsem 1,n \rsem}$ as well as an increasing sequence $\alpha\in [0,t_0]^{\lsem 0,n\rsem}$ such that :
 $\alpha_0=0, \alpha_n=t_0$ and $\forall i\in \lsem 1,n\rsem,\ \gamma([\alpha_{i-1},\alpha_i]) \subset \V_i$. 
 Define $\V_{n+1}:=\U$ and $\W_i$ the connected component of $\gamma(\alpha_i)$ in $\V_i\cap\V_{i+1}$ for $i\in \lsem 1,n\rsem$ and define 
 $$\mathcal O := \left\{\eta\in \Omega ~\left|~\begin{array}{l}
                                                     \forall i\in \lsem 1,n\rsem, \eta([\alpha_{i-1},\alpha_i])\subset \V_i)  \\ \forall i\in \lsem 1,n\rsem,\ \eta(\alpha_i)\in\W_i \\ 
                                                     \eta([t_0,1])\subset p(\U)
                                                      \end{array} \right.\right\}.$$
$\mathcal O$ is an open neighborhood of $\gamma$ for the compact-open topology of $\Omega$. Let $\eta\in \mathcal O$, forall $i\in \lsem 1,n\rsem $, choose a path  $\omega_i \in 
\Omega(\gamma(\alpha_i),\eta(\alpha_i),\W_i)$. We also define $\omega_0$ as the 
constant path at $y_0$. We notice that for all $i\in \lsem 1,n,\rsem$, the loop $
\gamma_{|[\alpha_{i-1},\alpha_i]}*\omega_i*(\eta_{|[\alpha_{i-1},\alpha_i]})^{-1}*\omega_{i-1} $ is homotopically trivial. Therefore, $\eta$ is homotopic to the loop 
$\eta':=\gamma_{|[0,t_0]}*\omega_n*\eta_{|[t_0,1]}$. Consider the lifts $\hat\gamma$
and $\hat\eta'$ of $\gamma$ and $\eta'$ respectively, both starting at $\hat y_0$. They are equal on $[0,t_0]$; since $\gamma$ and $\eta$ both stay in $p(\U)$, the paths $\hat\gamma$ and $\hat\eta'$ stay in the same connected component of $p^{-1}(p(\U))$, hence they both stay in $\U$. We then deduce that $\psi(\eta)=\hat\eta'(1)\in \U$ thus $\psi(\mathcal O)\subset \U$.
Finally, $\psi$ is continuous, then so is $\overline \psi$.

\item Let $\U\subset Y$ open and $K\subset [0,1]$ compact, define 
$\mathcal O := \{ \gamma\in \Omega~|~ \gamma(K)\subset \U \}$. Four cases can occur:
\begin{enumerate}
 \item If $0\in K$ and $y_0 \notin \U $, then $\psi(\mathcal O)= \emptyset$;
 \item If $0\notin K$ or $y_0 \in \U$;
    \begin{enumerate}
     \item If $1 \notin K$, then $\psi(\mathcal O)=\widetilde Y^S$;
     
     
     \item If $1\in K\neq [0,1]$ , then $\psi(\mathcal O)= p^{-1}(\U)$;
     \item If $K = [0,1]$ then $\psi(\mathcal O)$ is the connected component of $\hat y_0$ in $p^{-1}(\U)$.
    \end{enumerate}
\end{enumerate}
Case $(1)$ is trivial. In case $(2.a)$, consider $x\in \widetilde Y^S$, $\gamma\in \Omega$ such that $\psi(\gamma)=x$. Take some $t_0\in ]0,1[$ such that $[t_0,1]\cap K = \emptyset$ and consider the path $\eta(t)=y_0$ if $t\leq t_0$ and 
$\eta(t)=\gamma(\frac{t-t_0}{1-t_0})$ if $t\in [t_0,1]$. We indeed have $\psi(\eta)=\psi(\gamma)=x$. A similar argument shows case $(2.b)$. 
Then, in case $(2.c)$, we have $y_0\in \U$ and the connected component $\hat\U$ of $\hat y_0$ in $p^{-1}(\U)$  is path-connected. For $x\in \hat\U$, there thus exists some path $\hat\gamma \in \Omega(\hat y_0,x,\hat \U)$ and, for such a path, $\psi(p\circ \gamma)=x$. Furthermore $p\circ \gamma\subset \U$ thus $p\circ \gamma\in \mathcal O$. Therefore, $\hat \U \subset \psi(\mathcal O)$. Certainly, $\psi(\mathcal O)\subset p^{-1}(\U)$ and since $\hat\gamma$ is connected and $\hat\gamma(0)=\hat y_0\in \hat\U$, then $\psi(\mathcal O)\subset \hat \U$.
Finally, $\psi(\mathcal O)=\hat \U$, the connected component of $\hat y_0$ in $p^{-1}(\U)$.

We deduce from the previous analysis that $\psi$ is open, hence $\overline \psi$ is open.
\end{itemize}
Together, the points above show that $\overline \psi$ is an homeomorphism.
\qed\end{proof}

We end this section with two corollaries of Proposition \ref{prop:spread_fiber}, the proofs of which are left to the reader, and two examples illustrating possible behaviors of the fiber above a branching point. 
\begin{defi}\label{def:covering_branched_two_paths} 
 Let $X\xrightarrow{p} Y$ be a  covering branched over $S\subset Y$. 
 Given two paths $\gamma_1,\gamma_2$ with fixed start point $a\in Y\setminus S$ and free end point in an open connected subset $\U\subset Y\setminus S$; $\gamma_1$ and $\gamma_2$ are $p$-homotopic if their lifts, from the same start point, end in the same connected component of $p^{-1}(\U)$.
\end{defi}

  \begin{prop} \label{prop:general_fiber_path}
  Let $X\xrightarrow{p} Y$ be a  covering branched over $S\subset Y$. Let $a\in Y\setminus S$ and let $\U\subset Y$ be open and connected.  Define
$\pi_1(a,p,\U)$ the set of paths with fixed start at $a$ and free end in $\U$ up to $p$-homotopy.

Then, for all $b\in Y$ and all connected neighborhood basis $\boldsymbol \U$ of $b$ we have 
  $$ p^{-1}(b) \simeq \varprojlim_{\U\in \boldsymbol\U} \pi_1(a,p,\U)$$
  where the projective system is given by the natural maps $\pi_1(a,p,\U)\rightarrow \pi_1(a,p,\V)$  for $\U\subset \V$.
\end{prop}

\begin{prop} \label{prop:general_fiber_path_universal}
 Let $Y$ be a first countable connected locally path-connected Hausdorff topological space and let $S\subset Y$ skeletal such  that $Y\setminus S$ is semi-locally simply connected. Let $a\in Y\setminus S$ and let $\U\subset Y$ be open and connected.   Denote $\widetilde Y^S\xrightarrow{p} Y$ the maximal cover of $Y$ possibly branched over~$S$. Define
$\pi_1(a,\U)$ the set of paths with fixed start at $a$ and free end in $\U$ up to homotopy.

Then, for all $b\in Y$ and all connected neighborhood basis $\boldsymbol \U$ of $b$ we have 
  $$ p^{-1}(b) \simeq \varprojlim_{\U\in \boldsymbol\U} \pi_1(a,\U)$$
  where the projective system is given by the natural maps $\pi_1(a,\U)\rightarrow \pi_1(a,\V)$  for $\U\subset \V$.
\end{prop}
  

\begin{ex}\label{ex:consider_fiber_above_homeomorphic} Consider $Y = \RR^2$ and $S=\{(\frac{1}{n},0)~:~n\in \N^*\cup\{\infty\}\}$. The fiber above  $(0,0)$ in $\widetilde Y^S$ is homeomorphic to $\mathbb Z^\N$ endowed with the weak topology; hence is not locally compact. This answers positively Problem 10.8 of \cite{MR2198590}.
\end{ex}
\begin{proof}
 Consider the neighborhood basis given by the discs $D_n$ centered at $(0,0)$ of radius $\frac{1}{n-1/2}$ for $n\in \N^*$. Let $a = (2,0)$ and  $b:=(0,0)$, for $n\in \N^*$, the natural map $\pi_1(a,D_{n+1})\rightarrow\pi_1(a,D_{n})$ has an infinite fiber above every point. The result then follows from Proposition \ref{prop:general_fiber_path_universal}.
\qed\end{proof}



\subsection{Galoisian branched covering}\label{sec:galoisian_branched}

The usual notion of Galoisian covering can be extended in a natural way to branched covering via the group of automorphisms of the branched covering. Although, as we shall see, the completion of a Galoisian covering need not be Galoisian. The main result of this section is a topological criteria akin to semi-local simple connectedness which ensures maximal branched coverings are Galoisian, see Proposition \ref{prop:Galoisian_semilocal}. It allows to state a Galoisian correspondance for branched covering stated in Theorem \ref{theo:Galois_cor}.
\begin{defi}[Automorphisms of a spread]\label{def:automorphisms_spread} Let $X\xrightarrow{p} Y$ be a spread, an automorphisms of $p$ is an homeomorphism $\varphi$ of $X$ such that $p\circ \varphi = p$. 
Denote by  $\Gamma(X/Y)$ the group of spread automorphisms of $X$ above $Y$. 
We also denote by $\Gamma_x$ the stabilizer of $x$ if $x\in X$ and by $\Gamma_\U$ the set-wise stabilizer of $\U$ if $\U$  is a subset of $X$.
\end{defi}
Note that a spread automorphism is just an automorphism in the category $\CatSpread$.

\begin{defi}\label{def:branched_covering_galoisian_group} A branched covering $X\xrightarrow{p} Y$ is Galoisian if the group $\Gamma(X/Y)$ of automorphism of $p$ acts transitively on the fibers of $p$. It is called quasi-Galoisian if $p$ is the completion of a Galoisian covering.  


\end{defi}

\begin{rem}\label{rem:galois_quotient} If a branched covering $X\xrightarrow{p} Y$ is Galoisian then $p$ induces an homeomorphism $\Gamma(X/Y)\backslash X \rightarrow Y$. 
\end{rem}


\begin{lem}\label{lem:branched_covering_group_automorphisms} Let $X\xrightarrow{p} Y$ be a branched covering, the group of automorphisms of $p$ is exactly the group of automorphisms of the induced branched covering $p^{-1}(\U) \rightarrow \U$ for any $\U$ portly in $Y$:
$$\Gamma(X/Y) = \Gamma(p^{-1}(\U)/\U).$$
\end{lem}
\begin{proof}
 An automorphism $\phi$ of $X \rightarrow Y$ lifts the identity thus  for any $\U$ subset  of $Y$, the automorphism $\phi$ preserves $p^{-1}(\U)$. In particular, for any $\U$ portly subset of $Y$, such a $\phi$ induces an automorphism of $p^{-1}(\U) \xrightarrow{~~~} \U$. Furthermore, with $\U$ portly in $Y$, by functoriality of the spread completion and since the completion of $p^{-1}(\U)\rightarrow Y$ is $X\rightarrow Y$, an automorphism of $p^{-1}(\U) \xrightarrow{~~~} \U$ extends uniquely to a automorphism of $X \xrightarrow{~~~} Y$.   Therefore, $\Gamma(X/Y) = \Gamma(p^{-1}(\U)/\U))$.
\qed\end{proof}
\begin{cor}\label{cor:branched_covering}Let $X\xrightarrow{p} Y$ be a branched covering, then
 $$\Gamma(X/Y) =  \Gamma(\ord(X)/\ord(Y)).$$
\end{cor}



The maximal covering branched over some locus need not be Galoisian. Indeed, one can consider $Y=\RR^2$ and $S:=\{(1/n,0):n\in \N^*\}\cup \{(0,0)\}$. On the one hand, the automorphisms group of $\widetilde Y^S$ is countable (it is the fundamental group of $Y\setminus S$ thus a free group generated by countably infinitely many generators). On the other hand, there are uncountably infinitely many homotopy classes in of path from $(2,0)$ to $(0,0)$ and none of these classes are equivalent in the sense of Theorem \ref{theo:ramified_covering_path};
hence the fibre above $(0,0)$ is uncountable. The action of $\Gamma(X/Y)$ is then not transitive on the fiber above $0$ and $\widetilde Y^S \rightarrow Y$ is not Galoisian.


\begin{defi}[semi-locally simple connectedness]\label{def:semi_locally_simple_connectedness} Let $X$ be a Hausdorff, connected, locally path-connected topological space.  We say that an open subset $\U$ is semi-locally simply connected at $x \in \overline \U$ if there exists a connected open neighborhood $\V$ of $x$ such that every loop $\gamma$ of $\Omega(x,x,\V\cap \U)$ seen as an element of $\Omega(x,x,\U)$ is almost trivial . 

$\U$ is then semi-locally simply connected in $X$ if $\U$ is dense and semi-locally simply connected at every point of $X$
\end{defi}
\begin{rem}\label{rem:definition_semi_locally_simply_connected}
 The definition of semi-locally simply connected is coherent with the usual definition of semi-locally simply connected. Indeed a topological space $X$ is semi-locally simply connected if and only if it is semi-locally simply connected in  itself.
\end{rem}

\begin{rem} \label{rem:simple_cases} Using the same notations as in the definition above, there are generic situations in which dense open subsets are semi-locally simply connected at every point of the whole space.
\begin{itemize}
 \item If $X$ is a manifold and $X\setminus \U$ is a codimension at least 2 submanifold, more generally if $X$ is a simplicial complex and $X\setminus \U$ is a union of facets (such that $\U$ is still dense and locally connected);
 \item  or if $X$ admits a basis of neighborhoods such that for each $\V$ of said basis, $\V\cap \U$ is semi-locally simply connected in $\V$;
\end{itemize}
then $\U$ is locally simply connected at every point of $X$. 
\end{rem}



\begin{lem}\label{lem:dense_orbits}
 Let $X\xrightarrow{p} Y$ be a quasi-galoisian branched covering. Let $a\in Y$ and $\hat a \in p^{-1}(a)$, the orbit  $\Gamma(X/Y)\hat a$  is dense in $p^{-1}(a)$.
\end{lem}
\begin{proof} We use the same notation as in Proposition \ref{prop:spread_fiber}. 
  Since the branched covering $p$ is quasi-Galoisian, for every $\U\in \boldsymbol\U$, $\Gamma(X/Y)$ acts transitively on $X_\U$. Furthermore, $\Gamma(X/Y)$ acts by homeomorphisms on $X$, hence preserves connected components, its action thus commutes with the bonding maps of the projective system $(X_\U)_{\U\in \boldsymbol{\U}}$. Finally, the action of $\Gamma(X/Y)$ on each $X_\U$ lifts to an action on $\varprojlim_\U X_\U$ whose orbits are dense.
\qed\end{proof}



\begin{prop}\label{prop:Galoisian_fiber}
  Let $X\xrightarrow{p} Y$ be a quasi-Galoisian branched covering possibly branched over $S\subset Y$.
  
  For all $b\in Y$ and all connected neighborhood basis $(\U_i)_{i\in I}$ of $b$ we have 
  $$ p^{-1}(b) \simeq \varprojlim_{i\in I} \Gamma(X/Y)/\Gamma_{\hat\U_i} $$
  where $\hat\U_i$ is the connected component of $\hat b$ in $p^{-1}(\U_i)$ for $i\in I$.
  
\end{prop}

\begin{proof}
  Since $X\xrightarrow{p} Y$ is quasi-Galoisian, $\Gamma(X/Y)$ acts transitively on the connected components of $p^{-1}(\U)$ for every open $\U\subset Y$. The group $\Gamma_\V$ is by definition the stabilizer of $\V$ for $\V$ a connected of $p^{-1}(\U)$  for some open $\U\subset Y$. The discrete spaces $\Gamma(X/Y)/\Gamma_{\hat\U}$ and $X_\U$ are thus homeomorphic (using the same notations as Proposition \ref{prop:spread_fiber}). The result then follows from Proposition \ref{prop:spread_fiber}.

  \qed\end{proof}

\begin{lem} \label{lem:alternative_countable}
 Let $X\xrightarrow{p} Y$ be a quasi-Galoisian branched covering possibly branched over $S\subset Y$ and let $b\in Y$. 
 
 Assume that $\Gamma(X/Y)$ is countable, then the following are equivalent:
 \begin{enumerate}[(i)]
  \item $p^{-1}(b)$ is discrete;
  \item $p^{-1}(b)$ is countable;
  \item $\Gamma(X/Y)$ acts transitively on $p^{-1}(b)$.
 \end{enumerate}

\end{lem}
\begin{proof}
 Consider a neighborhood basis $(\U_n)_{n\in \N}$ of $b$ in $Y$ indexed over $\N$ such that $\U_n\subset \U_m$ whenever $n\geq m$. Denote by $X_n$ the discrete space of connected components of $p^{-1}(\U_n)$ and $(p_{n,m})_{n\geq m}: X_n \rightarrow X_m$ the bonding maps which associated a connected component $\V$ of $p^{-1}(\U_n)$
 its connected component in $p^{-1}(\U_m)$.
 For $n\in\N$, define $N_n := \#p_{n+1,n}^{-1}(x)$ for some $x\in X_n$. Since the action of $\Gamma(X/Y)$ is transitively on each $X_n$ and commutes with the bonding maps, $N_n$ does not depends on the choice of $x$ and $N_n\leq \aleph_0$.
 \begin{itemize}
  \item  If $N_n>1$ for infinitely many $n\in \N$, then the projective limit $\varprojlim_n X_n$ is uncountable, hence the action of $\Gamma(X/Y)$ is not transitive. Furthermore, a basis of the topology is given by the inverse image of   subsets of the $X_n$ by the natural projection $p_m : \varprojlim_n X_n \rightarrow X_m$. None of these preimage is reduced to a point, hence the topology of $\varprojlim_n X_n$ is not discrete.
  \item  If $N_n> 1$ only for finitely many $n\in \N$, then the projective limit  $\varprojlim_n X_n$ is  stationnary. Then $\varprojlim_n X_n \simeq X_m$ for some $m\in \N$ big enough. In particular $\varprojlim_n X_n$ is discrete and, by Lemma \ref{lem:dense_orbits}, the action of $\Gamma(X/Y)$ on it is transitive, since $\Gamma(X/Y)$ is countable then so is $\varprojlim_n X_n$. 
 \end{itemize}
\qed\end{proof}


\begin{prop} \label{prop:Galoisian_semilocal}
Let $Y$ be a connected, locally path-connected, second countable, Hausdorff topological space and let $S\subset Y$ skeletal such that $Y\setminus S$ is semi-locally simply connected.
Let $\widetilde Y^S\xrightarrow{p} Y$ be the maximal covering of $Y$ possibly branched  above $S$. The following are equivalent:
\begin{enumerate}[(i)]
\item $\widetilde Y^S\xrightarrow{p} Y$  is Galoisian;
  \item $Y\setminus S$ is semi-locally simply connected in $Y$;
\end{enumerate}
\end{prop}
\begin{proof}\
We write $\Gamma = \Gamma(\widetilde Y^S/Y)$ and $X = \widetilde Y^S$. By Lemma \ref{lem:alternative_countable}, we can replace $(i)$ by the equivalent statement that the fibers of $p$ are discrete. 

  \begin{itemize}

   \item $(ii)\Rightarrow (i)$. Assume that $\ord(Y)$ is semi-locally simply connected at every point of $Y$ and consider some $y\in Y$, we show that $p^{-1}(y)$ is discrete.
   
   If $y\in \ord(Y)$, it follows from standard theory of (unbranched) covering; henceforth we assume $y\notin \ord(Y)$ and consider some connected open neighborhood $\U$ of $y$ such that for every loops $\gamma\in\Omega(y,\ord(\U))$,  the constant loop at $y$ is in the adherence of the connected component of $\gamma$ in  $\Omega(y,\ord(Y))$. We choose a connected component $\V$ of $p^{-1}(\U)$ and consider $x_1,x_2 \in p^{-1}(y)\cap \V$.  Consider a path $\hat\gamma\in \Omega(x_1,x_2,\ord(X)\cap \V)$. Such a path exists since $\V$ is connected and $\ord(X)$ portly hence $\ord(X)\cap \V$ is open dense in $\V$ and path-connected. 
   
   For any open neighborhood $\W$ of $y$, by hypothesis we can choose some $\gamma_1\in \Omega(y,\W)$ in the connected component of $\gamma_0=p\circ \hat\gamma$ in $\Omega(y,\ord(Y)$.  Let $H:[0,1]\times[0,1]\rightarrow Y$ be an homotopy from $\gamma_0$ to $\gamma_1$ in $\Omega(y,\ord(Y))$ the ordinary part of $H$ defined on 
   $]0,1[\times [0,1]\rightarrow Y$ lifts to $\ord(X)$ and, by Lemma \ref{lem:complete_path_carac}, this lift extends continuously to a lift $\widetilde H:[0,1]\times [0,1]\rightarrow X$ of $H$. Since the fibers of $p$ are totally disconnected and $\forall s\in [0,1], H(0,s)=H(1,s)=y$, then $\widetilde H(0,\cdot)$ and $\widetilde H(1,\cdot)$ are constant. In particular, $x_1$ and $x_2$ are in the same connected component in $p^{-1}(\W)$. We deduce that $x_1=x_2$ by Proposition \ref{prop:spread_fiber}. Finally, $\V\cap p^{-1}(y)$  is a singleton.
   
 
 
 \item $(i)\Rightarrow (ii)$. Assume the fibers of $p$ are discrete and consider some neighborhood $\V$ of some $x\in X$ such that $\V\cap p^{-1}(p(x))=\{x\}$; denote by $\U:=p(\V)$. Since $p$ is a branched covering, by Corollary \ref{cor:spread_connected_image} we may choose $\V$ connected and such that $\V$ is a connected component of $p^{-1}(\U)$. 
 
 Consider some path $\gamma \in \Omega(y,\U)$ and some lift $\hat\gamma : [0,1]\rightarrow \V$. Since $\gamma(0)=\gamma(1)=y$ both $\hat\gamma(0)$ and $\hat \gamma(1)$ are in $p^{-1}(y)\cap \V$ and are thus equal to $x$ ie $\hat\gamma \in \Omega(x,\V)$. 
 
 Let $\W$ be some neighborhood of $y$, choose some $t_0,t_1 \in ]0,1[$ such that $\gamma(]0,t_0])\subset \W$ and $\gamma([t_1,1[)\subset \W$. The connected component $\hat\W$ of $x$ in $p^{-1}(\W)$ contains both $\hat\gamma([t_1,1[)$ and $\hat\gamma(]0,t_0])$; moreover $\hat\W\cap \ord(X)$ is path-connected; there thus exists a path $\omega:\hat\gamma(t_1)\rightsquigarrow \hat\gamma(t_0)$ in $\hat\W$. Since $\ord(X)$ is the universal covering of $\ord(Y)$, it is simply connected and the loop $\gamma_{|[t_0,t_1]}* \omega$ is trivial. Therefore, $\hat\gamma$ is homotopic to $\hat\gamma_{[t_1,1]}* \omega^{-1} * \hat\gamma_{[0,t_0]}$.  Hence, $\hat\gamma$ is homotopic in $\Omega(x,\ord(X))$ to a loop in $\Omega(x,\ord(\hat\W)$. Finally, $\gamma$ is homotopic in $\Omega(y,\ord(Y))$ to a loop in $\Omega(x,\ord(\W))$. The open subset $\W$ is arbitrary, then $\ord(Y)$ is semi-locally simply connected at $y$.
 
  \end{itemize}
\end{proof}





The following Theorem is then a direct consequence of Propositions \ref{prop:Galoisian_semilocal} and \ref{prop:universality}.
\begin{thm}[Galois correspondance]\label{theo:Galois_cor}
 Let $Y$ be a connected, locally path-connected Hausdorff topological space and let $S\subset Y$ skeletal.
 Assume $Y\setminus S$ is semi-locally simply connected in $Y$. 
 
 Then the maximal covering $\widetilde Y^S$ possibly branched over $S$ is Galoisian. 

 Furthermore, every covering $X$ of $Y$ possibly branched over $S$ is isomorphic to the quotient of $\widetilde{Y}^S$ by a subgroup of $\Gamma(\widetilde{Y}^S/Y)$. 
 Finally, this subgroup is normal if and only if $X$ is Galoisian.
\end{thm}

\begin{ex}\label{ex:following_situations_semi_locally_simply}  In the following situations,  $Y\setminus S$ is semi-locally simply connected in $Y$ 
\begin{enumerate}
 \item $Y$ is a $n$-manifold with $n\geq 2$ and $S$ is a finite subset;
 \item $Y$ is a 3-manifold and $S$ is a tame knot.
 \item $Y$ is a simplicial complex of dimension $n$ and $S$ is a $(n-2)$-skeleton.
\end{enumerate}

 
\end{ex}

\begin{ex}\label{ex:consider_denote_absolute_galois} Consider $Y = \RR^2$ and $S=\{(\frac{1}{n},0)~:~n\in \N^*\cup\{\infty\}\}$ and denote by $\Gamma$ the absolute Galois group $\Gamma(\widetilde Y^S/Y)$. We see that $Y\setminus S$ is not semi-locally simply connected in $Y$.  
Let $\Gamma_2:=\langle \Gamma[\gamma_n]^2\Gamma~:~ n\in \N \rangle$ where $\gamma_n$ is a simple loop around $(1/n,0)$. Consider $Y_2$, the completion of the spread $\Gamma_2\backslash\ord(\widetilde Y^S) \rightarrow Y$.
The fiber above $(0,0)$ in  $Y_2$ is  a Cantor set. This answers positively Problem 10.7 of \cite{MR2198590}. 


\end{ex}


\begin{ex}\label{ex:wild_cantor_general_antoine} If $S$ is a  wild Cantor \cite{general_antoine} embedded into $Y = \SS^3$, then $Y\setminus S$ is not semi-locally simply connected in $Y$.
\end{ex}



% 
% \subsection{On covering branched over a wild cantor}
% 
% Montesinos in \cite{MR2198590} asked wether one could construct a branched covering of the sphere by itself branched over a cantor set. In order to have such a branching locus, the complement has to be non simply connected hence the considered Cantor set is wildly embedded in $\SS^3$.  
% 
% By a Theorem of ??? \cite{finite_ram_cov}, such a covering is of finite degree.
% 
% \begin{lem}
%  Let $n\in \N$ and let $M\xrightarrow{p} N$ be a branched covering of $n$-manifolds $M,N$. Let $S\subset M$ be its branching locus.
% 
%  If $p$ is finite and galoisian, then forall $\hat x\in M$, $\Gamma_x$ is non trivial if and only if $\hat x \in S$
% \end{lem}
% \begin{proof}
%  
% \end{proof}
% 
% 
% \begin{lem}Let $n\in \N$ and let $M\xrightarrow{p} N$ be a branched covering of $n$-manifolds $M,N$. Let $S\subset M$ be its branching locus.
% 
% If $p$ is finite and galoisian then there exists $\hat x\in M$ and a standard neighborhood $\U$ of $x$ such that $\forall y\in \U\cap S, \Gamma_x=\Gamma_y$.
%  
% \end{lem}
% \begin{proof}
%  Let $\hat x$ be a
% \end{proof}


\section{Singular $(G,X)$-manifolds (II) : tame and wild}\label{sec:GXsing2}

In this section we utilize elements laid out on branched covering to construct singular equivalents to tools of the standard theory of $(G,X)$-manifolds. A fundamental property of regular $(G,X)$-manifolds is the developement Theorem: let $M$ be a connected $(G,X)$-manifold and let $\widetilde M$ be its universal covering, there exists a couple $(\D,\rho)$ with $\rho: \pi_1(M)\rightarrow G$ and $\D:\widetilde M\rightarrow X$ a $\rho$-equivariant $(G,X)$-morphism; furthermore the couple is unique up to the action of $G$. 

One of the goal is to allow for efficient proof of uniformization results. For instance, if $\D$ is injective $\widetilde M$ naturally identifies \textbf{set-wise} to a domain in $X$ and we may identify set-wise $M$ to $\rho\backslash\D(\widetilde M)$. As will become clear in the examples given in the last section of the present work, topology is a core difficulty. We may  define a developing map a merely an a.e. $(G,X)$-morphism from $\widetilde M$ to $X$, this leads to {\it tamely singular} $(G,X)$-manifolds. 

\begin{ex} \label{ex:Tame}
 Let $n\in \N$ and $\theta\in \RR_+^*\setminus \{2\pi\}$ and  $M = \mathbb E^{n+2}_{\theta}$ so that $\sing(M) = \{r=0\}$ in cylindrical coordinates. 
 We observe that $\widetilde{\mathbb E}^{n+2}_{\theta} := \widetilde M^{\sing(M)}$ can be identified with $\susp(\RR)\times \RR^{n}$.  In cylindrical coordinates $(r,\theta,z_1,\cdots,z_{n}) \in \RR_+\times \RR\times \RR^{n}$, the covering $\widetilde{\mathbb E}^{n+2}_{\theta}$ can be endowed with the singular Euclidean metric $\d r^2 + r^2 \d \theta^2 + \sum_{i=1}^{n}\d z_i^2$ and we can define the developing map 
 $$\fonction{\D}{\widetilde{\mathbb E}^{n+2}_{\theta} }{\E^{n+2}}{(r,\theta,z_1,\cdots,z_{n})}{ (r\cos\theta,r\sin(\theta),z_1,\cdots,z_n)}.$$
\end{ex}

However, the following example  shows that many interesting examples are not tamely singular.


\begin{ex}\label{ex:WTame}
 Using notation from \ref{ex:quotient_base}, the extreme BTZ-like whitehole model has maximal branched covering but no developing map. Indeed, one can observe that the developing map 
$$\fonction{\D}{\reg(\cyl{3})}{\E^{1,2}}{\left(\tau,\r,\frac{\theta}{2\pi}\right)}{
\begin{pmatrix}t \\ x \\ y\end{pmatrix}=\begin{pmatrix} \tau + \frac{1}{2}\r\theta^2 \\ \tau + \frac{1}{2}\r\theta^2 -r \\ -\r\theta  \end{pmatrix}
}.$$
does not extends continuously at $\r = 0$. For instance, the sequence $(\tau_n,\r_n,\theta_n)_{\N^*}$ with 
$$\forall n\in \N^*, \quad \tau_n=0,~\r_n=\frac{1}{n},~ \theta_n = n$$
is such that $(\tau_n,\r_n,\theta_n)\xrightarrow{n\rightarrow +\infty} (0,0,0)$ in $\cyl{3}$ but $\D(\tau_n,\r_n,\theta_n)$ is unbounded.
\end{ex}

 Aiming at a uniformization for this last example  leads to a notion of {\it virtually tame} singular $(G,X)$-manifolds by relaxing the continuity assumption on the developing map. 
 In this context we are able to construct uniformization results.

Of course, a natural definition of a developing map (even discontinuous) is not guaranteed in general as the following example shows.

\begin{ex}\label{ex:eucl_sphere} Let $M = \mathbb S^2$  and $N\in M$ endowed with a singular $(\isom^+(\E^2),\mathbb E^2)$-structure induced by a stereographic projection $\mathbb S^2\setminus \{N\} \rightarrow \mathbb E^2$. 
On the one hand, $M$ is a singular Euclidean manifold with $\sing(M)=\{N\}$ and  $\widetilde M^{\sing(M)} = M$. On the other hand, the  developing map $\D:\reg(M)\rightarrow \mathbb E^2$ given by  the stereographic projection  cannot be extended continuously and no sequence $x_n\xrightarrow{n\rightarrow +\infty} N$ with $\forall n\in\N, x_n\in \reg(M)$ is such that $(\D(x_n))_{n\in \N}$ converges.
\end{ex}

The present section lays out distinctions between singularities in $(G,X)$-manifolds based on the existence or not of developing maps. We thus introduce notions of tame and virtually tame singular $(G,X)$-manifolds, the section revolves around the localization
of these notions. Localization achieved in Theorems \ref{theo:local_tame} and \ref{theo:local_virt_tame}.



In the whole section, we need a guarantee that we are allowed to consider the maximal covering possibly branched over any skeletal subset of any open subset. We thus make an extra assumption on $X$.
\begin{defi}\label{def:topological_space_locally_unloopable} A topological space $X$ is locally unloopable if it is locally path-connected and if for all 
$x\in X$ and all neighborhood $\U$ of $x$, there exists an open neighborhood $\V\subset \U$ of $x$ which is semi-locally simply connected.
\end{defi}
\begin{lem}\label{lem:analytical_structure_locally_unloopable} Let $(G,X)$ be an analytical structure with $X$ locally unloopable.  Then every connected singular $(G,X)$-manifold $M$ admits a maximal covering possibly branched over some given skeletal subset $S$ as long as $S\supset \sing(M)$.
\end{lem}
\begin{proof} Since $S\supset \sing(M)$, we have
 $M\setminus S = \reg(M)\setminus S$. Since $X$ is locally unloopable and since $\reg(M)\setminus S$ is locally homeomorphic to $X$, in particular $\reg(M)\setminus S$ is locally unloopable. Since $M$ is connected and $\reg(M)\setminus S$ portly then $\reg(M)\setminus S$ is connected thus unloopable. 
\qed\end{proof}

\textbf{
In the whole section $(G,X)$ is an analytical structure with $X$ locally unloopable.}


% \footnote{The developing property is a key ingredient in the classification of closed manifolds of dimension $n\leq 3$. Though there are other ways for manifolds of dimension 1 or 2, to the author knowledge, it is the only known route for the classification of 3-manifolds. With 4 "general steps":
% the description of natural analytical structures $(G,X)$, a (very difficult) geometrization step which endow (prime) manifolds  with one of the natural $(G,X)$-structures, a uniformization step describing any manifold as a quotient $\Gamma\backslash\Omega$ for some $\Omega\subset X$, finally a classification of couple $(\Gamma,\Omega)$. 
% The author finds the use of this route for showing that closed 1-manifolds are circle very enlightning.}
% 
% 
%                   

\subsection{Ramified coverings, group actions}

In this section, we prove natural technical properties one expects to hold for the theory to work smoothly.




  \begin{prop}\label{prop:singular_manifold_branched_covering} Let $M$ be a singular $(G,X)$-manifold and let $N\xrightarrow{\pi} M$ be a branched covering.
                    
                    $N$ has a unique a.e. $(G,X)$-structure for which $\pi$ is an a.e. $(G,X)$-morphism.
                    and $\reg(\pi)\supset\reg(N)\cap \ord(\pi)$.
                    
                     Furthermore, the Galois group $\Gamma(\widetilde M^S/M)$ acts by a.e. $(G,X)$-morphisms.
                    \end{prop}
                    \begin{proof}
                        By Lemma \ref{lem:aepullback} there exists a unique a.e. $(G,X)$-structure on $N$ such that
                        $\pi$ is an a.e. $(G,X)$-morphism. Lemma \ref{lem:regf1} implies that $$\reg(\pi)=\reg(N)\cap\{x ~|~ f \text{ local homeo. at }x\}\supset \reg(N)\cap \ord(\pi)$$
                     
                     The Galois group $\Gamma(N/M)$ acts by homeomorphisms on $\widetilde M^S$ and $\pi$ is an a.e. local homemorphism. Denoting $\A,\B$ the a.e. $(G,X)$-atlases of $N,M$ respectively, we may thus apply Proposition \ref{lem:aepullback} 
                     $$\gamma^*\A = \gamma^* \pi^*\B=(\pi\circ \gamma)^*\B = \pi^*\B = \A.$$
                     In other words, $\gamma$ is an a.e. $(G,X)$-morphism. 
                    \qed\end{proof}

					\begin{cor}\label{cor:singular_manifold_skeletal_subset} Let $M$ be a singular $(G,X)$-manifold and let $S$ be a skeletal subset of $M$ containing $\sing(M)$. Then $\widetilde M^S$ admits a unique a.e. $(G,X)$-structure such that the natural projection $\widetilde M^S\rightarrow M$ is an a.e. $(G,X)$-morphism.
					
					Furthermore, the Galois group $\Gamma(\widetilde M^S/M)$ acts by a.e. $(G,X)$-morphisms. 
					 
					\end{cor}








\subsection{Developing map and tameness}


In this section $M,N$ will denote connected singular $(G,X)$-manifolds. 
For brievety sake, whenever it makes sense, we will denote by $\widetilde M$ the maximal covering of $M$ possibly branched over $\sing(M)$ which can also be denoted by $\widetilde M^{\sing(M)}$ in the notation of the previous section. 




% 
% \edt{
% \begin{prop} Consider $\mathcal C$ the category of singular $(G,X)$-manifolds whose morphisms are local homeomorphisms.
% The application $M\mapsto \widetilde M$ is a fully faithful functor from $\mathcal C$ to itself.
% \end{prop}
% \begin{proof}
%  Let $M_1\xrightarrow{\phi} M_2$ be an a.e. $(G,X)$-morphism which is a local homeomorphism with $M_1$ and $M_2$ two $(G,X)$-manifolds. From Proposition \ref{prop:GX_sing_morph_local}, $\phi$ sent regular locus to regular locus; then its restriction to the regular locii $\phi$ lifts to a map $ \widetilde \reg(M_1)\xrightarrow{~~~} \widetilde \reg(M_2)$. Since spread completion is functorial, $\phi$ lifts  to a map $\widetilde M_1\rightarrow \widetilde M_2$.
% \end{proof}
% }

\subsubsection{Developing maps}
\begin{defi}[developing map]\label{def:developing_map}
  Let $N\xrightarrow{p} M$ be a branched covering of singular $(G,X)$-manifold.
  A developing map $\D$ of $N\rightarrow M$ is any a.e. $(G,X)$-morphism $N \rightarrow X$. 
  
  We say that $M$ admits a developing map if there exists a branched covering $\widehat M \rightarrow M$ possibly branched above $\sing(M)$ and a developing map of $\widehat M \rightarrow M$
  
   For simplicity sake, whenever it makes sense, a developing map of $\widetilde M\rightarrow M$ is simply called a developing map of $M$. 
  \end{defi}





% \begin{lem} 
%   Let $N\xrightarrow{p} M$ be a branched covering of singular $(G,X)$-manifold and let $\D$ be a developing map  of $N\rightarrow M$. Then $\reg(p)\subset \reg(\D)$.
% \end{lem}
% \begin{proof}
%  \edt{pas tout à fait certain TBD}
% \end{proof}


Up to reduction to portly subset, the proofs of the following three Lemmas are identical to the usual proof in the regular context and are thus skipped.  
\begin{lem}\label{lem:MtoXG}
  Let $N\xrightarrow{f_1,f_2} X$ be a.e.  $(G,X)$-morphisms,
then there exists a unique $g\in G$ such that $f_1=gf_2$.
\end{lem}

\begin{lem} \label{lem:holonomy} Let $M$ be a singular $(G,X)$-manifold and $\Gamma$ be a group acting on $N$ by a.e. $(G,X)$-morphisms. 

For any a.e. $(G,X)$-morphism $M\xrightarrow{f} X $  there exists a unique morphism $\rho:\Gamma\rightarrow G$ such that $f$ is $\rho$-equivariant.
\end{lem}
% \begin{proof}
%  For $\gamma \in \Gamma(N/M)$, the map $\D\circ \gamma$ is also a developing map of $N\rightarrow M$. Let $x\in \reg(N)$ and let $(\U,\V,\varphi)$ be a chart around $x$. 
%  There exists $g,h \in G$ such that $\forall y\in \V, \D\circ \varphi^{-1}(y) = g(y)$ and  $\D\circ \gamma\circ \varphi^{-1}(y)= h(y)$. Therefore, $\forall u\in \U,\  \D\circ \gamma(u) = h\circ g^{-1} \circ \D(u) $; then, by analyticity, $\D\circ \gamma = h\circ g^{-1}\circ \D$ ie $\forall \gamma\in \Gamma(N/M),\exists g_\gamma\in G, \D\circ \gamma = g_\gamma\circ \D$. Note that such a $g_\gamma$ is unique and we can define $\rho : \Gamma(N/M)\rightarrow G, \gamma \mapsto g_\gamma$. Note that $$\rho(\gamma_1\gamma_2)\circ \D = \D\circ (\gamma_1\gamma_2)= \D\circ \gamma_1\circ \gamma_2 = \rho(\gamma_1)\circ \D\circ \gamma_2 = \rho(\gamma_1)\circ \rho(\gamma_2)\circ \D $$ and $\rho$ is thus a morphism.
% \end{proof}
\begin{defi}\label{def:quasi_galoisian_branched_covering_singular} 
 Let $N\xrightarrow{p} M$ be a quasi-Galoisian branched covering of singular $(G,X)$-manifold admitting a developing map $\D$. The holonomy of $N\rightarrow M$ (associated to $\D$) is the morphism $\rho$ of Lemma \ref{lem:holonomy}.
\end{defi}
 \begin{rem}\label{rem:developing_map_quasi_galoisian_covering}
    A developing map $\D$ of a quasi-Galoisian covering $M\rightarrow N$ sends points to some fix point of their stabilizer : $\D(x) \in \fix(\rho(\Gamma_x))$.
  \end{rem}
  
\begin{rem}\label{rem:singular_manifold_holonomy_indeed} 
 Let $M$ be a singular $(G,X)$-manifold. The holonomy of $\widetilde M \rightarrow M$ is indeed the holonomy of $\reg(M)$ in the usual sense.
\end{rem}
\begin{lem}\label{lem:quasi_galoisian_branched_covering_singular}
 Let $N\xrightarrow{p} M$ be a quasi-Galoisian branched covering of singular $(G,X)$-manifold. If it admits a couple $(\D,\rho)$ of developing map and holonomy, then all other such couples are obtained via conjugation ie:
 if $(\D',\rho')$ is another couple developing map and holonomy, then there exists a unique $g\in G$ such that $(\D',\rho')= (g\circ\D,\rho^g)$ with $\rho^g  : \gamma \mapsto g\rho(\gamma)g^{-1}$.
\end{lem}


%    @Guiserix Sacha  Par contre ce qui est vrai, c'est que j'ai le luxe de ne pas être focalisé sur l'immédiat. Cela me laisse plus de marge pour apprécier des propositions de fond qui auront un effet dans 4-5 ans. 

  
\begin{lem}\label{lem:descend_developing_map} Let $M\xrightarrow{p} N$ be a quasi-Galoisian branched covering of singular $(G,X)$-manifolds. Then, the following are equivalent:
\begin{enumerate}[(i)]
 \item there exists an a.e. $(G,X)$-morphism $M\rightarrow X$ and the holonomy of $M\rightarrow N$ is trivial,
 \item there exists an a.e. $(G,X)$-morphism $N\rightarrow X$. 
\end{enumerate} 
\end{lem}
\begin{proof}
 Assume $(i)$ and consider some a.e. $(G,X)$-morphism $M\xrightarrow{\D} X$ as well as the equivalence relation $x\sim y$ if $p(x)=p(y)$ for $x,y\in M$.
 By Lemma \ref{lem:dense_orbits}, the action of $\Gamma(M/N)$ has dense 
 orbit in the fiber of $p$ then the equivalence classes of $\sim$ are 
 the closure of orbits of $\Gamma(M/N)$. Furthermore, since $\rho=1$ 
 then $\D$ is constant on the orbit of $\Gamma(M/N)$ and continuous, hence constant on the equivalence classes of $\sim$. Therefore, $\D$ induces a continuous map $\sim\backslash M \rightarrow X$. Moreover, $p$ induces a continuous map $\sim\backslash M \rightarrow N$ which is open by Corollary \ref{cor:spread_connected_image}, injective by definition and surjective since $p$ is surjective; hence $\sim\backslash M \rightarrow N$ is an homeomorphism. We thus constructed a continuous map $\overline \D:N \rightarrow X$. Notice that $\overline \D$ is a $(G,X)$-morphism on $\ord(N)$ by usual results on $(G,X)$-manifolds. Finally, $\D$ is an a.e. $(G,X)$-morphism. 
 
 Assume $(ii)$ and consider some $N\xrightarrow{\overline \D} X$, the map $ M \xrightarrow{\overline \D \circ p} X$ is thus an a.e. $(G,X)$-morphism and 
 $$\forall \gamma\in \Gamma(M/N),\quad \D\circ\gamma =  \overline \D\circ p\circ \gamma = \overline \D\circ p = \D $$
 so $\rho= 1$ i.e. the holonomy of $M\rightarrow N$ is trivial. 
\qed\end{proof}


\begin{lem}\label{lem:morphism_singular_manifolds_branched}
 Let $f:M\rightarrow N$ be a morphism of singular $(G,X)$-manifolds and let $\hat M$ and $\hat N$ be branched coverings of $M$ and $N$ respectively. 
If $\hat f:\hat M \rightarrow \hat N$ if a lift of $f$, then $\hat f$ is a morphism of singular $(G,X)$-manifolds.
\end{lem}
\begin{proof}
One can restrict to portly subsets to reduce the Lemma to the case were $M,N$ are regular, $\hat M,\hat N$ are unbranched coverings and $f$ is a $(G,X)$-morphism. Standard results apply. 
\qed\end{proof}

\begin{prop} \label{prop:branched_devmap}Let $M$ be a singular $(G,X)$-manifold. The following are then equivalent:
\begin{enumerate}[(i)]
 \item there exists an a.e. $(G,X)$-morphism $\widetilde M\rightarrow X$,
 \item there exists a branched covering of $N\xrightarrow{p} M$ which admits a developing map.
\end{enumerate} 
\end{prop}
\begin{proof}
 Clearly, $(i)\Rightarrow (ii)$.  Assume $(ii)$ and consider some branched covering $N\xrightarrow{p}M$  and a developing map $N\xrightarrow{\D} X$. Denote by $S:=M\setminus \ord_p(M)$ and $\widetilde S$  the lift of $S$ to $\widetilde M:=\widetilde M^{\sing(M)}$.   
 Since  $S\subset S\cup \sing(M)$, by Proposition \ref{prop:universality}, the  map $N \xrightarrow{\D} X$ lifts to an a.e. $(G,X)$-morphism $\widehat M := \widetilde M^{S\cup \sing(M)}\xrightarrow{~~\widehat \D~~} X$. 
 The developing map of $\reg(M)$ is a $(G,X)$-morphism $\widetilde{\reg}(M)\xrightarrow{\D_0} X $ thus a $(G,X)$-morphism $ \widehat M' \xrightarrow{\widehat \D_0} X$ were $\widehat M'$ is the universal covering of $\widetilde{\reg}(M)\setminus \widetilde S$ which embeds naturally in $\widehat M$ since it is the universal covering of $\reg(M)\setminus S$. By Lemma \ref{lem:descend_developing_map}, the holonomy 
 $$\Gamma:=\Gamma(\widehat M'/\widetilde\reg(M)) = \Gamma(\widehat M/\widetilde M) \xrightarrow{~~~\rho~~~} G$$ of $ \widehat M' \rightarrow \widetilde{\reg}(M)$ is trivial. 
 Futhermore, there exists $g\in G$ such that $g\circ\widehat \D_{|\widehat M'} = \widehat \D_0$ ; therefore $\widehat \D_0$ extends continuously to $\widehat M \xrightarrow{\widehat \D_0} X$.  Since the action of $\Gamma$ on $\widehat\D_{0|\widehat M'}$ is trivial  and  since $\widehat M'$ is dense in $\widehat M$, the action of $\Gamma$ on $\widehat \D_0$ is trivial. We can thus apply Lemma \ref{lem:descend_developing_map} again to obtain a continuous extension of $\D_0$ to $\widetilde M$.
\qed\end{proof}


\subsubsection{Tame singularities}


Following the discussion of the previous sections, we introduce the notion of tame singular $(G,X)$-manifolds. 
\begin{defi}\label{def:singular_manifold_tame_admits}
 A singular $(G,X)$-manifold $M$ is tame if it admits a  developing map.
\end{defi}
Since the notion of developing map is global, we need to localize tameness to ease manipulation of this property.

\begin{lem}\label{lem:embedding_tame}
Let $M,N$ be singular $(G,X)$-manifolds and let $M\rightarrow N$ be an a.e. $(G,X)$-morphism. If $N$ is tame then $M$ is tame.
\end{lem}
\begin{proof}
 The morphism $M\rightarrow N$ lifts to a morphism $\widetilde M\rightarrow \widetilde N$. We can compose the latter by the developing map $\widetilde N \rightarrow X$ of $N$ to obtain a developing map of $M$. 
\qed\end{proof}
\begin{cor}\label{cor:branched_covering_tame_singular} A branched covering of a tame  singular $(G,X)$-manifold is tame.
\end{cor}
\begin{cor}\label{cor:open_subsets_tame_singular} Open subsets of tame singular $(G,X)$-manifolds are tame.
\end{cor}



\begin{thm} \label{theo:local_tame}  Let $M$ be a singular $(G,X)$-manifold the following are equivalent:
\begin{enumerate}
 \item $M$ is tamely singular 
 \item $M$ locally tamely singular.
\end{enumerate}



\end{thm}


\begin{proof} Denote by $\widetilde M \xrightarrow{p} M$ the maximal covering of $M$ ramified above $\sing(M)$.


Assume $M$ tamely singular, for any $x\in M$ and any neighborhood $\U$ of $x$, by local connectivity, there exists an connected open neighborhood $\V\subset \U$  of $x$. Take any connected component $\W$ of $p^{-1}(\V)$ so that $\W\rightarrow \V $ is a branched covering, by universality there exists a branched covering $\widetilde \V\xrightarrow{q} \W$. Take a developing map $\D$ of $M$, the map $\D_{|\W}\circ q$ is a developing map of $\V$ hence $\V$ is tamely singular.

Assume $M$ locally tamely singular.
Let $\widetilde x\in \widetilde M$ and let $\U$ be an open connected neighborhood of $x:=p(\widetilde x)$ admitting a developing map. Define $\widetilde \U \xrightarrow{\pi} \U$ the natural branched covering. 
  The natural injection $\U \xrightarrow{\iota} M$ lifts to an a.e. $(G,X)$-morphism $\widetilde \U\xrightarrow{\widetilde \iota} \widetilde M$. 
  Let $\reg(\widetilde M) \xrightarrow{\D_M} X$ be  a $(G,X)$-morphism,  the map $\reg(\widetilde \U) \xrightarrow{\D_M\circ \widetilde \iota} X$ is a $(G,X)$-morphism and extends continuously to an a.e. $(G,X)$-morphism $\widetilde \U\xrightarrow{\D_\U} X$. Let $\hat \U$ be the connected component of $\widetilde \iota (\widetilde \U)$ in $p^{-1}(\U)$; by Lemma \ref{lem:branched_covering_restriction},  $\widehat \U \xrightarrow{p} \U$ is a branched covering, by Proposition \ref{prop:universality} there thus exists a  map $\widetilde \U \xrightarrow{\pi'} \widehat \U$  such that $p\circ \pi' = \pi $. Since $\widetilde \U/\widehat\U$ is quasi-Galoisian, this map is unique up to action of $\Gamma(\widetilde \U/\widehat \U)$. Since $p \circ \widetilde \iota = \pi$, we conclude that $\pi' = \widetilde \iota \circ \gamma$ for some $\gamma \in \Gamma(\widetilde M/\widehat M)$. In particular, $\widetilde \iota^{|\widehat \U}$ is a branched covering.
  
  Using same argumentation as in the proof of Proposition \ref{prop:branched_devmap} we show that $\D_{\U}$ descends to an a.e. $(G,X)$-morphism  $\widehat \U \rightarrow X$ which extends continuously $\D_{M|\widehat \U}$.  Finally $\D_M$ extends continuously at $\widetilde x$ and, since $\widetilde x$ is arbitrary, to the whole $\widetilde M$.
 \qed\end{proof}
 
 The following remark though trivial is important.
 \begin{rem}\label{rem:tamely_singular_manifold_admits} Let $M$ be a tamely singular $(G,X)$-manifold. If $M$ admits an injective developing map $N\xrightarrow{\D}X$ for some branched covering $N$ of $M$ then $M$ is isomorphic to $\Gamma(N/M)\backslash\D(N)$ where 
 \begin{itemize}
  \item $\Gamma(N/M)$ acts on $X$ via the holonomy of $\D$;
  \item $\D(N)$ is endowed with the topology making  $\D$ an homeomorphism.
 \end{itemize} 
 \end{rem}

 The specification of the topology is important as the developing map $\D$, though continuous, may not be open in neighborhoods of singular points.
 \begin{ex}\label{ex:take_closed_unit_disc}
  Take $X$ the closed unit disc in $\RR^2$ and identify the open disc with the pointcaré disc $\H$. Take $G$ the group of isometries of the pointcaré disc. Note that $G$ acts by homeomorphisms on $X$.
  Consider $M$ a genus $g$ surface and $S\subset M$ a finite subset such that $2g-2+s>0$. Endow $M\setminus S$ with a finite volume complete hyperbolic metric. 
  On the one hand, $M$ is a tamely singular $(G,X)$-manifold with injective developing map $\widetilde M \xrightarrow{\D} X$. On the other hand, an open neighborhood $\U$ of a singular point $x$ of $\widetilde M$ is homeomorphic to a domain $\{r<r_0\}$ in the  ``infinite angle disc'' $\DD_\infty$ introduced in Example \ref{ex:quotient_base}. We may choose $\U$ so that $\D(U)$ is a horodisc centered at $\D(x)$. Such a domain is not open for the topology of $X$.
  
  
 \end{ex}





\subsection{Singular models}

In common settings, singularities are not arbitrary : they are assumed to be isomorphic to some model such as $(\mathbb E^2_{\theta})_{\theta>0}$ for singular locally Euclidean surfaces or $(\E^{1,2}_\alpha)_{\alpha>0}$ for flat 3-spacetimes with massive particles.

\subsubsection{Analytic family}

\begin{defi}\label{def:singular_manifold_automorphisms_group}
 Let $M$ be a singular $(G,X)$-manifold. The automorphisms group of $M$ denoted $\Aut(M)$ is the group of isomorphisms of singular $(G,X)$-manifolds from $M$ to itself.
\end{defi}

\begin{defi}\label{def:family_singular_manifolds_endowed} A family $(X_\alpha)_{\alpha \in A}$ of singular $(G,X)$-manifolds each endowed with a non-empty subset $S_\alpha\subset \sing(X_\alpha)$ is an analytic family of model spaces if for all $\alpha_1,\alpha_2\in A$ and all a.e. $(G,X)$-isomorphism $\U_1\xrightarrow{\varphi} \U_2$ where $\U_1,\U_2$ are connected open subsets of $X_{\alpha_1}$ and $X_{\alpha_2}$ respectively we have:
$$\U_1\cap S_{\alpha_1}\neq \emptyset \quad \Rightarrow \quad \alpha_1=\alpha_2 ~~\text{and}~~\exists \phi \in  \mathrm{Aut}(X_{\alpha_1}), \varphi=\phi_{|\U_1}^{|\U_2}.$$


\end{defi}

\begin{defi}[$X_A$-atlas]\label{def:atlas}
  Let $(X_\alpha)_{\alpha\in A}$ be an analytic family of model spaces and let $M$ be a singular $(G,X)$-manifold.
    
  For $\alpha\in A$, an $X_\alpha$-chart around $x\in \sing(M)$ is a triplet $(\U,\V,\phi)$ with $x\in \U\subset M,\phi(x)\in S_\alpha, \V\subset X_\alpha$ open and connected;  and $\U\xrightarrow{\phi} \V$ an isomorphism. If such a chart exists, we say that $x$ has type $\alpha$. 
  
  An $X_A$-chart is an $X_\alpha$-chart for some $\alpha\in A$.  
  An $X_A$-atlas is a collection of $X_A$-charts covering $M$.
\end{defi}
\begin{defi}\label{def:analytic_family_model_spaces}
  Let $(X_\alpha)_{\alpha\in A}$ be an analytic family of model spaces and let $M$ be a singular $(G,X)$-manifold.
  
  A $X_A$-manifold is a singular $(G,X)$-manifold $M$  admitting  a $X_A$-atlas
\end{defi}



\begin{rem}\label{rem:singular_point_manifold_has}
    A singular point $x\in \sing(M)$ of a $X_A$-manifold $M$ has exactly one type. Indeed, let $\U_i\xrightarrow{\varphi_i} \V_i\subset X_{\alpha_i}$ be isomorphisms with $\U_i,\V_i$ open for $i\in \{1,2\}$ and with $x\in \U_1\cap \U_2$. Restricting $\varphi_1,\varphi_2$ if necessary, we may assume $\U_1=\U_2$ connected  so that $\V_1\xrightarrow{\varphi_2\circ \varphi_1^{-1}}\V_2$ is an isomorphism, Moreover, $\V_1$ contains the singular point $\varphi_1(x)\in S_{\alpha_1}$ so by analyticity $\alpha_1=\alpha_2$.
\end{rem}
\begin{rem}\label{rem:theorem_theo_local_tame}
 By Theorem \ref{theo:local_tame}, a $X_A$-manifold is tame if and only if all models are tame.
\end{rem}
\begin{ex}\label{ex:follows_proposition_brunswic_cauchy} It follows from Proposition 1.3 of  \cite{brunswic2021cauchy} that the family $(\mass{\theta})_{\theta\geq 0}$ is analytic.
 
\end{ex}
\begin{cex}\label{cex:start_singular_manifold_homeomorphic} Start from $\Sigma$ a singular $\SS^2$-manifold homeomorphic to a sphere with $s\geq 3$ conical singularities of respective angles $(\theta_i)_{i\in \lsem 1,s\rsem}$. We know from Troyanov \cite{MR1005085}  that one can construct such a surface for arbitrary  $(\theta_i)_{i\in \lsem 1,s\rsem}$ as long as $$0<\sum_{i=1}^s( \theta_i -2\pi) + 4\pi\leq \min(4\pi, \theta_i : i\in \lsem 1,s\rsem).$$
Apply the Thurston suspension method described in Section \ref{sec:susp}, the constructed manifold $M$ is  a singular $(\SO(4),\SS^3)$-manifold (we may either apply  by Proposition \ref{prop:susp_portly}  or check that directly). We call $O$ its center.

An automorphism of $M$ has to fix $O$ since singular locus is homeomorphic to a line except at $O$, furthermore it preserves distances and geodesics. By considering the sphere of some radius $\varepsilon>0$ centered at $O$ we see that the natural map $\Aut(\Sigma)\rightarrow \Aut(M)$ is a group isomorphism.  As a consequence, except for very particular $\Sigma$, $\Aut(M)$ is trivial. For instance, if all the $\theta_i$ are different then $\Aut(M)=\{Id\}$. These remarks hold for all neighborhoods of $O$.

For a generic $\Sigma$, since locally we may construct a rotation around any of the axes and since  those cannot be extended as automorphism of $M$, we conclude that for $S\subset \sing(M)$, $(M,S)$ is  analytic if and only if $S=\{O\}$.

 
\end{cex}


\subsubsection{A motivational example : the BTZ model space}
\label{sec:mot_BTZ}
The BTZ model space $\mass{0}$ is defined as $\RR^3$ endowed with the singular $\mass{}$-structure induced by the flat lorentzian metric $\d s^2=-2\d\tau \d \r +\d \r^2+ r^2 \d \theta$ in cylindrical coordinates.
The singular locus of $\mass{0}$ is then $\sing(\mass{0})=\{\r=0\}$ and the regular locus 
is $\reg(\BTZ):= \{\r>0\}$. 
Following example \ref{ex:quotient_base}, from which we reuse notations, we denote  the maximal covering of $\BTZ$ possibly branched over $\sing(\BTZ)$  by $\BTZtilde$. Since $\BTZ$ is simply $\RR^3$ and $\sing(\BTZ)$ is a simply a line, $\BTZtilde$ is homeomorphic to the maximal covering of $\RR^3$ branched over a line denoted by  $\cyl{3}$ and parameterize using cylindrical coordinates $(\tau,\r,\theta)\in \RR\times\RR_+\times \RR$. A natural candidate for a developing map is given by 
$$\fonctiondeux{\D}{\BTZtilde}{\E^{1,2}}{\left(\tau,\r,\frac{\theta}{2\pi}\right)}{
\begin{pmatrix}t \\ x \\ y\end{pmatrix}=\begin{pmatrix} \tau + \frac{1}{2}\r\theta^2 \\ \tau + \frac{1}{2}\r\theta^2 -\r \\ -\r\theta  \end{pmatrix}
}{(\tau,0,\cdot)}{\begin{pmatrix}\tau \\ \tau \\ 0\end{pmatrix}}$$
which can be checked to be a $\E^{1,2}$-morphism on $\reg(\BTZtilde)=\{\r>0\}$ but is discontinuous at every points of $\sing(\BTZtilde)=\{\r=0\}$. Furthermore, one may check that for all  $x=(\tau,0,0)\in \sing(\BTZtilde)$ we have  $$\bigcap_{\U\ni x}\overline {\D(\reg(\U))}=\{(\lambda,\lambda,0)~:~\lambda \geq \tau\}$$ so that there is no way to extend $\D$ continusouly to $\{\r=0\}$ keeping the standard topologies of $\BTZtilde$ and $\mass{}$.

Consider the topology $\T$ of $\BTZtilde$ generated by the standard topology $\T_0$ and the sets of the form $$\U_\lambda = \left\{(\tau,\r,\theta)\in \BTZtilde ~|~ \r\theta^2< \lambda  \right\},\quad \lambda \in \RR_+^*.$$
Then $\D:(\BTZtilde,\T)\rightarrow \mass{}$ is continuous  and element of $\Aut(\BTZtilde)$ are $\T$-continuous, furthermore the quotient topology induced on $\BTZ$ is the usual topology.
 Finally, we note that $\reg(\BTZtilde)$ is $\T$-portly.



\subsubsection{Virtually tame models}
Some model spaces of interest such as the BTZ model space are not tame  but are almost tame in the sense that we may extend the developing map in a unique fashion allowing to construct developing map on $X_A$-manifolds. 
\begin{defi}\label{def:singular_manifold_virtually_tame} A singular $(G,X)$-manifold $M$ is virtually tame if there exists  a first countable locally path connected (strong) topology $\T$ on $\widetilde M$ thinner than the natural (weak) topology such that :
\begin{enumerate}[(A)]
\item The quotient topology induced by $\T$ on $M$ is the natural one.
 \item the developing map $\D:\reg(\widetilde M)\rightarrow X$ admits a $\T$-continuous extension ;
 \item elements $\Gamma(\widetilde M/M)$ are $\T$-continuous;
 \item  the lift of $\reg(M)$ to $\widetilde M$ is $\T$-portly in $\widetilde M$.
\end{enumerate}
We call such a $\T$ a tamifying topology.
\end{defi}
\begin{defi}\label{def:singular_manifold_locally_galoisian}A singular $(G,X)$-manifold $M$ is locally Galoisian  
if for every $\U\subset \M$ open, $\reg(\U)$ is semi-locally simply connected in $\U$.
\end{defi}
The goal of this section is to prove the following localization property.

\begin{thm}\label{theo:local_virt_tame} Let $(X_\alpha)_{\alpha\in A}$ be an analytical family of locally Galoisian model spaces.
If for every $\alpha\in A$, there exists a tamifying topology $\T_\alpha$ such that the group $\Aut(\widetilde X_\alpha)$ acts $\T_\alpha$-continuously,
 then the following are equivalent:
\begin{enumerate}[(i)]
 \item every $X_A$-manifold is virtually tame
 \item for every $\alpha \in A$, $X_\alpha$ is virtually tame
\end{enumerate}
\end{thm}

To this end, we will need technical results on branched covering involving two topologies on the same underlying space. In every instance, we shall refer to the strong and weak topologies, the former being the thinner of the two.
\begin{defi}\label{def:topological_spaces_continuous_map} Let $(M,\T_M),(N,\T_N)$ be topological spaces and Let $M\xrightarrow{p} N$ be a continuous map. For any topology $\T$ on $N$ define $p^*\T$ the topology of $M$ generated by $\T_M$ and the $\T_M$-connected components of $p^{-1}(\O)$ for $\O\in \T$. 
 
\end{defi}



\begin{lem} \label{lem:spread_top} Let $M$ be a first countable, connected locally path connected Hausdorff topological space. Let $S\subset M$ be a skeletal subset such that $M\setminus S$ is semi-locally simply connected in $M$;
consider a sub-covering $\widetilde M^S\xrightarrow{p}N\xrightarrow{q}M$.

Let $\T$ be a (strong) thinner first countable and locally path connected topology on $N$ for which elements of $\Gamma(N/M)$ are continuous. Consider the induced strong topology $p^*\T$ on $\widetilde M^S$.
Assume in addition there exists some $\R \subset \ord_p(N)$ strongly portly  in $N$ on which strong and weak topologies agree.

 Then, we have the following
 \begin{enumerate}[(a)]
 \item $\widetilde M^S$ is strongly first countable connected locally path connected;
 \item strong and weak topologies agree on $p^{-1}(\R)$; 
 \item $p^{-1}(\R)$ is strongly portly;
  \item elements of $\Gamma(\widetilde M^S/M)$ are  strongly continuous;
  \item $\widetilde M^S\xrightarrow{p}N$ is a branched covering for the strong
  topologies.
  \end{enumerate}
\end{lem}
\begin{rem}\label{rem:intuitively_spread_thinner_than}
 Intuitively, if $(M,\T_M)\xrightarrow{p}(N,\T_N)$ is a spread and $\T$ is thinner than $\T_N$, then the strong topology $p^*\T$ is the coarsest topology on $M$ such that $(M,p^*\T)\xrightarrow{p}(N,\T)$ is a spread. However, the definition does not garantee that $\T_M$-connected components are $p^*\T$-connected components. Point $(e)$ above is thus non-trivial.
\end{rem}

\begin{proof} To begin with, we define $\widehat \R := p^{-1}(\R)$ and note that by definition of the ordinary locii, $\widehat \R \xrightarrow{p}\R$ is a covering for the weak topologies. Then, it is not difficult to see that $(p_{|\widehat \R}^{|\R})^*\T_{|\R} = (p^*\T)_{|\R}$ so that point {\it (b)} is clear. We also note that it follows from the definitions that $p$ is continuous for the strong topologies.



To show point {\it (d)}, consider some $\gamma\in \Gamma(\widetilde M^S/M)$ and take some $\overline \gamma\in \Gamma(N/M)$ such that $p\circ \gamma = \overline \gamma\circ p$. For $\widehat \O$ a weak connected component of some $p^{-1}(\O)$ with $\O$ strong open, since $\gamma$ is a weak homemorphism $\gamma^{-1}(\widehat \O)$ is a weak connected component of $(p\circ \gamma)^{-1}(\O)=(\overline \gamma\circ p )^{-1}(\O)$. Since $\overline \gamma$ is strongly continuous $\overline \gamma^{-1}(\O)$ is strongly open so that $\widehat \O$  is strongly open  since it is a weak connected component of $p^{-1}(\overline \gamma^{-1}(\O))$. 


\

\underline{Claim :}  For all strong open $\widehat \O\subset \widetilde M^S$ that is a weak connected component  $p^{-1}(\O)$ for some strong open $\O\subset N$ then $ \widehat \O\cap \widehat \R\neq \emptyset$; Furthermore, for all $\hat y\in \widetilde M^S$ and for all strong neighborhood $\widehat O'\ni \hat y$, there exists $\widehat \O\ni \hat y$  such that $\widehat \O\subset \widehat \O'$ and such that for all $\hat x\in \widehat \O\cap \widehat \R$, there exists a a weakly continuous path $\hat \gamma : \hat x\rightsquigarrow \hat y$ with $\hat\gamma([0,1[)\subset \widehat \R\cap \widehat \O$. 
All such path are strongly continuous.  
We may assume that $\widehat\O\cap \widehat \R$ is connected.



Indeed, take $\widehat \O,\O$ as above and take $\hat y\in\widehat \O$ and $y=p(\hat y)$. Without loss of generality, we may assume $\O$ strongly connected. Since  $N$ is strongly first countable connected locally path connected, since $\O$ is strongly connected and $\R$ strongly portly, $\O\cap \R\neq\emptyset$ so we may choose $x\in \O\cap \R\neq\emptyset$ and by Lemma \ref{lem:path_covering} there exists a strongly continuous path $\gamma : x\rightsquigarrow y$ such that $\gamma([0,1[)\subset \R$. By Proposition \ref{prop:Galoisian_semilocal}, $\widetilde M^S\rightarrow M$ is Galoisian and thus so is $\widetilde M^S\rightarrow N$. There thus exists some weakly continuous lift $\hat \gamma$ of $\gamma$ such that $\gamma(1)=\hat y$. Therefore, $\hat y=\hat\gamma(1)$ and $\gamma(0)$ are in the same weak connected component of $p^{-1}(\O)$ hence $\hat\gamma(0)\in \widehat \O\cap \widehat\R$ eg   $\widehat \O\cap \widehat \R\neq \emptyset$.

Choose $\hat y\in \widetilde M^S$ and $\widehat \O'\ni \hat y$ and define $y:=p(\hat y)$, since $\widetilde M^S\rightarrow N$ is Galoisian, there exists some weak open neighborhood $\U\ni y$ such that the connected component of $p^{-1}(\U)$ has exactly one preimage of $y$ in each weak connected componen. Hence, choosing $\widehat \O$ and $\O$ with the additional property that $\O\subset \U$ we have $\widehat \O\cap p^{-1}(y)=\{\hat y\}$. Then, for  $\hat x\in \widehat \O\cap \widehat \R$ arbitrary and $x:=p(\hat x)$ we may consider the same $\gamma : x\rightsquigarrow y$ as before and then its weakly continuous lift $\hat \gamma$ such that $\hat \gamma(0)=\hat x$. Therefore, $\hat\gamma(1)$ is in the same weak connected component of $p^{-1}(\O)$ as $\hat x$ so in $\widehat \O$ and thus $\hat \gamma(1)=\hat y$.
Since  weak and strong topologies agree on $\widehat\R$, $\hat\gamma$ is strongly continuous on $[0,1[$. Take any open $\widehat \O'\ni \hat y$  being a weak connected component of $p^{-1}(\O')$ for some strong open $\O'$. Since $\gamma$ is strongly continuous, for some $\varepsilon>0$, we have $\gamma([1-\varepsilon,1])\subset \O'$ so that $\hat\gamma([1-\varepsilon,1])\subset p^{-1}(\O') $. Therefore for all $t\in [1-\varepsilon,1[$, $\hat \gamma(t)$ and $\hat\gamma(1)=\hat y$ are in the same weak connected component of $p^{-1}(\O')$ thus $\hat\gamma([1-\varepsilon,1])\subset \widehat \O'$. The path $\hat \gamma : \hat x\rightsquigarrow \hat y$ is thus strongly continuous.

Finally, if we consider in addition some $\hat z\in \widehat \O\cap \widehat \R$ and $z:= p(\hat z)$. Since $x,z \in \O
\cap \R$ which is path connected, we may choose a path $\gamma :x \rightsquigarrow z$ which lifts to a path $\hat x\rightsquigarrow \hat z$ in $\widehat \R \cap \widehat\O$ hence $\widehat \O \cap \widehat \R$ is connected.



\

We now continue the proof of the Lemma.
\begin{itemize}
\item the first part of the claim shows that $\widehat R$ is strongly dense.
 \item Let $\hat y\in \widetilde M^S$ and $y=p(\hat y)$. Take a countable basis of neighborhood $(\O_i)_{i\in I}$ of $y$. A basis of neighborhood of $\hat y$ is provided by weak connected components $\widehat \O_i$ of $\hat y$ in  $p^{-1}(\O_i)$ for $i
 \in I$. Hence, $\widetilde M^S$ is strongly first countable.
 
 \item Let $\hat y\in \widetilde M^S$ and $y=p(\hat y)$. Take $\widehat \O\ni \hat y$  such that the second part of the claim is satisfied (say for $\widehat\O'=\widetilde M^S$). The statement implies that $\{\hat y \}\cup \widehat \O\cap\widehat \R $ is path connected. Take any $\hat z\in \widehat \O$ and choose $\widehat O_1\subset \widehat \O$  such that $\{\hat z \}\cup \widehat \O_1\cap\widehat \R $ is path connected. 
 Since $\emptyset\neq\widehat \O_1\cap\widehat \R\subset \widehat \O\cap\widehat \R$ we may find a strongly continuous path $\hat \gamma$ from $\hat y$ to $\hat z$ via some arbitrary point $\hat x\in \widehat \O_1\cap\widehat \R$. Hence, $\widetilde M^S$ is  strongly locally path connected.
 
 \item Since every point of $\widetilde M^S$ admits a strongly connected  neighborhood, since $\widehat \R$ is weakly  connected (thus strongly connected by point {\it(b)}) and $\widehat\R$ is strongly dense, $\widetilde M^S$ is strongly connected.
 \item $\widehat \R$ is strongly open as the only weak connected component of $p^{-1}(\R)$. It is also dense. 
 
 Let  $\W$ be a strongly connected open subset of $\widetilde M^S$, let $\hat x,\hat y\in \W\cap\widehat \R$ and let $\hat \gamma_0:\hat x\rightsquigarrow \hat y$ be a strongly continuous path. For all $t\in [0,1]$, there exists a strong neighborhood $\widehat\O_t$ of $\hat\gamma(t)$ statisfying the second part of the claim (with $\widehat\O'=\W$). Choose a family of such open sets, by compactness of $[0,1]$ there exists  sequences $t_0=0<t_1<\cdots<t_n=1$ and $\tau \in [0,1]^{\lsem 0,n-1\rsem}$ such that for all $i\in \lsem 0,n-1\rsem$, $\hat \gamma([t_i,t_{i+1}])\subset \widehat\O_{\tau_i}$. For each $i\in \lsem 1,n-1\lsem$, we  choose some $\hat x_i \in \widehat\O_{\tau_{i-1}}\cap \widehat\O_{\tau_i}\cap \widehat \R$ and define $\hat x_0=\hat x$ and $\hat x_n=\hat y$.  We can construct a strongly continuous path $\hat \gamma':\hat x_0\rightsquigarrow \hat x_n$ by choosing for each $i\in \lsem 0, n-1\rsem $ a strongly continuous  path $\hat x_i\rightsquigarrow \hat x_{i+1}$ taking its parameter in $[t_i,t_{i+1}]$ and its values in $\widehat\O_{\tau_i}\cap \widehat \R$.  The path $\hat \gamma'$  takes its values in $\W\cap \widehat \R$ and joins $\hat x$ to $\hat y$, we conclude that $\W\cap \widehat \R$ is path connected.
 
 Finally $\widehat \R$ is strongly portly.

 \item Consider the topology $\T'$ on $\widetilde M^S$ generated by strong connected components of preimage of strong open of $N$. 
 Since  for any strong open $\U$ the strong connected components are included in a weak connected component, the topology $\T'$ is thinner than the strong topology on $\widetilde M^S$. On the other hand, the strong topology on $\widetilde M^S$ is locally connected so connected the generator of $\T'$ are strong open.  We conclude that $p$ is a spread for the strong topologies.
 
 \item Apply Lemma \ref{lem:complete_path_carac} with $X=\widetilde M^S,Y=N, X_1=\widehat \R,Y_1=\R$ for the strong topologies; to show that $p$ is a branched covering for the strong topologies, it suffices to proves that for any path $\gamma : x\rightsquigarrow y$ such that $\gamma([0,1[)\subset \R$, any partially lifts  $\hat \gamma_{|[0,1[}$ extends strongly continuously at $1$. 
 
 Since $p$ is a branched covering for the weak topologies, $\hat \gamma_{|[0,1[}$ admits a weak limit at $1$, we denote $\hat \gamma$ its weak continuous extension so $[0,1]$. The claim ensures that $\hat \gamma$ is actually strongly continuous. We conclude that $p$ is a branched covering.
 
 \end{itemize}
\qed\end{proof}






\begin{proof}[of Theorem \ref{theo:local_virt_tame}]
 $(i)\Rightarrow (ii)$ is clear since each $X_\alpha$ is an $X_A$-manifold.
 Assume $(ii)$ and let $M$ be a $X_A$-manifold.
 Denote by $\widetilde M \xrightarrow{p} M$ the maximal covering of $M$ possibly branched over $\sing(M)$. For each $\alpha\in A$, we give ourselves a tamifying topology $\T_\alpha$ on $\widetilde X_\alpha$ and a $\T_\alpha$-continuous developing map $\D_\alpha$ given by hypothesis {\it (b)}.   We choose a developing map $\reg(\widetilde M)\xrightarrow{\D} X$.
 
  Let  $(\U,\V,\phi)$ be a  $X_A$-chart with $\V\subset X_\alpha$, let $\widetilde{X_\alpha} \xrightarrow {q} X_{\alpha}$ the maximal branched covering and let $\widehat \U$ be a connected component of $p^{-1}(\U)$. Take any connected component $\widehat \V$ of $q^{-1}(\V)$.
By Lemma \ref{lem:branched_covering_restriction}, $ \widehat \U\xrightarrow{p} \U$ and $\widehat \V\xrightarrow{q}\V$ are branched coverings. Define $\widetilde \V\xrightarrow{r}\widehat \V$ and $\widetilde \U \xrightarrow{\ell} \widehat \U$ 
the natural maximal coverings and lift $\phi$ to an homemorphism $\widetilde \U\xrightarrow{\widetilde \phi}\widetilde \V$.
  Consider the topology $r^*\T_\alpha$ defined in Lemma \ref{lem:spread_top} on $\widetilde \V$, we pull it back to $\widetilde \U$ by $\widetilde \phi$ to get $(r\circ\widetilde \phi)^*\T_\alpha$ and then consider the quotient topology on $\widehat \U$. Using universality of maximal branched covering,  one may check this topology does not depends on $r,\widetilde \phi$ or $\ell$. We call $\T_{\U,\V,\phi}$ the topology constructed this way on $p^{-1}(\U)$.
  We define the topology $\T$ on $\widetilde M$ generated by the topology of $\reg(\widetilde M)$ and the elements of $\T_{\U,\V,\phi}$ on each $ $ for $(\U,\V,\phi)$ going through all $X_A$-charts.


 \begin{figure}[h!]
  $$\xymatrix{
  &&&\widetilde \U \ar@{=}[r]^{\widetilde\phi}\ar[d]^{\ell}&\widetilde \V\ar[d]^{r}&\\
  X&\ar[l]_{\D}\reg(\widetilde M)\ar@{^(->}[r]&\widetilde M\ar[d]^{p}&\widehat \U\ar@{_(->}[l]\ar[d]^{p} &\widehat \V\ar[d]^{q} \ar@{^(->}[r]& \widetilde X_\alpha\ar[d]^{q}\ar[r]^{\D_\alpha}&X\\
  &&M& \U\ar@{_(->}[l]\ar@{=}[r]^{\phi}&\V\ar@{^(->}[r]&X_\alpha
  }$$
 \end{figure}

 Since $X_\alpha$ is locally Galoisian, the branched coverings $\widetilde \V\rightarrow \V$, $\widetilde \V\rightarrow \widehat\V$, $\widetilde \U\rightarrow \U$, $\widetilde \U\rightarrow\widehat \U$  and $\widetilde X_\alpha\rightarrow X_\alpha$ are all Galoisian. Lemma \ref{lem:spread_top} applies to $\widetilde \V \rightarrow \widehat \V \rightarrow \V$ taking $\R= \widehat  \V \cap q^{-1}(\reg(X_\alpha))$. Indeed, $\R$ is strongly portly in $\widehat \V$ by assumption (D) and  since $\widetilde \V := \widetilde \V^{\sing(\V)}$ we have $$\R:=\widehat \V \cap q^{-1}(\reg(X_\alpha))=  q^{-1}(\reg(\V))\subset \ord_r(\widehat \V).$$
 
 \begin{itemize}
   \item The topology $\T_{\U,\V,\phi}$ on $\widehat \U$ does not depend on $(\V,\phi)$
   
   If we consider two $X_\alpha$-charts $(\U,\V_1,\phi_1)$ and $(\U,\V_2,\phi_2)$ then $ \V_1\xrightarrow{\phi_1^{-1}\circ \phi_2} \V_2$ is an isomorphism. By analyticity, $\phi_1^{-1}\circ \phi_2$ is the restriction of an automorphism of $X_\alpha$ which lifts to an automorphism of $\widetilde X_\alpha$. Therefore, there exists $\psi\in \Aut(\widetilde X_\alpha)$ such that $\psi(\widehat \V_1)= \widehat \V_2$. By hypothesis, $\psi$ is $\T_\alpha$-continuous so $(\widehat \V_1,\T_\alpha) \simeq (\widehat \V_2,\T_\alpha) $ then $\T_{\U,\V_1,\phi_1}=\T_{\U,\V_1,\phi_2}$.

  \item The strong topology on $\widetilde M$ is  first countable and locally path connected:
  
  By Lemma \ref{lem:spread_top}, the strong topology on a given $\widetilde \U$ is first countable and locally path connected.  Since $\widetilde \U \xrightarrow{\ell}\widehat \U $ is Galoisian, $\ell$ is open for the strong topologies so the strong topology on $\widehat \U$ is first countable and locally path connected.
  
 
  \item[(A)] The quotient topology induced by $\T_\U$ on $\U$ is by construction the pull back by $\phi$ of the quotient topology induced by $r^*\T_\alpha$ on $\V$. Since $\widetilde \V\xrightarrow{r} \V$ is a branched covering for the strong topology, the quotient topology of $r^*\T_\alpha$ on $\widehat \V$ is $\T_\alpha$ which induces the natural topology on $\V$.

We deduce that the quotient topology induced by the strong topology on $\widetilde M$ is the natural topology $M$.
 



\item[(B)] 
\begin{itemize}
\item $\Gamma(\widehat \V/\V)$ is the set of the restriction of the elements of $\Gamma(\widetilde X_\alpha/X_\alpha)$ stabilizing $\widehat \V$ setwise, these acts strongly continuously by hypothesis (C) on $X_\alpha$. By Lemma \ref{lem:spread_top}, the Galois group $\Gamma(\widehat \V/\V)$ acts strongly continuously  on $\widetilde \V$. 



\item  

Denote 
$\D_\U := \D_\alpha\circ r\circ \widetilde \phi$ which is strongly continuous and note that $\D_{\U|\reg(\widetilde \U)}$ is a $(G,X)$-morphism.
The holonomy of $\widetilde \U\rightarrow \widehat \U$  is trivial  
hence for all $\gamma\in \Gamma(\widetilde \U/\widehat \U)$, $(\D_\U \circ  \gamma)_{|\reg(\widetilde \U)} = \D_{\U|\reg(\widetilde \U)}$.

By the previous point and strong density of $\reg(\widetilde \U)$   we deduce that $$\forall \gamma \in \Gamma(\widetilde \U/ \widehat \U),\quad  \D_\V \circ  \gamma = \D_{\V}.$$



\item Since $\widetilde \U\rightarrow \widehat\U$ is Galoisian, $\Gamma(\widetilde \U/\widehat \U)$ acts transitively on the fibers of $\ell$.
We can thus descend $\D_{\U}$ to an a.e $(G,X)$-morphism  $\widehat \U\xrightarrow{\D_{\U,\V,\phi}} X$ with $\widehat \U$ endowed with the strong topology.

Since $\D_{|\reg(\widehat \U )}=g\D_{\U|\reg(\widehat \U)}$ for some $g\in G$, the developing map $\D_{|\reg(\widehat \U)}$ extends to a strongly continuous map on  $\widehat \U$. Finally, $\D_{|\reg(\widetilde M)}$ extends to an a.e. $(G,X)$-morphism for the strong topology on $\widetilde M$.
\end{itemize}
\item[(C)] Let $\gamma\in \Gamma(\widetilde M/M)$, for every $X_A$-chart $(\U,\V,\phi)$, it sends connected component $\widehat \U$ of $p^{-1}(\U)$ to connected component so the local model of such a $p(\widehat \U)$ is the same as $p\circ \gamma (\widehat \U)$.  The topology $\T_\U$ on $\gamma (\widehat \U)$ can then be constructed choosing the maximal branched covering $\widetilde \U\xrightarrow{\gamma\circ \ell} \gamma(\widehat \U)$. As a consequence, the strong topology on $\gamma(\widehat \U)$ is the image by $\gamma$ of the strong topology on $ \widehat \U$. Therefore, $\gamma^{-1}$ is strongly continuous. We deduce that $\Gamma(\widetilde M/M)$  acts strongly continuously on $\widetilde M$.
   \item[(D)] Recall that we applied Lemma \ref{lem:spread_top} with $\R=q^{-1}(\reg(\V))$ so that
  $$r^{-1}(\R)\subset \ord_{q\circ r}(\widetilde \V)\cap  \reg(\widetilde \V) = \ord_{p\circ \ell\circ \widetilde \phi}(\widetilde \V)\cap  \reg(\widetilde \V)$$ and then its image by $\ell\circ \widetilde \phi$ is a subset of $p^{-1}(\reg(\U))\subset \widehat \U$. Since $r^{-1}(\R)$ is strongly portly  its image by the strongly open surjective map $\ell \circ \widetilde \phi$ is also strongly portly.  
  
  Since $p^{-1}(\reg(\U))\cap \widehat \U$ is weakly open thus strongly open, and since it contains $\ell\circ \widetilde \phi(r^{-1}(\R))$; then $p^{-1}(\reg(\U))\cap \widehat \U$ is strongly portly in $\widehat \U$. Therefore, $p^{-1}(\reg( M))$ is locally strongly portly hence strongly portly.

  \end{itemize}
\qed\end{proof}
\begin{rem}\label{rem:hypothesis_model_spaces_necessary}   Hypothesis (C) for the model spaces is necessary for the proof.  Take $\theta \in \RR_+\setminus \{0,2\pi\}$ and consider $\mass{\theta}$ as a singular model.
We may choose as tamifying topology either the natural one or the tamifying topology $\T$ we defined on $\mass{0}$. Indeed, properties (A),(C) and (D) are purely topological and $\BTZtilde \rightarrow \BTZ$ is isomorphic to $\widetilde \E_\theta^{1,2}\rightarrow \mass{\theta}$ in $\CatSpread$. 
Property (B) is certainly satisfied since $\mass{\theta}$ is tame.
We may thus construct a messed up tamifying topology on $\widetilde \E_\theta^{1,2}$ by taking the topology $\T'$ generated by the natural one and the restriction of $\T$ to $\{t>0\}$. Since  the Galois group acts by 'rotation' around the 'axis' $\{r=0\}$ it preserves $\T'$. However, the automorphism group also contains translations along the $t$-axis which send points from $t<0$ to $t>0$. 

The proof above thus fails at the first step since the induced topology $\T_{\U,\V,\phi}$ for $\sing(\U)\neq\emptyset$ depends on whether $\V\subset \{t>0\}$ or $\V \subset \{t>0\}$ or any other intermediary cases with $(0,0,0)\in \V$.



That being said, the author ignores whether one could remove this hypothesis by showing that from  a given tamifying topology, the coarsest topology that makes the developing map and the automorphisms continuous would still satisfy (A) and (D).


\end{rem}



 Finally, the following remark though trivial is important.
 \begin{rem}\label{rem:quotient_virtually_tame} Let $M$ be a virtually tamely singular $(G,X)$-manifold. If $M$ admits an injective developing map $N\xrightarrow{\D}X$ for some branched covering $N\rightarrow M$ possibly branched above $\sing(M)$ then $M$ is isomorphic to $\Gamma(N/M)\backslash\D(N)$ where 
 \begin{itemize}
  \item $\Gamma(N/M)$ acts on $X$ via the holonomy of $\D$;
  \item $N$ is endowed with the quotient of a tamifying topology $\T$;
  \item $\D(N)$ is endowed with the topology generated by that of $X$ and the images by $\D$ of $\T$.
 \end{itemize} 
 \end{rem}


\section{An application : uniformization of $\BTZ$-manifolds}\label{sec:BTZ_uniform}
In this section we present  an 'application' of the theory developped. Our aim is to present a uniformization result for $\BTZ$-manifolds eg singular $(\isom(\mass{}),\mass{})$-manifolds with singularities modeled on $\BTZ$. 
Those manifolds are not tame but we will show they are virtually tame. 
We can thus define $\Gamma(\widetilde M/M)$-equivariant developing maps on any $\BTZ$-manifold and we wish to apply Remark \ref{rem:quotient_virtually_tame} to some reasonable class of $\BTZ$-manifolds.

We begin by a short introduction to the geometry of Minkowski space is necessary.

\subsection{Geometry of Minkowski space}\label{sec:Minkowski}
We choose $X = \E^{1,2}$, the affine space $\RR^3$ endowed with cartesian coordinates $(t,x,y)$ and a quadratic form $Q(t,x,y)= -t^2+x^2+y^2$ of signature $(1,2)$;
together with  $G=\GG\ltimes \RR^3$ the identity component of its group of affine isometry. 
We will denote $(\E^{1,2},\GG\ltimes \RR^3)$-manifolds by $\E^{1,2}$-manifolds for simplicity sake. 


We choose the point $O=(0,0,0)\in \E^{1,2}$ and call it the origin of $\E^{1,2}$,  we have a natural projection $\GG\rtimes \RR^3 \ltimes \E^{}\xrightarrow{L}  \GG$ we will refer to as {\it linear part}. $\GG$ acts naturally on the paraboloid model of the hyperbolic plane $\H := \{(t,x,y)\in \E^{1,2}~|~ t>0, ~~Q(t,x,y)=-1 \}$ and $\GG \simeq \PSL(2;\RR)$. We say that $\phi\in \GG\ltimes \RR^3$  is parabolic (resp. elliptic, resp. hyperbolic) if the action of its linear part $L(\phi)$ on  $\H$ is parabolic (resp. elliptic, resp. hyperbolic).
A useful characterisation is that $\phi$ is 
\begin{itemize}
 \item parabolic if $\phi$ stabilises set-wise exactly one lightlike line, 
 \item elliptic if $\phi$ stabilises set-wise exactly one timelike line,
 \item hyperbolic if $\phi$ stabilises set-wise two lightlike lines.
\end{itemize}



\subsection{Causality of Minkowski space}\label{sec:mink_caus}


Minkowski space $\E^{1,2}$ is naturally endowed with two partial order relations induced by strict orders defined as follows : for any $p:=(t,x,y), q:=(t',x',y')$  we say that 
$p<q$ (resp. $p\ll q $) if $Q(q-p)\leq 0$ (resp. $Q(q-p)<0$) and $t'>t$. 
A non-trivial vector $u$ is lightlike if $Q(u)=0$, causal if $Q(u)\leq 0$, timelike if $Q(u)<0$ and spacelike if $Q(u)>0$. A line is lightlike (resp. spacelike, resp. timelike) if its direction is lightlike (resp. spacelike resp. timelike).
The relation $<$ is the causal order while $\ll$ is the chronological order. 
A curve is future causal  (resp. chronological) if it is increasing of the order $<$ (resp. $\ll$). 



One can check that the group $G$ acts by non-decreasing homeomorphisms for both orders eg : 
$$ \forall p,q \in \E^{1,2}, \forall \phi\in \GG,\quad 
 p< q \Rightarrow \phi p < \phi q.$$
$$ \forall p,q \in \E^{1,2}, \forall \phi\in \GG,\quad 
 p\ll q \Rightarrow \phi p \ll \phi q.$$

 \subsection{Locally Minkowski manifolds and Mess-Bonsante-Barbot Theorem}
\label{sec:mess}
Every $\E^{1,2}$-manifolds is naturally endowed with two sheafs of pre-order relations inherited from the order relations of $\E^{1,2}$.  
Future causal or chronological curves are curves that are locally increasing for respective pre-order (or equivalently that are increasing in charts).

Properties of these sheafs of pre-order provides a hierarchy of $\E^{1,2}$-manifolds \cite{MR2196783}. 
 Let $M$ be a $\E^{1,2}$-manifold. 
 \begin{itemize}
  \item 
 $M$ is causal if its induces pre-ordre $\leq$ is an order. 
 
 \item $M$ is globally hyperbolic if  there exists a surface $\Sigma\subset M$ such that every future causal curves intersect $\Sigma$ exactly once. Such a surface is called a Cauchy surface of $M$.

 
 \item $M$ is globally hyperbolic Cauchy-compact if it admits a compact Cauchy-surface
 \item $M$ is globally hyperbolic Cauchy-compelte if it admits a smooth Cauchy-surface on which the semi-Riemannian metric of $M$ induces a complete Riemannian metric. 
 \end{itemize}

Following notations of \cite{oneil}, for $M$ a $\E^{1,2}$-manifold and $p\in M$, we denote 
$$J^+(p) := \{q\in M~|~ p\leq q\} \quad J^-(p) := \{q\in M~|~ p\geq q\}$$
$$I^+(p) := \{q\in M~|~ p\ll q\} \quad I^-(p) := \{q\in M~|~ p\gg q\}.$$

A classical result of Geroch \cite{MR0270697,MR2196783} implies that a $\E^{1,2}$-manifolds $M$ is globally hyperbolic if and only if $M$ is causal and if $J^+(p)\cap J^-(q)$ is compact of all $p,q\in M$. Another classical result of Choquet-Bruhat and Geroch \cite{MR0250640,sbierski_geroch} shows that  any globally hyperbolic $\E^{1,2}$-manifolds $M$ embeds naturally into a Cauchy-maximal globally hyperbolic extension. 
The precise definition of Cauchy-maxi\-mality is not relevant for the following except as an hypothesis, we therefore omit to further develop it.  


Following the work of Mess \cite{mess}, Bendetti, Bonsante \cite{MR2499272,MR2170277} and Barbot \cite{barbot_globally_2004}, a uniformization of globally hyperbolic Cauchy-maximal Cauchy-complete $\E^{1,2}$-manifolds is known.  Let $M$ be such a manifold of Cauchy surface $\Sigma$. For simplicity sake, we restrict ourselves to $\Sigma$ of finite type with genus $g$ and $s$ punctures such that $2g-2+s>0$.   
In this case, there exists a representation  $\pi_1(\Sigma)\xrightarrow{~~\rho~~} \GG\ltimes \RR^3$ whose linear part is discrete and faithful as well as a $\rho$-invariant convex domain $\Omega\subset \E^{1,2}$ such that $M\simeq \rho\backslash \Omega$ and $\widetilde M \simeq \Omega$. Furthermore, up to reversing the causal order,  there exists a $\rho$-invariant family $(\Delta_i)_{i\in I}$ of 
lightlike lines such that $$\Omega =  \bigcap_{i\in I} I^+(\Delta_i).$$
Note that for each $i\in I$, the set $I^+(\Delta_i)$ is an open halfspace of $\E^{1,2}$.

\subsection{Some complements on the BTZ model}

Recall that we introduced the developing map
$$\fonctiondeux{\D}{\BTZtilde}{\E^{1,2}}{\left(\tau,\r,\frac{\theta}{2\pi}\right)}{
\begin{pmatrix}t \\ x \\ y\end{pmatrix}=\begin{pmatrix} \tau + \frac{1}{2}\r\theta^2 \\ \tau + \frac{1}{2}\r\theta^2 -\r \\ -\r\theta  \end{pmatrix}
}{(\tau,0,\cdot)}{\begin{pmatrix}\tau \\ \tau \\ 0\end{pmatrix}}$$

continuous for the topology $\T_{BTZ}$ generated  by the natural topology of $\BTZtilde$ and open sets of the form

$$\U_\lambda = \left\{(\tau,\r,\theta)\in \BTZtilde ~|~ \r\theta^2< \lambda  \right\},\quad \lambda \in \RR_+^*.$$


The image of this map is the  causal future of a lightlike line $\Delta$ ie 
$\D(\BTZtilde)=J^+(\Delta)$. More precisely, $\Delta$ is the line directed by $
\overrightarrow u = (1,1,0)$ through the origin of Minkowski space, $\D(\BTZtilde)$ is then the union of $\Delta$ with the open half space of $\mass{}$ above the plane directed by $\overrightarrow u^\perp$ through the origin.
The Galois group $\Gamma:=\Gamma(\BTZtilde/\BTZ)$ is isomorphic to $\ZZ$ and stabilizes point-wise the singular locus; taking $\gamma$ a generator of $\Gamma(\BTZtilde/\BTZ)$, the holonomy $\rho$ sends $\gamma$ to $\phi:=\rho(\gamma)$ which thus stabilizes point-wise the image of $\sing(\mass{0,\infty})$ ie $\Delta$. Since $\Delta$ is lightlike, $\phi$ is parabolic.

We did not quite prove that $\BTZ$ is virtually tame, let us correct this.
\begin{prop}\label{prop:BTZ_virt_tame}The $\BTZ$ model space is virtually tame.
 
\end{prop}
\begin{proof} 
 From the motivational example given in section \ref{sec:mot_BTZ}, the only missing property is that the induced quotient topology on $\BTZ$ is indeed the natural one. 
 
 For any $\lambda>0$,  the image of $\U_\lambda$ in $\BTZ$
contains $\left\{(\tau,\r,\theta)~|~ \r < \frac{\lambda}{4\pi^2} \right\}$ and is thus open. Therefore, the quotient topology on $\BTZ$ is the natural topology.
\qed\end{proof}

\begin{prop} \label{prop:BTZ_manifold_virt_tame} All $\BTZ$-manifolds are virtually tame.
\end{prop}
\begin{proof}
 By Proposition \ref{prop:BTZ_virt_tame}, the model space $\BTZ$ is virtually tame. From the motivational example given in section \ref{sec:mot_BTZ} we have that $\Aut(\BTZtilde)$ acts strongly continuously. Topologically, $\BTZ$ is homeomorphic to $\RR^{3}$ and $\sing(\BTZ)$ is a line. Therefore $\BTZ$ is locally Galoisian. Theorem \ref{theo:local_virt_tame} thus applies.
\qed\end{proof}


We now identify the right topology to put on $\BTZtilde$. To begin with  the holonomy is faithful and that $\D$ is injective, therefore the developing map induces a bijective map 
$$\overline \D :  \BTZ \xrightarrow{} \Gamma\backslash J^+(\Delta).$$
using the tamifying topology of $\BTZ$ we deduce that 
 $\overline \D$ is continuous.


However, the map $\overline \D$ is \textbf{not open} thus not a homeomorphism if the righthand side is endowed with the quotient topology of the usual topology of $J^+(\Delta)$. Actually, $\Gamma\backslash J^+(\Delta)$ is not even Hausdorff. The situation is very similar to that of cuspidal singularities viewed as point on the unit circle in the closure of the disc model of $\H$ embedded into $\RR^2$. 
To see that, consider the sequence $(\tau_n,\r_n,\theta_n)_{n\in \N}$ in $\BTZtilde$ with $\tau_n=-1$, $\r_n=1/n$ and $\theta_n=\sqrt{2n}$. Observe that $\D(\tau_n,\r_n,\theta_n)\xrightarrow{n\rightarrow +\infty} (0,0,0)$ but that $(\tau_n,\r_n,\theta_n) \xrightarrow{n\rightarrow +\infty} (-1,0,0)$. Furthermore, choose some sequence $(k_n)_{n\in \N}$ such that $\theta_n-k_n$ is bounded so that $\D(\tau_n,\r_n,\theta_n-k_n)\xrightarrow{n\rightarrow +\infty} (-1,0,0)$. Since $(\tau_n,\r_n,\theta_n)$ and $(\tau_n,\r_n,\theta_n-k_n)$ are $\Gamma$-conjugate  for all all $n\in \N$ they induce the same sequence in $\BTZ$. 
Its image by $\overline \D$ thus has two distincts limits. 

To force this map to be a homeomorphism, one has to add to the topology of $J^+(\Delta)$ the image by $\D$ of neighborhood basis of points in $\{\r=0\}\subset \BTZtilde$ which leads to the following definition.

\begin{defi}[BTZ topology]\label{defi:BTZtop} Let $\Delta$ be a lightlike line in $\mass{}$.
The BTZ topology on $J^+(\Delta)$ is the topology generated by the one induced by the natural topology of $\mass{}$ and open subsets of the form $I^+(p)\cup]p,+\infty[$  for  $p\in \Delta$ where $]p,+\infty[$ denote the relatively open future half-ray from $p$.
\end{defi}

With this definition we have the following.
\begin{prop} \label{prop:quotient_chart}
If $J^+(\Delta)$ is endowed with the BTZ topology,
then the map $\overline \D : \mass{0}\rightarrow \Gamma\backslash J^+(\Delta)$  is an isomorphism of singular $\E^{1,2}$-manifold.
\end{prop}


This last remark will be useful.
\begin{rem}\label{rem:tamifying_topology_generated_open} The tamifying topology $\T_{BTZ}$ is generated by $\T_{BTZ}$-open diamond eg the sets of the form $Int_{\T_{BTZ}}\left(J^{-}(p)\cap J^{+}(q)\right)$ for $p,q\in \BTZtilde$.
\end{rem}






\subsection{The maximal branched covering of $\mass{0}$-manifolds}

Consider $M$ a $\mass{0}$-manifold, ie a singular $\mass{}$-manifold whose singular locus is locally modelled on $\mass{0}$. As notice above, the $\BTZ$ is locally Galoisian so $\widetilde M\rightarrow M$ is Galoisian.
And by Proposition \ref{prop:BTZ_manifold_virt_tame}, $M$ is virtually tame 

We may add that the singular locus being a 1-submanifold of non trivial holonomy, $\widetilde M$ is exactly branched over $\sing(M)$.
We summarize this as follow.
\begin{prop}\label{prop:manifold_admits_galoisian_maximal} Let $M$ be a $\mass{0}$-manifold, then $M$ admits a Galoisian maximal covering exactly branched above $\sing(M)$ and admits a developing map continuous for the tamifying topology $\T_{BTZ}$.
Furthermore, $\widetilde M$ is a $\BTZtilde$-manifold.
\end{prop}
For simplicity sake, in the rest of the section, ``Developing map'' of some $\BTZ$-manifold $M$ will always refer to the one for the $\T_{BTZ}$-topology.


The goal of this section is to show, similarly to Cauchy-complete Cauchy-maximal $\mass{}$-manifolds, that Cauchy-complete Cauchy-maximal $\mass{0}$-manifolds can be realized as quotients of convex domains of Minkowski space.

Causal notions introduced in the section \ref{sec:mess} above extend naturally to $\BTZtilde$-manifolds and $\BTZ$-manifolds. Although, most of the proofs of elementary causal properties for Lorentzian manifold extend {\it mutatis mutandis} to such manifolds, a cautious reader will have to re-prove most properties from scratch. A proper framework for general topological spacetimes applying our manifolds of interests can be found in \cite{thesis} and the author aims at publishing them. 
\begin{defi}\label{def:manifolds_future_distinguishing_map} A $\BTZtilde$-manifolds $M$ is future distinguishing if the map $x\mapsto I^+(x)$ defined on $M$ is injective.
\end{defi}
\begin{defi}\label{def:singular_manifold_endowed_causal}
 A singular $\mass{}$-manifold $M$ endowed with some causal pre-order $\leq$ extending the one on $\reg(M)$ is causally connected if 
for all $x,y\in M$ such that  $x\leq y$, there exists a future causal path from $x$ to $y$.
\end{defi}
The proof of the following two Lemmas are omitted as they fall in the elementary causal property mentionned above and proving them (though logically necessary) is not particularly useful for our current exposition.

\begin{lem}\label{lem:globally_hyperbolic_manifolds_manifolds} Globally hyperbolic  $\BTZ$-manifolds and $\BTZtilde$-manifolds are future distinguishing.
\end{lem}


\begin{lem}\label{lem:manifolds_manifolds_causally_connected} $\BTZ$-manifolds and $\BTZtilde$-manifolds are causally connected.
\end{lem}



\begin{lem}\label{lem:increasing_map} Let $M$ be a  $\BTZ$-manifold and let $\D:\widetilde M\rightarrow \mass{}$ be a developing map.
Then $\D$ is  increasing on $M$ eg
$$\forall x,y\in M, \quad x<y \Rightarrow \D(x)<\D(y).$$
\end{lem}
\begin{proof} 
 The restriction of such a map to a chart neighborhood is the restriction of a developing map of the local model space. Since the developing maps of each model spaces are  increasing, $\D$ is locally increasing. Furthermore, since $M$ is connected, it is causally connected therefore $\D$ is increasing.
\qed\end{proof}


\begin{defi}\label{def:manifold_future_complete_all}
 A $\BTZ$-manifold $M$ is future complete if for all $x\in \widetilde M$ 
 $$ \D(J^+(x)) = J^+(\D(x)) \text{ and }  \D(I^+(x)) = I^+(\D(x)).$$
\end{defi}

\begin{prop}\label{prop:globally_hyperbolic_cauchy_complete_manifolds} Globally hyperbolic Cauchy-complete $\BTZ$-manifolds are future complete.
\end{prop}
\begin{proof}
From Theorem 3 of \cite{brunswic2021cauchy} the regular par of $M$ is globally hyperbolic and Cauchy-complete. Barbot uniformization Theorem  \cite{barbot_globally_2004} then implies future completeness of the regular part. 
Finally Lemma 2.4 of \cite{brunswic2021cauchy} implies that the singular locus contains full future BTZ rays.
\end{proof}


\begin{defi}\label{def:manifolds_future_distinguishing_map_2} A $\BTZtilde$-manifolds $M$ is future distinguishing if the map $x\mapsto I^+(x)$ defined on $M$ is injective.
\end{defi}

\begin{lem}\label{lem:globally_hyperbolic_manifolds_maximal} Let $M$ be a Globally hyperbolic $\BTZ$-manifolds. its maximal branched covering $\widetilde M$ is future distinguishing.
\end{lem}
Although, one has to re-prove most causal properties from scratch, the proof of this Lemma it not different from the proof of the similar statement for Lorentzian manifolds. 

\begin{thm}\label{theo:globally_hyperbolic_cauchy_complete_cauchy_maximal} Let $M$ be a globally hyperbolic Cauchy-complete Cauchy-maximal $\mass{0}$-manifold, let $\widetilde M$ its maximal covering branched over $\sing(M)$, of  developing map and holonomy  $(\D,\rho)$. We note $\Gamma:=\Gamma(\widetilde M/M)\simeq\pi_1(\reg(M))$ its Galois group.  Then, up to time reversal of $M$:
\begin{itemize}
 \item $\D$ is injective
 \item $\rho$ is faithful, discrete and torsion-free; 
 \item there exists a $\Gamma$-invariant family of relatively open future complete lightlike  rays $(\Delta_i)_{i\in I}$ and a $\Gamma$-invariant family of lightlike planes $(\Pi_j)_{j\in J}$ such that 
 $$\D(\widetilde M) = \bigcap_{i\in I} J^+(\Delta_i) \cup \bigcap_{j\in I} I^+(\Pi_j) $$
 
 \item endowing $\D(\widetilde M)$ with the BTZ topology in the neighborhood of the lightlike rays $(\Delta_i)_{i\in I}$, the developing map induces an isomorphism
 $$ M\xrightarrow{~~\overline \D~~} \rho\backslash\D(\widetilde M
 ) $$
\end{itemize}


\end{thm}
\begin{proof}
 From Theorem 3 of \cite{brunswic2021cauchy}, the regular part of $M$ is globally hyperbolic, Cauchy-maximal and Cauchy-complete; therefore, by general result of Barbot \cite{barbot_globally_2004}, the restriction $\D_{|\reg(\widetilde M)}$ is injective, $\rho$ is faithful, discrete and without torsion and $$ \Omega:=\D(\reg(\widetilde M)) =\bigcap_{j\in I} I^+(\Pi_j)$$ for some $\Gamma$-invariant family $(\Pi_j)_{j\in J}$ of lightlike planes. Furthermore, $\Gamma$ acts totally discontinuously and freely on $\Omega$. Notice that $\Omega$ is future complete ie for all $x\in \Omega, J^+(x)\subset \Omega$.
 
 Note that from Lemma \ref{lem:increasing_map}, $\D$ is increasing.  
 Let $x\in \sing(\widetilde M)$, its stabilizer $\Gamma_x$ is isomorphic to $\ZZ$ and for $\gamma\in \Gamma_x\setminus\{1\}$, $\rho(\gamma)$ is parabolic. Therefore, the set of fix points of $\Gamma_x$ in $\mass{}$ is a lightlike line $\Delta$. The restriction of $\D$ to the connected component  $S_x$ of $\sing(\widetilde M)$ containing $x$ is increasing with image $\Delta_x:=\D(S_x)\subset \Delta$. 
 Since $S_x$  is a totally ordered relatively open future half line, so is its image $\D(\Delta_x)\subset\Delta$.
 
 Since $\Gamma$ does not act freely on $S_x$, it does not acts freely  via $\rho$ on $\Delta$, therefore $\Delta$ does not intersects $\Omega$. Since $\D$ is $\T_{BTZ}$-continuous
 $\D(\Delta_x)$ lie in the closure of $\Omega$ hence $\D(\Delta_x)\subset \partial \Omega$. 
 Finally, since $J^+(\D(\Delta_x)) = I^+(\D(\Delta_x))\cup \D(\Delta_x) $ and since $M$ is future complete,
 $$\Omega\supset \D(I^+(S_x)) =I^+(\D(S_x))$$ so $\D(\widetilde M)$ has the wanted form. 
 
 Let $x,y\in \widetilde M$ such that $\D(x)=\D(y)$. Since $M$ is future complete, we have $$\D(I^+(x))=I^+(\D(x))=I^+(\D(y))=\D(I^+(y))\subset \Omega$$
 and since $I^+(x)$ and $I^+(y)$ are subsets of $\reg(M)$ and $\D_{|\reg(M)}^{|\Omega}$ is a bijection we deduce that $I^+(x)=I^+(y)$. Since $\widetilde M$ is globally hyperbolic, it is future distinguishing and $x=y$. Finally, $\D:\widetilde M \rightarrow \mass{}$ is injective.
 
 The map $\D_{|\reg(\widetilde M)}$ is open since it is a local homeomorphism. The image of a neighborhood basis of a singular point $x\in \sing(\widetilde M)$ is obtains by considering sets of the form $I_{p,q}:=\mathrm{Int}(J^+(p)\cap J^-(q))$ with $p\in J^-(x)$ and $q\in I^+(x)$. By taking $p,q$ sufficiently close to $x$, such a neighborhood can be chosen in a causally convex in $M$ chart neighborhood $\U$ of $x$.  This way $\mathrm{Int}(J^+(p)\cap J^-(q)) = \mathrm{Int}(J_\U^+(p)\cap J_\U^-(q))$ and the image by $\D$ of such a domain is the exactly $(I^+(\D(p))\cup \D(]p,+\infty[)\cap I^-(\D(q))$ with $]p,+\infty[$ the future half BTZ ray from $p$ in $\widetilde M$. Therefore, the image of such $\D(I_{p,q})$ together with the topology induced by $\mass{}$ generates the BTZ topology on $\D(\widetilde M)$.
 
 For the BTZ topology on $\D(\widetilde M)$, the map $\D : \widetilde M \rightarrow \D(\widetilde M)$ is then a $\Gamma$-invariant homeomorphism and induces a isomorphism of singular $\mass{}$-manifolds $\Gamma\backslash\widetilde M\simeq \rho\backslash\D(\widetilde M) $. Since $\widetilde M \rightarrow M$ is Galoisian, $M\simeq \Gamma\backslash\widetilde M$ and 
$M\simeq\rho\backslash\D(\widetilde M)$.
\qed\end{proof}


\bibliographystyle{alpha}   
\bibliography{note}	

\end{document}


 
% \begin{prop} Let $M$ be a singular $(G,X)$-manifold such that $\reg(M)$ is semi-locally simply connected and let $x_0\in \reg(M)$. 
% The following are equivalent:
% 
% \begin{enumerate}[(i)]
%  \item the developing map $\D:\Omega(x,\reg(M),\reg(M)) \rightarrow X$ extends continuously to a map $\Omega(x,M,\reg(M))\rightarrow X$;
%  \item $M$ is tamely singular.
% \end{enumerate}
% \end{prop}
% 
% \subsection{Tame singularities in length spaces}
% Recall that a length space is a metric space $(X,d)$ such that the distance between $a,b\in X$ is the infimum of the length of curves from $a$ to $b$. 
% We take the convention that distances can take value $+\infty$. 
% The main object of this section are analysed length structure in the following sense.
% \begin{defi} An analytical length structure is a triple $(G,X,d)$ where $(G,X)$ is an analytical structure, $d$ a length space structure on $X$ and $G$ acts on $X$ by isometries.
% \end{defi}
% We will frequently omit to name the metric and simply refer to an analytical length structure $(G,X)$ instead of $(G,X,d)$.
% Let us begin by a simple remark whose proof is left to the reader.
% \begin{rem}
%  Let $(G,X)$ an analytical length structure and let $M$ be a singular $(G,X)$-manifold $M$.
%  For $\gamma \in \Omega(x,y,\reg(M))$, define the length $Len_M(\gamma)$ as the infinimum 
%  $$Len_M(\gamma) = \inf \sum_{i=1}^p Len_X(\varphi_i\circ\gamma_{|]t_{i-1},t_i[}) $$
%  where $Len_X$ is the length in $X$ in the usual sense and the infimum is taken over all tuple  of charts $(\U_i,\V_i,\varphi_i)_{i\in \lsem 1,p\rsem }$ and all  subdivision $t_0=0<\cdots<t_p=1$ such that $\forall i\in \lsem 1,p\rsem, \gamma(]t_{i-1},t_i[) \subset \U_i$.
%  And define $$\forall x,y\in M,\quad d(x,y) = \inf_{\gamma\in \Omega(x,y,\reg(M))}Len_M(\gamma).$$
% Then $(M,d)$ is a length space.
% \end{rem}
% In the whole section, we will refer to this distance as the {\it natural length space structure}.  For $M$ a singular $(G,X)$-manifold with topology $\mathcal T$, if no distance or topology is specified, we shall implicitely endow $M$ with its {\it given topology} $\mathcal T$.
% As the following example shows, there is no guarantee that the natural length structure is compatible with the given topology. 
% 
% \begin{ex}
%  Consider again $M=\mathbb S^2$ endowed with the singular $\E^2$-structure defined in Example  \ref{ex:eucl_sphere} and denote $\T$ its given topology. We see that $(\E^2, \SO_2(\RR)\times \RR^2)$ is an analytical length structure and that for the natural length structure $d$ on $M$, the space $(M,d)$ has two connected  components : $\reg(M)$ and $\sing(M)$. In particular, $(M,d)$ is not homeomorphic to $(M,\T)$
% \end{ex}
% 
% In this section, we caracterise tame singularities by the concordance between the length topology and the given topology. 
% 
% 
% We'll need the following classical Lemma of which we provide a proof since we could not find a reference to point to.
% \begin{lem} Let $(M,d)$ be a semi-locally simply connected  length space. If $d$ takes finite values, then every homotopy class with fixed ends contains a finite length path.
% \end{lem}
% \begin{proof}
% Consider a path $\gamma: [0,1]\rightarrow \reg(M)$; by compactness, we can choose $t_0=0,\cdots,t_n=1$ such that for all $k\in \lsem 1,n\rsem$, $\gamma([t_{k-1},t_k])\subset \U_k$ with $\U_k$ simply connected in $M$. For each $k\in \lsem 1,n\rsem $, again by a compactness argument, we can find $\varepsilon>0$ and $t_{k,0}=t_{k-1},\cdots, t_{k,\ell_k}=t_{k}$ such that $\gamma([t_{k,i-1},t_{k,i}])\subset B(\gamma(t_{k,i}),\varepsilon) \subset \U_k$.
% Since $d(\gamma(t_{k,i-1}),\gamma(t_{k,i}))<\varepsilon$, there exists a path $\omega_{k,i}:\gamma(t_{k,i-1})\rightsquigarrow \gamma(t_{k,i})$ of length stricly lesser than $\varepsilon$. In particular, $\omega_{k,i}$ is included in $B(\gamma(t_{k,i}),\varepsilon)$.  
% For  $k\in \lsem 1,n\rsem$, consider the path $\omega_k = \omega_{k,\ell_k}*\cdots*\omega_{k,1}$; 
% $\omega$ has same end points as $\gamma_{|[t_{k-1},t_k]}$ and takes its values in $\U_k$. Therefore, $\gamma_{|[t_{k-1},t_k]}$ is homotopic to $\omega_k$ in $M$. 
% We deduce that $\gamma$ is homotopic to $\omega_n*\cdots *\omega_1$ which has finite length.
% \end{proof}
% 
% 
% 
% 
% 
% \begin{lem} 
%  Let $(G,X)$ be an analytical length structure and let $M\xrightarrow{f} N$ an a.e. $(G,X)$-morphism. 
%  Then, $f$ preserves the length of curves in $\reg(f)$. 
% \end{lem}
% \begin{proof}
%  Recall that $\reg(f)$ is the biggest portly subset of $\reg(M)$ such that $f(\reg(f))\subset \reg(N)$ and $f_{|\reg(f)}^{|\reg(N)}$ is a $(G,X)$-morphism.  If $\gamma$ is a curve in $\reg(f)$ included in a connected chart $(\U,\V,\varphi)$ of $M$ such that $f(\U)$ is included in a connected chart $(\U',\V',\varphi')$  of $N$ then there exists $g\in G$ such that for all $x\in \U$ we have $\varphi'(f(x))= g \varphi(x)$. Therefore,
%  $$L_N(f\circ \gamma)):=L_X(\varphi'\circ f\circ \gamma)=L_X(g\varphi\circ \gamma) = L_X(\varphi\circ \gamma)=:L_M(\gamma).$$ 
%  If $\gamma\in \Omega(x,y,\reg(f))$ for some $x,y\in M$,then subdivisions $t_0=0<\cdots<t_p=1$ of $\gamma$ with $\max_{i}|t_i-t_{i-1}|$ small enough satisfy the previous assumptions. Therefore, $Len_M(\gamma) =Len_N(f\circ \gamma)$. 
% 
% \end{proof}
% \begin{cor} \label{cor:lip}With the same notations as before,
%  if $\reg(f)=\reg(M)$ then $f$ is 1-lipchitz on $\reg(M)$.
% \end{cor}
% \begin{proof} Let $x,y\in \reg(M)$,
%   $$d(f(x),f(y))  \leq \inf_{\gamma \in \Omega(x,y,\reg(M))} L_N(f\circ \gamma) =  \inf_{\gamma \in \Omega(x,y,\reg(M))} L_M( \gamma) = d(x,y)$$
% \end{proof}
% 
% \begin{prop}
%  Let $(G,X)$ be an analytical length structure with $X$ complete, let $M$ be a connected singular $(G,X)$-manifold with $\reg(M)$ semi-locally simply connected in $M$. Denote by $\mathcal T$ the topology of $M$ and by $d$ the length structure on $M$ induced by $X$.
%  Then the following are equivalent: 
%  \begin{enumerate}[$(i)$]
%   \item $M$ is tamely singular;
%   \item $(M,d)$ is homeomorphic to $(M,\mathcal T)$.
%  \end{enumerate}
% \end{prop}
% \begin{proof} For simplicity sake, we denote $\widetilde M^{\sing(M)}$ by $\widetilde M$, we give ourselves a couple $(\D,\rho)$  developement/holonomy of $\reg(M)$.  Recall that every homotopy classes in $\reg(M)$ with fixed endpoint contains a finite length path.
% 
%  \begin{itemize}
%   \item $(ii)\Rightarrow (i)$. The map $\D$ is a $(G,X)$-morphism $\reg(\widetilde M) \rightarrow X$. Therefore, $\D:\reg(M)\rightarrow X$ is 1-Lipchitz by Corollary \ref{cor:lip}. Since $(M,d)$ is homeomorphic to $(M,\T)$, $\reg(M)$ is dense in $M$ for the distance $d$.  Finally, $X$ is complete.
%   Therefore, $\D$ extends continuously to $M$.
%   \item $(i)\Rightarrow (ii)$
%   
%   
%  \end{itemize}
% \end{proof}



% \subsection{Tame Geometrical Completion of a singular $(G,X)$-manifold}
% 
% In the list of construction given in section \ref{sec:construction_example} only has been left out until now: completions. These completions have many avatar: cuspidal ends of complete hyperbolic surfaces, BTZ-extension \cite{brunswic2021cauchy}, ... \edl{TO BE COMPLETED}.
% 
% The goal of this section is the construction of a completion of a $(G,X)$-manifold that adds singular points to obtain a tamely singular $(G,X)$-manifold.
% 
% \subsubsection{Simultaneous Spread Completion with respect to a function}
% 
% The first step is the extension of the following diagram 
% $$\xymatrix{
% \widetilde M \ar[rr]^{\D}\ar@{-->}[dr]^{\iota} \ar[dd]^{p}&& X \ar@{=}[dr]\\
% & \widetilde{\overline M} \ar@{-->}[dd] \ar@{-->}[rr]^{\overline \D}&& X
% \\
% M\ar@{-->}[dr]& \\
% & \overline M
% }$$
% ie starting from some regular $(G,X)$-manifold $M$ and its developing map $\widetilde M \xrightarrow{\D} X$ we aim at constructing a complete spread $\widetilde{\overline M}\rightarrow \overline M$ together with a spread embedding so that $\D$ lifts to $\widetilde{\overline M}$.
% 
% 
% \begin{thm} Let $\widetilde M \xrightarrow{p} M$ be a universal covering together with a local homeomorphism $\widetilde{M} \xrightarrow{\D} X$. 
% Assume that these topological spaces are first countable and  $\sigma$-compact.
% 
% Then there exists a branched covering $\widetilde{\overline{M}}\xrightarrow{p} \overline M$ extending $\widetilde M \xrightarrow{p} M$ such that $\D$ extends continuously to $\widetilde{\overline{M}}$
% \end{thm}
% 
% We thus give ourselve such a $\widetilde M \xrightarrow{p} M$ with $\widetilde{M}\xrightarrow{\D} X$.  We give ouselves an increasing sequence $(K_n)_{n\in \N}$ of compact subsets of $M$ whose union is the whole $M$.
% \begin{defi}
% Consider $\mathcal E$ the set of sequences of couple $(\U_n,\widehat \U_n)_{n\in \N}$  such that 
% \begin{enumerate}[(a)]
% \item for all $n\in \N$, $\U_n\subset M$ is connected open
% \item for all $n\in \N,$ $\widehat \U_n$ is  a connected component of $p^{-1}(\U_n)$
%  \item $\forall \in \N, K_n \cap \overline\U_n =\emptyset$
%  \item $\mathrm{Card}\left(\bigcap_{n\in \N}\overline {\D(\widehat \U_n)}\right) = 1$
% \end{enumerate}
% \end{defi}
% \begin{defi}
%  Define on $\mathcal E$ the relation 
%  $(\U_n,\widehat \U_n)_{n\in \N}\leq (\V_n,\widehat \V_n)_{n\in \N}$
% if 
% $$\forall n\in \N,\exists m\in \N,\quad  \U_m \subset  \V_n~~\text{and}~~\widehat \U_m \subset \widehat \V_n.  $$
%  \end{defi}
% 
% \begin{rem}\label{rem:preorder} The relation $\leq$ is a preorder relation thus induces an order relation of $\mathcal E/\sim$ with $\sim$ the equivalence relation generated by $\leq$. 
% 
% We denote by $[\boldsymbol\U]$ the equivalence class of $\boldsymbol \U = (\U_n,\widehat \U_n)_{n\in \N} \in \mathcal E$.
% \end{rem}
% \begin{defi} Let $\boldsymbol \U,\boldsymbol\V \in \mathcal E$, define 
% the union: $$ \boldsymbol \U \cup \boldsymbol\V := (\U_n\cup \V_n,\widehat \W_n)_{n\in \N}$$
% with $\widehat\W_n$ the connected component of $p^{-1}(\U_n\cup\V_n)$ containing $\widehat \U_n\cup \widehat \V_n$. 
% \end{defi}

% \begin{rem} For any familly $(\boldsymbol \U^{(i)})_{i\in I} $ and any $\boldsymbol \V$ in $\mathcal E$. 
% we have the following:
% \begin{enumerate}[(a)]
%  \item if $(\boldsymbol \W^{i})_{i\in I}\in \mathcal E$ is such that $\forall i\in I, \boldsymbol \W^{(i)} \sim \boldsymbol\U^{(i)} $
%  then 
%  $$ \bigcup_{i\in I}\boldsymbol\W^{(i)} \sim \bigcup_{i\in I} \boldsymbol\U^{(i)};$$
%  \item if for all $i\in I, \boldsymbol\U^{(i)}\leq \boldsymbol\V$
%  then $\bigcup_{i\in I} \boldsymbol\U^{(i)} \leq \V$
% If for all $i\in I, \boldsymbol\U^{(i)}\geq \boldsymbol\V$
%  then $\bigcup_{i\in I} \boldsymbol\U^{(i)} \geq \V$
%  \end{enumerate}
% \end{rem}
\subsubsection{Weak developing map}
As the example described in the section \ref{sec:Lorentzian} below will show, some interesting singularities are not tame. We thus extend the notion of developing map to {\it weak developing map}.  

In this section we add the topological hypothesis that \textbf{topological spaces considered are normal} (the model space $X$ as well as all singular $(G,X)$-manifolds $M,N,...$)  ie that for every point $x$ and every neighborhood $\U$ of $x$ there exists a neighborhood $\V\subset \U$ of $x$ whose closure $\overline \V \subset \U$.
The reader can check the following Lemma.
\begin{lem}\label{lem:spread_hausdorff_locally_connected} Let $N\xrightarrow{p}M$ be a spread with $N,M$ Hausdorff and locally connected. 
If $M$ is normal, then so is $N$.
\end{lem}


\begin{defi}[Weak a.e. $(G,X)$-morphism]\label{def:weak_morphism}
Let $M,N$ be singular $(G,X)$-manifolds, a weak a.e. $(G,X)$-morphism  is a possibly discontinuous map 
$N\xrightarrow{f} X$ such that 
\begin{enumerate}[(a)]
 \item there exists a portly subset $\U\subset M$ such that $f_{|\U}$ is an a.e $
 (G,X)$-morphism;
 \item for any $\U$ such as in $(a)$ we have 
 $$\forall x\in M, f(x) \in \bigcap_{\V\ni x} \overline {f(\V\cap \U)}$$
 where the intersection runs over the open neighborhoods $\V$ of $x$. 
\end{enumerate}

Two such weak a.e $(G,X)$-morphisms $N\xrightarrow{f_1,f_2}M$ are equivalent if they are equal on a portly subset.
\end{defi}



\begin{lem}\label{lem:singular_manifolds_weak_morphism}
 Let $M,N$ be singular $(G,X)$-manifolds and $M\xrightarrow{f} N$ a weak a.e. $(G,X)$-morphism.

 Point $(b)$ of definition above can be replaced by 
 \begin{enumerate}
  \item[(b')] There exists a $\U$ such as in $(a)$ such that 
 $$\forall x\in M, f(x) \in \bigcap_{\V\ni x} \overline {f(\V\cap \U)}$$
 where the intersection runs over the open neighborhoods $\V$ of $x$. 
 \end{enumerate}

\end{lem}
\begin{proof} Assume $(a)$ and take $\U_1,\U_2$ as in $(a)$. 
We notice that
$$\forall \V\subset M \text{ open}, \quad \overline {f(\V\cap \U_1)} =   \overline {f(\V\cap \U_2)}$$
which is a direct consequence of continuity of $f$ on $\U_1\cup \U_2$ and density of both $\U_1$ and $\U_2$.
The result follows.

Let $x\in M$ and let $y\in \overline {f(\V\cap \U_1)}$ and let $\W$ be an open neighborhood of $y$. Consider some $x'\in \V\cap \U_1$ such that $f(x')\in \W \cap f(\V\cap \U_1)\neq \emptyset$. Since $f$ is continuous at $x'$ the set $f^{-1}(\W)$ is a neighborhood of $x'$. Then $f^{-1}(\W)\cap \V$ is a neighborhood of  $x'$ and by density of $\U_2$ we have $f^{-1}(\W)\cap\V\cap\U_2 \neq \emptyset$, hence $\W\cap f(\V\cap\U_2) \neq \emptyset$ and $y\in \overline {f(\V\cap \U_2)}$. By symmetry, the result follows.

\end{proof}



\begin{lem}\label{lem:Tamecont}
Let $M,N$ be singular $(G,X)$-manifolds and $M\xrightarrow{f} N$ a weak a.e. $(G,X)$-morphism.

There exists a unique maximal domain $\Tame(f)$ containing a portly subset and such that $f_{|\Tame(f)}$ is an a.e. $(G,X)$-morphism.
Furthermore, for any subset $\U\subset\Tame(f)$ portly in $M$:  $$\Tame(f) = \{x\in M ~|~ f \text{ continuous at } x\} \subset \left\{x\in M ~\left|~ \mathrm{card}\bigcap_{ \V \ni x} \overline{f(\V \cap \U)}=1 \right. \right\}$$
where the intersection runs through the neighborhoods of $x$.
\end{lem}
\begin{proof}



Any domain $M'$ containing a portly subset is connected by 
Lemma \ref{lem:path_covering} and naturally endowed with an a.e. $(G,X)$-structure induced by the one of $M$. Furthermore, $f_{|M'}$ is an a.e. $(G,X)$-morphism if and only if $f_{|M'}$ is continuous.
Define $\Tame(f):= \{x\in M ~|~ f \text{ continuous at } x\}$, by property $(a)$ it contains a portly subset, moreover $f_{|\Tame(f)}$ is continuous thus an a.e. $(G,X)$-morphism.
Furthermore, for any $M'$ as above, $f_{|M'\cup\Tame(f)}$ is still continuous and contains a portly subset. 
Therefore, we need to show that $\Tame(f)$ is maximal. We give ourselves some $\U_0 \subset \Tame(f)$ portly in $M$.

Let $x\in M$ assume that $f_{|\Tame(f)\cup \{x\}}$ is continuous. Take some open $\V\ni f(x)$. Since $N$ is normal, we take some open $\V_0\ni f(x)$ such that $\overline \V_0 \subset \V$. Since $f_{|\Tame(f)\cup\{x\}}$ is continuous, for $\U\ni x $ open small enough,  $f(\U\cap \Tame(f))\subset \V_0$. By  hypothesis $(b)$ we have 
$$f(\U) \subset \overline{f(\U\cap\Tame(f))} \subset \overline \V_0 \subset \V.$$
We deduce that $f$ is continuous at $x$ so $x\in \Tame(f)$.
Assume $f$ is continuous at $x\in M$ and take $\U$ portly subset of $\Tame(f)$, since $N$ is normal Hausdorff 
$\bigcap_{\V\ni x}\overline{f(\V\cap \U)}$ contains exactly one element. 

\end{proof}
\begin{defi}[Tame domain]\label{def:tame_domain} Let $M,N$ be singular $(G,X)$-manifolds and $M\xrightarrow{f} N$ a weak a.e. $(G,X)$-morphism. The tame domain of $f$ is the set $\Tame(f)$ given by Lemma \ref{lem:Tamecont}.
\end{defi}

\begin{lem}\label{lem:Tame_equ} If $f$ is a weak a.e. $(G,X)$-morphism then $\Tame(f)$ only depends on its equivalence class.
\end{lem}
\begin{proof}
 Consider some $M\xrightarrow{g} M$ equivalent to $f$, some portly subset $\U\subset M$ such that $f_{|\U}=g_{|\U}$ are continuous.  Consider $x\in M$ such that $f$ is continuous at $x$. By Lemma \ref{lem:Tamecont}, the only possible accumulation value of $g$ at $x$ is $f(x)$. 
 
 If $g$ is discontinuous at $x$ then there exists a neighborhood $\V$ of $f(x)$ as well as a sequence $(x_n)_{n\in \N} \in M^\N$ such that $x_n\xrightarrow{n\rightarrow +\infty} x$ and $\forall n\in \N, g(x_n)\notin \V$. Since $N$ is normal, there exists a neighborhood $\V_0$ of $x$ such that $\overline \V_0 \subset \V$. 

 For  all open neighborhood $\W$ of $x$ and all $n\in \N$ big enough,  $x_n\in \W$ so  $$g(x_n)\in \overline {g(\U\cap \W )} = \overline{f(\U\cap \W)} \quad g(x_n) \in M\setminus \overline \V_0$$
then $f(\W)\setminus \V_0 \supset f(\U\cap \W) \cap (M\setminus \overline \V_0) \neq \emptyset$. Since $\W$ is arbitrary, this contradicts continuity of $f$ at $x$. 
We deduce that $g$ is also continuous at $x$ and that $\Tame(g)\subset \Tame(f)$. 
By symmetry, $\Tame(g)=\Tame(f)$.

 
\end{proof}


\begin{prop}\label{prop:weak_morphism_group_acting} Let $M\xrightarrow{f} X$ be a weak a.e. $(G,X)$-morphism
and let  $\Gamma$ be a group acting on $M$ via a.e. $(G,X)$-morphisms.
Then
\begin{enumerate}[(i)]
 \item $\Tame(f)$ only depends on $M$ and is then denoted by $\Tame(M)$;
 \item $\Tame(M)$ is $\Gamma$-invariant;
 \item there exists a unique morphism $\Gamma \xrightarrow{\rho} G$ such that $f_{|\Tame(M)}$ is $\rho$-equivariant. $\rho$ is called the holonomy of $f$.
\end{enumerate}
\end{prop}
\begin{proof}

Let $f_1=f$ and $f_2$ be two weak a.e. $(G,X)$-morphisms from $M$ to $X$ and let $\U_1,\U_2$ be two portly open subsets such that $f_{1|\U_1},f_{2|\U_2}$ are a.e. $(G,X)$-morphisms. 
  By Lemma \ref{lem:MtoXG}, there thus exists some $g\in G$ such that $f_{1|\U_3} = gf_{2|\U_3}$. Consider such a $g \in G$, on the one hand  $f_{1}$ and $gf_{2}$ are equivalent; on the other hand Lemma \ref{lem:Tamecont} implies that $\Tame(f_2)=\Tame(gf_2)$. Using  Lemma \ref{lem:Tame_equ}, $$\Tame(f_1)=\Tame(gf_2)=\Tame(f_2).$$

Let $\gamma \in \Gamma$, clearly $f_1\circ \gamma^{-1} $ is a weak a.e. $(G,X)$-morphism, then 
$$\mathrm{Tame}(M)=\mathrm{Tame}(f_1\circ \gamma^{-1}) =  \gamma \mathrm{Tame}(f_1)=\gamma \mathrm{Tame}(M)$$ 
where the second equality comes from Lemma \ref{lem:Tamecont}.
Since  $f_{|\Tame(M)}$ is an a.e. $(G,X)$-morphism with values in $X$ the last property follows from Lemma \ref{lem:holonomy}.
\end{proof}



\begin{defi}[Weak developing map] \label{defi:weakdevmap}
Let $N\xrightarrow{p} M$ be a quasi-Galoisian branched covering of singular $(G,X)$-manifold of Galois group $\Gamma$. 
A weak developing map for  $N\xrightarrow{p} M$ is a weak a.e. $(G,X)$-morphism $N\xrightarrow{\D} X$ of holonomy $\Gamma \xrightarrow{\rho} G$ such that
\begin{enumerate}[(a)]
 \item $\D$ is $\rho$-equivariant.
 \item $\D : N\rightarrow (M,\T^\rho)$ is continuous  where $\T^\rho$ is the topology of $X$ given by $\Gamma$-invariant open subsets (with $\Gamma$ acting via $\rho$).
\end{enumerate}
\end{defi}
\begin{rem}\label{rem:same_notations_most_one} With the same notations, there is at most one equivalence class of weak a.e. $(G,X)$-morphisms $M\rightarrow X$. In such a class, there is always at least one equivariant map but the existence of a developing map is not garanteed. 
\end{rem}

\begin{rem}\label{rem:quasi_galoisian_branched_covering_singular}
 Let $N\xrightarrow{p}M$ be a quasi-Galoisian branched covering of singular $(G,X)$-manifolds of Galois group $\Gamma$ and let $N\xrightarrow{\D} X$ be a weak developing map. Then : 
$\forall x\in N, \forall \gamma \in \Gamma_x, \D(x) \in \mathrm{Fix}(\rho(\gamma))$.
\end{rem}




Several weak versions of results obtained for developing maps can be proved. The proof are left to the reader. 
\begin{lem}\label{lem:galoisian_branched_covering_manifolds}
 Let $N\xrightarrow{p}M$ be a Galoisian branched covering of a.e. $(G,X)$-manifolds. Then the following are equivalent:
 \begin{enumerate}[(i)]
  \item There exists a developing map $N\xrightarrow{\D}X$ and its holonomy is trivial;
  \item There exists a weak developing map $N\xrightarrow{\D}X$ and its holonomy is trivial;
  \item There exists an a.e $(G,X)$-morphism $M\xrightarrow{\overline \D}X$.
 \end{enumerate}
\end{lem}
\begin{proof}
$(i)\Rightarrow (ii)$ is trivial and $(iii)\Leftrightarrow (i)$ follows from Lemma \ref{lem:descend_developing_map}.

Assume $(ii)$. Recall that $N\rightarrow M$ is Galoisian so $\Gamma \backslash N=M$. Assumption $(b)$ in Definition \ref{defi:weakdevmap} implies that $N\xrightarrow{\D} X$ is continuous and thus induces a continuous map $\overline \D : \Gamma \backslash N \rightarrow   X$.


 Since $\mathrm{Tame}(N)$ is $\Gamma(N/M)$-invariant $p$ induces a branched covering $\mathrm{Tame}(N)\rightarrow p(\mathrm{Tame}(N))$ of which $\D_{|\mathrm{Tame}(N)}$ is a developing map.  From Lemma \ref{lem:descend_developing_map} we deduce that $\overline \D_{|p(\mathrm{Tame}(N))}$ is an a.e. $(G,X)$-morphism. Since $p(\mathrm{Tame}(N))$ is portly, $\overline \D$ is an a.e. $(G,X)$-morphism.
\end{proof}

\begin{rem}\label{rem:weak_dev_map_comp} If $N\xrightarrow{\D} X$ is a developing map of $N\rightarrow M$ and $\widehat N\xrightarrow{p} N$ is a branched covering then the composition $\widehat \D=\D\circ p$ is a natural candidate for a weak developing map for $\widehat N\rightarrow M$. 
\begin{itemize}
 \item For all $\hat\gamma\in \Gamma(\widehat N/M)$, $\gamma$ its projection in $\Gamma(N/M)$, and for all $x_0\in \widehat N$ we have:
 $$ \widehat \D(\hat \gamma x_0 ) := \D\circ p(\hat\gamma x_0) = \D(\gamma p(x_0)) = \rho(\gamma)\D( p(x_0))$$
In particular, $\widehat \D$ is $\widehat \rho$-equivariante with $\widehat \rho$ the composition of $\rho$ by the natural projection $\Gamma(\widehat N/M) \rightarrow \Gamma(N/M)=\Gamma(\widehat N/M)/\Gamma(\widehat N/N)$.

\item Since $N\xrightarrow{\D} (X,\T^\rho)$ and $\widehat N \xrightarrow{p} N$ are continuous, so is $\widehat N \xrightarrow{\widehat \D} (X,\T^\rho)$. 
\end{itemize}
Therefore, $\widehat \D$ is indeed a developing map for $\widehat N\rightarrow M$.

However, it is not necessarily a developing map for $\widehat N\rightarrow N$. Beware that hypothesis $(b)$ for weak developing map may not be satisfied. A counter-example can be constructed by considering $N\rightarrow M$ such that $\Tame(N)\neq N$ and then taking $\widehat N$ as the maximal covering of $N$ possibly branched over some skeletal subset $S\subset \reg(N)$ such that $\pi_1(\reg(N)\setminus S)\neq \{1\}$. In such a situation, condition $(b)$ is
``$\widehat N \xrightarrow{\widehat\D} (X,\T^{\hat\rho_{|\Gamma(\widehat N/N)}})$ is continuous''. But $\hat\rho_{|\Gamma(\widehat N/N)}$ is trivial so $(b)$ is equivalent to 
``$\widehat N \xrightarrow{\widehat\D} X$ is continuous''. Which is not the case since $\D$ is not continuous. In fact one can check directly that is such a situation,  $(i-iii)$ in the previous Lemma all fail (as they should).
\end{rem}



\begin{prop} \label{prop:branched_weak_devmap}Let $M$ be a singular $(G,X)$-manifold. Assume that $\reg(M)$ is semi-locally simply connected in $M$.
The following are then equivalent:
\begin{enumerate}[(i)]
 \item  $\widetilde M\rightarrow M$ admits a weak developing map,
 \item there exists a Galoisian branched covering of $N\xrightarrow{q} M$ which admits a weak developing map and such that $\reg(\ord_{q}(M))$ is semi-locally simply connected in $M$.
\end{enumerate} 
\end{prop}
\begin{proof} The proof is similar to the one of Proposition \ref{prop:branched_devmap}. The implication
  $(i)\Rightarrow (ii)$ is simple enough.  Assume $(ii)$ and consider some branched covering $N\xrightarrow{q}M$  and a weak developing map $N\xrightarrow{\D_N} X$ of holonomy $\Gamma(N/M)\xrightarrow{\rho} G$. Denote by $S:=M\setminus \ord_p(M)$ and $\widetilde S$  the lift of $S$ to $\widetilde M:=\widetilde M^{\sing(M)}$.   
 Apply Proposition \ref{prop:universality} to  construct the maximal branched covering $\widehat M := \widetilde M^{S\cup \sing(M)}\xrightarrow{\hat p} M$ as well as the maximal branched covering $\widetilde M:=\widetilde M^{\sing(M)} \xrightarrow{p} M$. Denote by $\widehat M \xrightarrow{\hat p_{\cdot}} \cdot$ the branched covering induced by $\hat p$ on subcoverings of $\widehat M\rightarrow M$. Since $\sing(M)\subset S\cup \sing(M)$ and $S\subset S\cup \sing(M)$ both $N$ and $\widetilde M$ are subcoverings of $\widehat M$. We wish to construct a weak developing map $\widetilde M \xrightarrow{\D} X$. 
 \begin{itemize}
 \item  From the weak developing map $\D_N$ we define a weak developing map $\widehat M \xrightarrow{\widehat \D :=\D\circ \hat p_N} X$ of $\widehat M\rightarrow M$ by remark \ref{rem:weak_dev_map_comp}.
\item  Since $\reg(\ord_p(M))$ is semi-locally simply connected in $M$, by Proposition \ref{prop:Galoisian_semilocal} the maximal covering $\widehat M$ possibly branched over $\sing(M)\cup S = M\setminus \reg(\ord_p(M))$ is Galoisian. In particular, $\widetilde M = \Gamma(\widehat M/\widetilde M)\backslash\widehat M$ by remark \ref{rem:galois_quotient}.
\item Consider $\U_M:= \reg(\ord_{\hat p}(M))$, $\U_{\widetilde M}:=p^{-1}(\U_M)$, $\U_{N}:=q^{-1}(\U_M)$ and $\U_{\widehat M}:=\hat p^{-1}(\U_M)$. The weak developing maps $\widehat \D$ and $\D_N$ restricted to $\U_{\widehat M}$ and $\U_{N}$ respectively are continuous developing maps. Moreover, there exists a developing map $\U_{\widetilde M}\xrightarrow{\D}X$ which can be chosen so that $\D\circ \hat p_{\widetilde M|\U_{\widehat M}}= \widehat \D_{|\widehat M}$. 

By Lemma \ref{lem:descend_developing_map}, the holonomy of $\U_{\widehat M} \rightarrow \U_{\widetilde M}$ is trivial, since this holonomy is given by $\hat \rho_{|\Gamma(\widehat M/\widetilde M )}$ we dedude that the continuous map $\widehat M \xrightarrow{\widehat \D} (X,\T^{\hat\rho})$ is invariant under the action of $\Gamma(\widehat M/\widetilde M)$.

\item The previous points allow to descend $\widehat \D$ to a weak developing map $\widetilde M \xrightarrow{\D} X$ 
\end{itemize}

\begin{figure}
  $$
  \xymatrix{ 
  &&X&&\\
 &&\widehat M \ar[drr]_{\hat p_N} 
 \ar[dd]^{\hat p} \ar[dll]^{\hat p_{\widetilde M}}
 \ar[u]^{\widehat \D}
 &&\\
 \widetilde M\ar[drr]_{p} \ar@{-->}[uurr]^{\D}&&&&N \ar[dll]^{q}\ar[uull]_{\D_N}
 \\&& M&& 
  }\quad \quad 
  \xymatrix{ 
  &&X&&\\
 &&\U_{\widehat M} \ar[drr]_{\hat p_N} 
 \ar[dd]^{\hat p} \ar[dll]^{\hat p_{\widetilde M}}
 \ar[u]^{\widehat \D}
 &&\\
 \U_{\widetilde M}\ar[drr]_{p} \ar[uurr]^{\D}&&&&\U_N \ar[dll]^{q}\ar[uull]_{\D_N}
 \\&& \U_M&& 
  }
$$

\end{figure}

 
 
 
\end{proof}







%Finally, every parabolic isometries are conjugated and the  centralizer of a parabolic isometry is the 1-parameter Lie subgroup is generates; therefore, 

		
